\documentclass[journal=tosc,submission,notanonymous]{iacrtrans}

% Figures and pseudocode.
\usepackage{tikz}
\usetikzlibrary{positioning,arrows.meta,fit,backgrounds,calc}
\usepackage{algorithm}
\usepackage{algpseudocode}
% Disambiguate algpseudocode line anchors so multiple algorithms with
% line-numbering don't collide under hyperref.
\makeatletter
\renewcommand{\theHALG@line}{\thealgorithm.\arabic{ALG@line}}
\makeatother

% Tables.
\usepackage{booktabs}
\usepackage{comment}

% Lists.
\usepackage{enumitem}

% The IACR text block is narrow; allow TeX to loosen dense math prose before
% reporting overfull lines.
\emergencystretch=2em

% IACR loads hyperref at the end of the preamble; load cleveref after it.
\AtEndPreamble{%
  \usepackage[capitalize,noabbrev]{cleveref}%
  \crefname{section}{\S}{\S\S}%
  \Crefname{section}{Section}{Sections}%
  \crefname{lemma}{Lemma}{Lemmas}%
  \Crefname{lemma}{Lemma}{Lemmas}%
  \crefname{corollary}{Corollary}{Corollaries}%
  \Crefname{corollary}{Corollary}{Corollaries}%
  \crefname{proposition}{Proposition}{Propositions}%
  \Crefname{proposition}{Proposition}{Propositions}%
  \crefname{theorem}{Theorem}{Theorems}%
  \Crefname{theorem}{Theorem}{Theorems}%
  \crefname{algorithm}{Algorithm}{Algorithms}%
  \Crefname{algorithm}{Algorithm}{Algorithms}%
  \crefname{figure}{Figure}{Figures}%
  \Crefname{figure}{Figure}{Figures}%
  \crefname{remark}{Remark}{Remarks}%
  \Crefname{remark}{Remark}{Remarks}%
}

% Scheme.
\newcommand{\TW}{\textsf{TW128}}
\newcommand{\Enc}{\mathsf{Enc}}
\newcommand{\Dec}{\mathsf{Dec}}
\newcommand{\Ver}{\mathsf{Ver}}
\newcommand{\Rand}{\mathsf{Rand}}
\newcommand{\SE}{\mathsf{SE}}
\newcommand{\WiC}{\mathsf{WiC}}

% Sponge / duplex / permutation.
\newcommand{\Keccak}{\textsc{Keccak}-\textit{p}}
\newcommand{\Keccakp}{\mbox{\Keccak[1600, 12]}}
\newcommand{\Sponge}{\mathsf{Sponge}}
\newcommand{\Duplex}{\mathsf{Duplex}}
\newcommand{\duplexing}{\mathsf{duplexing}}
\newcommand{\padten}{\mathsf{pad10^*}}
\newcommand{\padtenone}{\mathsf{pad10^*1}}
\newcommand{\RO}{\mathsf{RO}}
\newcommand{\ROflat}{\RO_{\mathsf{flat}}}
\newcommand{\flatmap}{\mathsf{flat}}
\newcommand{\RF}{\mathsf{RF}}
\newcommand{\KeyedTWOut}{\mathsf{KeyedTWOut}}
\newcommand{\bridge}{I}

% Resources.
\newcommand{\qv}{q_{\mathsf{v}}}
\newcommand{\Qv}{Q_{\mathsf{v}}}
\newcommand{\qRO}{q_{\RO}}
\newcommand{\QRO}{Q_{\RO}}
\newcommand{\Nleaf}{N_{\mathsf{leaf}}}
\newcommand{\nleaf}{n_{\mathsf{leaf}}}
\newcommand{\Nprim}{N_{\mathsf{prim}}}
\newcommand{\Ntot}{N_{\mathsf{tot}}}

% Advantages.
\newcommand{\Adv}{\mathbf{Adv}}
\newcommand{\Advpriv}{\Adv^{\mathsf{priv}}}
\newcommand{\Advintctxt}{\Adv^{\mathsf{int\text{-}ctxt}}}
\newcommand{\Advcmtds}[1]{\Adv^{\mathsf{CMTDs}\text{-}#1}}
\newcommand{\Advcmts}[1]{\Adv^{\mathsf{CMTs}\text{-}#1}}
\newcommand{\Advtagcmtds}[1]{\Adv^{\mathsf{tag\text{-}CMTDs}\text{-}#1}}
\newcommand{\Advcdy}{\Adv^{\mathsf{CDY}}}
\newcommand{\Epspriv}{\varepsilon^{\mathsf{priv}}_{\RO}}
\newcommand{\Epsintctxt}{\varepsilon^{\mathsf{int\text{-}ctxt}}_{\RO}}
\newcommand{\Advae}{\Adv^{\mathsf{ae}}}
\newcommand{\Advmuind}{\Adv^{\mathsf{mu\text{-}ind}}}

% Simulator and ideal verification oracles.
\newcommand{\Sim}{\mathcal{S}}
\newcommand{\Replay}{\mathsf{Replay}}

% Bad events and loss terms.
\newcommand{\badkey}{\mathsf{bad}_{\mathsf{key}}}
\newcommand{\badleaf}{\mathsf{bad}_{\mathsf{leaf}}}
\newcommand{\Lift}{\mathsf{Lift}}
\newcommand{\KeyGuess}{\mathsf{KeyGuess}}
\newcommand{\RootColl}{\mathsf{RootColl}}
\newcommand{\RootPre}{\mathsf{RootPre}}
\newcommand{\LeafColl}{\mathsf{LeafColl}}

% Tags / parameters.
\newcommand{\tauroot}{\tau_{\mathsf{root}}}
\newcommand{\tauleaf}{\tau_{\mathsf{leaf}}}
\newcommand{\prfxroot}{p_{\mathsf{root}}}
\newcommand{\prfxleaf}{p_{\mathsf{leaf}}}
\newcommand{\TagProj}{\gamma_{\mathsf{tag}}}

% Phase suffixes.
\newcommand{\siginit}{\ensuremath{\sigma_{\mathsf{init}}}}
\newcommand{\sigadmore}{\ensuremath{\sigma_{\mathsf{ad}}^-}}
\newcommand{\sigadlast}{\ensuremath{\sigma_{\mathsf{ad}}^+}}
\newcommand{\sigmsgmore}{\ensuremath{\sigma_{\mathsf{msg}}^-}}
\newcommand{\sigmsglast}{\ensuremath{\sigma_{\mathsf{msg}}^+}}
\newcommand{\sigaggmore}{\ensuremath{\sigma_{\mathsf{agg}}^-}}
\newcommand{\sigagglast}{\ensuremath{\sigma_{\mathsf{agg}}^+}}
\newcommand{\sigph}[1][\relax]{\ensuremath{\sigma_{\mathsf{ph}\ifx\relax#1\else,#1\fi}}}

% Duplex-call output-point classes.
\newcommand{\Rinit}{\ensuremath{\mathcal{R}_{\mathsf{init}}}}
\newcommand{\Rad}{\ensuremath{\mathcal{R}_{\mathsf{ad}}}}
\newcommand{\Rmsg}{\ensuremath{\mathcal{R}_{\mathsf{msg}}}}
\newcommand{\Ragg}{\ensuremath{\mathcal{R}_{\mathsf{agg}}}}
\newcommand{\Linit}{\ensuremath{\mathcal{L}_{\mathsf{init}}}}
\newcommand{\Lmsg}{\ensuremath{\mathcal{L}_{\mathsf{msg}}}}

% Transcripts.
\newcommand{\Rtag}{\ensuremath{\mathcal{R}_{\mathsf{tag}}}}
\newcommand{\Ltag}{\ensuremath{\mathcal{L}_{\mathsf{tag}}}}
\newcommand{\Flat}{\mathsf{Flat}}

% Misc.
\newcommand{\Perm}{\mathsf{Perm}}
\newcommand{\bits}{\{0,1\}}
\newcommand{\sample}{\mathrel{\mathord{\leftarrow}\mkern1mu\vcenter{\hbox{$\scriptstyle\$$}}}}
\newcommand{\concat}{\mathbin{\|}}
\newcommand{\bigconcat}{\mathop{\Big\Vert}\limits}
\DeclareMathOperator{\leaves}{leaves}


\title[TW128]{TW128: A Tree-Structured Duplex Authenticated Encryption Mode}
\author{Coda Hale}
\institute{Hale Institute for Bicycle Appreciation}

\begin{document}

\maketitle

\keywords{authenticated encryption \and committing security \and sponge constructions \and duplex constructions \and Keccak}

\begin{abstract}
  AEAD tags are often treated as compact commitments to the full encryption
  context, but standard deployed modes such as AES-GCM and
  ChaCha20-Poly1305 do not provide this property. We present \TW{}, a
  nonce-based authenticated encryption mode with associated data designed around
  the stronger tag-centric requirement: given a candidate opening, the
  authentication tag should itself behave like a digest of the full
  $(K, U, A, M)$ context, including the key.

  \TW{} adapts the KangarooTwelve tree-hashing pattern to keyed duplex
  encryption over \Keccakp{}. A serial root duplex binds the key, nonce,
  associated data, and first message chunk, then aggregates tags from
  independently keyed leaf duplexes for the remaining chunks. The leaf
  decomposition is canonical, while the number of leaves evaluated at once is an
  implementation choice; scalar, SIMD, and multicore implementations therefore
  produce identical ciphertexts and tags without exposing the degree of
  parallelism as a mode parameter.

  The final root tag is both the AEAD authentication tag and the object of the
  main security analysis. In the public permutation model, via an intermediate
  random oracle and the flat sponge lift, we prove that the wrapper returning
  this tag is a secure hash of the full input: collision-, second-preimage-, and
  preimage-resistant in the standard hash-function sense. The same root tag
  analysis yields AEAD-facing committing security, including full-input
  \textsf{CMTDs-4}. The root-centered structure also gives integrity under
  release of unverified plaintext in the sense of \cite{ABLMMY14}. As baseline
  AEAD properties, we prove concrete multi-user indistinguishability and derive
  all-in-one authenticated encryption as the single-user corollary. At the
  implemented parameters, all stated notions meet a $128$-bit security target.
  Vectorized \texttt{amd64} and \texttt{arm64} implementations recover the
  scalable performance profile of the tree design: in bulk they reach
  $48$\,Gbit/s, with the AVX-512 and \texttt{arm64} paths delivering
  $51$--$76\%$ of hardware-accelerated AES-128-GCM throughput.
\end{abstract}

\section{Introduction}
\label{sec:introduction}

Authenticated encryption with associated data (AEAD) is often treated in
practice as if its tag were a compact commitment to the full encryption context:
a short string that binds the key, nonce, associated data, and message. One
source of this intuition is that many practitioners first meet message
authentication through hash-based MACs such as HMAC, whose tags look and behave
much like keyed digests. Standard modern AEADs do not provide this full
commitment property at all. AES-GCM and ChaCha20-Poly1305, for example, are not
key-committing: an adversary can craft a single ciphertext that decrypts to
valid plaintexts under two or more keys \cite{DGRW18, ADGKLS22}. The cause is
structural. Their authenticity rests on algebraic MACs whose outputs can often
be made to validate under a second context, whereas a digest is expected to
resist both collisions and preimages.

The resulting attacks are not edge cases in the formalism; they are failures of
the compact-commitment intuition that users and protocols naturally assign to
tags. The ``invisible salamander'' attack defeats Facebook's attachment
franking scheme \cite{DGRW18}; related constructions produce ciphertexts that
open to different well-formed files under different keys \cite{ADGKLS22};
partitioning oracles recover passwords from systems built on non-committing AEAD
\cite{LGR21}; and context discovery attacks break \textsf{GCM}, \textsf{CCM},
\textsf{EAX}, \textsf{SIV}, and \textsf{OCB3} by exploiting authentication
layers that are not preimage-resistant \cite{MLGR23}. These attacks differ in
surface form, but all expose the same mismatch: the tag is used as if it were a
full-context commitment, while the mode only proves ordinary authenticity.

The AEAD literature has only recently begun to model this expectation directly.
Compactly committing AEAD was introduced for message franking \cite{GLR17}, and
the encryptment approach of \cite{DGRW18} builds commitment into a dedicated
primitive. The committing hierarchy of \cite{BH22} makes the collision-style
requirement explicit: no two distinct contexts should open the same ciphertext
and tag. Context discoverability \cite{MLGR23} captures the preimage-facing
counterpart: given a target ciphertext, it should be infeasible to find a
decrypting context. These notions explain when a ciphertext can have another
valid opening, or any valid opening at all, but they are still phrased as AEAD
games over ciphertexts. The practitioner's assumption is more tag-centric: the
tag itself is treated as a compact commitment string. We therefore analyze the
tag with hash notions directly, and then recover the AEAD commitment notions as
game-specific consequences of the same tag analysis.

\TW{} (\emph{TreeWrap}, with $128$ denoting the $128$-bit security target)
takes the digest intuition literally. If an AEAD tag is meant to behave like a
compact digest, the mode should be shaped like a modern high-performance digest
rather than like a serial stream cipher with an algebraic MAC. The relevant
performance lesson from modern hashing is that large inputs are processed by
splitting them into chunks, evaluating independent work in parallel, and
aggregating short chaining values into a final root value.
KangarooTwelve \cite{BDPVVV16} is the SHA-3-family exemplar of this pattern: it
hashes fixed-size chunks independently under \Keccakp{}, absorbs their chaining
values into a root node, runs leaves across SIMD lanes and cores, and keeps
memory bounded independently of the input length.

\TW{} asks what this tree looks like when the sponge nodes are replaced by
keyed duplexes that encrypt while they absorb. A serial root duplex, initialized
with the key $K$ and nonce $N$, absorbs the associated data, encrypts the first
message chunk, absorbs a tag from each remaining chunk, and emits the final root
tag. Each remaining chunk is processed by an independent leaf duplex keyed by
$(K, N, j)$, where $j$ is the chunk index; the leaf encrypts its chunk and
returns a leaf tag to the root. Thus every input reaches the root directly or
through a leaf tag, while every duplex state remains bounded at $1600$ bits. As
in KangarooTwelve's \emph{kangaroo hopping}, a single-chunk message uses only
the root duplex and pays no tree overhead (\Cref{sec:latency}).

This gives a mode that is not merely parallelizable but scalable. The leaf
decomposition and transcript are canonical, so there is only one ciphertext and
tag for a given input, but the number of leaves evaluated at once is left to the
implementation. A scalar implementation, an AVX2 implementation, an AVX-512
implementation, a NEON implementation, and a multicore implementation all
produce identical ciphertexts and tags. The degree of parallelism is therefore
not a mode parameter and is not visible on the wire. This separates \TW{} from
prior duplex-based AEADs with parallel variants, where the parallelism degree is
fixed by the mode instance or otherwise shared by the communicating parties.

Because \TW{} adapts the KangarooTwelve tree to encryption, the tag emitted by
the root is naturally the output of a keyed tree hash. We formalize this by
considering the wrapper $\HTW$ that returns the root tag, and prove that $\HTW$
is a secure hash in the sense of \cite{RS04}: collision-, second-preimage-, and
preimage-resistant as a function of the full $(K, N, A, M)$ input. The $256$-bit
root tag is therefore a compact digest-like commitment to the full context. The
familiar AEAD commitment notions are recovered from the same root tag candidate
analysis: full-input $s$-way \textsf{CMTDs-4} \cite{BH22} is the
collision-facing AEAD consequence, while \cite{MLGR23} context discoverability
$\mathsf{CDY}^\ast$ is the preimage-facing consequence, with an additional
hidden-component guessing term. The proofs use the structure of the AEAD games
to keep the stated bounds tight.

We also prove the standard AEAD security of the mode. In the public permutation
model, \TW{} has a concrete multi-user indistinguishability bound in the sense
of \cite{BT16}, with all-in-one authenticated encryption in the sense of
\cite{RS06} as its single-user corollary. Each bound is proved first against an
intermediate random oracle and then transferred to the real permutation by the
flat sponge lift of \cite{BDPV08} (\Cref{app:flat-sponge-lift}). At the
implemented parameters ($k = c = \tauroot = \tauleaf = 256$), the bounds meet a
$128$-bit security target for adversaries running up to $\Ntot \leq 2^{63}$
\Keccakp{} calls---on the order of $2^{71}$ bytes of processed data.

The mode's design-level parallelism translates into measured implementation
performance. Optimized \texttt{amd64} (AVX-512 and AVX2) and \texttt{arm64}
(NEON with the SHA-3 extension) implementations evaluate leaf duplexes in
batches while preserving the canonical transcript. On contemporary hardware,
\TW{} exceeds $10$\,Gbit/s and reaches $50$--$75\%$ of hardware-accelerated
AES-128-GCM
(\Cref{sec:performance}), while using one permutation for confidentiality,
authentication, and commitment.

\subsection{Our Contributions}
\label{sec:contributions}

\begin{itemize}
\item \textbf{A scalable tree AEAD mode.} We specify \TW{}, a nonce-based AEAD
  with associated data built as a two-level tree of strict \cite{BDPV11}
  duplexes over \Keccakp{} (\Cref{sec:tw128}). The mode fixes the root/leaf
  transcript but delegates the number of leaves evaluated simultaneously to the
  implementation, so scalar and vectorized implementations interoperate without
  changing ciphertexts or tags.

\item \textbf{Digest-like commitment.} We prove that the wrapper $\HTW$
  returning the root tag is a secure hash of the full $(K, N, A, M)$ input in
  the sense of \cite{RS04}: collision-, second-preimage-, and
  preimage-resistant (\Cref{thm:coll,thm:asec,thm:epre}). Thus the
  authentication tag behaves like the compact digest practitioners often expect
  an AEAD tag to be.

\item \textbf{AEAD commitment consequences.} We instantiate the same root tag
  analysis in the existing AEAD vocabulary. The all-bodies-equal structure of
  \textsf{CMTDs-4} gives the full-input $s$-way committing bound without the
  leaf-collision slack of the general hash collision theorem \cite{BH22}
  (\Cref{cor:cmtds4}), and the remaining \textsf{CMT}/\textsf{CMTD} notions
  for $\ell \in \{1, 3, 4\}$ follow from the committing hierarchy and direct
  source-game arguments (\Cref{cor:committing-combined}). The same fixed-target
  root-preimage analysis, plus hidden-component guessing, gives
  $\mathsf{CDY}^\ast$ context discoverability security \cite{MLGR23}
  (\Cref{thm:cdy-full}).

\item \textbf{Multi-user and AE security.} We give a concrete multi-user
  indistinguishability bound \cite{BT16} for $D$ users (\Cref{thm:mu-ind}).
  The $D = 1$ case yields a concrete all-in-one AE bound \cite{RS06} in the
  public permutation model (\Cref{cor:ae}). The proof proceeds through a
  dedicated \TW{} random oracle model and is transferred to the public
  permutation setting by the \cite{BDPV08} flat sponge lift.

\item \textbf{Optimized implementations.} We describe and evaluate vectorized
  \texttt{amd64} (AVX-512 and AVX2) and \texttt{arm64} (NEON with the SHA-3
  extension) implementations whose batched leaf cores realize the mode's
  scalable parallelism in practice (\Cref{sec:performance}).
\end{itemize}

\subsection{Related Work}
\label{sec:related}

\paragraph{Committing AEAD.}
The study of committing authenticated encryption was initiated by \cite{GLR17},
who introduced compactly committing AEAD for message franking, and by
\cite{DGRW18}, who built it from a dedicated primitive (encryptment).
\cite{BH22} organized the collision-style notions into the
\textsf{CMT-1}/\textsf{3}/\textsf{4} and \textsf{CMTD} hierarchy and gave the
\textsf{UtC}, \textsf{RtC}, and \textsf{HtE} transforms; \cite{CR22} gave the
\textsf{CTX} transform, which commits a nonce-based AE scheme to its full
context by adding a hash. These transforms are efficient and broadly
applicable, but they treat commitment as an external layer. \TW{} instead makes
the authentication tag itself the hash-like commitment: the root tag is the
compact commitment, and the root tag analysis gives both the \cite{BH22}
full-input committing consequence and the \cite{MLGR23} context discoverability
bound. Closest in spirit is the recent work
of \cite{DHMVV24}, which also builds nonce-based AE natively on SHA-3 functions
and obtains \textsf{CMT-4} from collision resistance; it targets session-based
communication through a sequential duplex chain, whereas \TW{} is a single
scalable tree for standalone AEAD
(\Cref{sec:compare-commit}).

\paragraph{Sponge and duplex AEADs.}
Permutation-based authenticated encryption built on the sponge and duplex
constructions of \cite{BDPV08, BDPV11} is well established. Keyak
\cite{BDPVV16} uses the Motorist mode over round-reduced \Keccak{} and already
admits parallel duplex instances; Ascon \cite{DEMS21} and Xoodyak
\cite{DHPVV20} are serial duplex AEADs over smaller permutations. \TW{} shares
the single-permutation, constant-state duplex foundation but differs in where
parallelism lives. Keyak fixes the degree of parallelism as a mode parameter,
while Ascon and Xoodyak expose no mode-level parallelism. In \TW{}, the
decomposition into indexed leaves is fixed by the transcript, but the number of
leaves evaluated at once is chosen by the implementation and is invisible to the
wire format (\Cref{sec:compare-sponge}).

\paragraph{Tree hashing.}
\TW{}'s two-level shape is inspired by KangarooTwelve \cite{BDPVVV16}, a tree
hash over the same permutation. \TW{} adopts its choice of permutation and
final-node/leaf split and adapts them to authenticated encryption: each leaf
chunk is encrypted and tagged independently, leaf tags replace unkeyed chaining
values, and the root is keyed by $(K, N)$. The construction also inherits
KangarooTwelve's short-message behavior: inputs fitting in one chunk use only
the root duplex. Because \TW{} has a fixed two-level topology with independently
keyed, indexed leaves and phase-separated duplex calls, it does not need
KangarooTwelve's Sakura tree encoding \cite{BDPV14} to make its transcripts
unambiguous (\Cref{sec:tw128}).

\subsection{Paper Organization}
\label{sec:organization}

\Cref{sec:prelims} recalls the notation, sponge and duplex constructions,
nonce-based AEAD security, committing security \cite{BH22}, context
discoverability \cite{MLGR23}, and the hash notions of \cite{RS04}.
\Cref{sec:tw128} specifies the \TW{} construction and parameters.
\Cref{sec:security-model} fixes the public permutation games, intermediate
random oracle model, resource accounting, and loss terms.
\Cref{sec:aead-security} proves multi-user indistinguishability and all-in-one
AE security. \Cref{sec:root-tag-hash-security} proves root-tag hash security,
committing security, and context discoverability, and
\Cref{sec:concrete-security} summarizes the concrete bounds.
\Cref{sec:performance} reports implementation performance, and
\Cref{sec:discussion} compares \TW{} to prior AEAD and committing constructions
and discusses limitations. The appendices collect structural lemmas, the flat
sponge lift, vectorized kernel details, and test vectors.

\section{Preliminaries}
\label{sec:prelims}

\subsection{Notation}
\label{sec:notation}

\paragraph{Strings.}
Write $\bits$ for the binary alphabet. For an integer $n \geq 0$, $\bits^n$
denotes the set of $n$-bit strings, $\bits^* = \bigcup_{n \geq 0} \bits^n$
the set of all finite-length bit strings, and $\varepsilon$ the empty
string. For $x \in \bits^*$, $|x|$ is the bit length and $x \concat y$ is
the concatenation; the bitwise XOR of two equal-length strings $x, y \in
\bits^n$ is $x \oplus y$. A \emph{byte string} is an element of
$(\bits^8)^*$; its byte length is denoted $\|x\| = |x|/8$. For an
integer $0 \leq n < 2^{64}$, $\langle n \rangle_8 \in \bits^{64}$ is its
eight-byte little-endian encoding. For a finite or infinite bit string
$x$ and $0 \leq \tau \leq |x|$ when $x$ is finite,
$\lfloor x \rfloor_\tau \in \bits^\tau$ denotes the leading $\tau$ bits
of $x$.

\paragraph{Sampling and probability.}
We write $x \sample S$ for sampling $x$ uniformly at random from a finite
set $S$, with implicit coins independent across samples. For an event $E$
in a probability experiment, $\Pr[E]$ denotes its probability.

\paragraph{Sets of maps.}
For a set $S$, $\Perm(S)$ denotes the set of permutations on $S$; for a
positive integer $b$, $\Perm(\bits^b)$ is the set of $b$-bit
permutations.

\paragraph{Games and advantages.}
We write $\mathcal{A}^{O_1, \dots, O_k}$ for a (possibly randomized)
adversary with oracle access to $O_1, \dots, O_k$, and
$\Pr[\mathsf{G}^{\mathcal{A}} \Rightarrow 1]$ for the probability that a
game $\mathsf{G}$---an experiment that runs $\mathcal{A}$ against
specified oracles and returns a Boolean outcome---returns $1$. For a
distinguishing experiment between two games $\mathsf{G}_0, \mathsf{G}_1$,
the advantage of $\mathcal{A}$ against a scheme $\Pi$ in notion
$\mathsf{X}$ is the absolute difference
\[
  \Adv^{\mathsf{X}}_\Pi(\mathcal{A})
  = \bigl| \Pr[\mathsf{G}_0^{\mathcal{A}} \Rightarrow 1]
         - \Pr[\mathsf{G}_1^{\mathcal{A}} \Rightarrow 1] \bigr|.
\]
For a win-event experiment without a hidden challenge bit it is instead
$\Pr[\mathsf{G}^{\mathcal{A}} \text{ wins}]$, the probability the game
returns $1$; a bit-guessing experiment with a hidden challenge bit uses
the rescaled form
$\bigl|2 \Pr[\mathsf{G}^{\mathcal{A}} \text{ wins}] - 1\bigr|$, as in the
multi-user notion of \Cref{sec:mu-ind}.

\paragraph{Concrete security.}
All bounds in this paper are concrete: each advantage is an explicit
function of the resource caps introduced in
\Cref{sec:resources}. We make no use of asymptotic or negligibility
conventions.

\paragraph{Reference summary.}
\Cref{tab:notation} summarizes the resource caps and loss terms used
throughout. The caps are defined formally in \Cref{sec:resources}, the
loss terms in \Cref{sec:loss-terms}.

\begin{table}[ht]
  \centering
  \small
  \begin{tabular}{l l l}
    \toprule
    Symbol & Meaning & Defined \\
    \midrule
    $N$        & Construction interface cost (\Keccakp{} calls)       & \cref{sec:resources} \\
    $\Nprim$   & Direct primitive-interface queries                   & \cref{sec:resources} \\
    $\Ntot$    & $N + \Nprim$ (total query-sequence cost)             & \cref{sec:resources} \\
    $\qv$, $\Qv$ & Valid verification queries / cap                   & \cref{sec:resources} \\
    $\qRO$, $\QRO$ & Direct random oracle queries / cap (RO-model)    & \cref{sec:resources} \\
    $\nleaf$, $\Nleaf$ & Distinct construction-generated final leaf transcripts / cap & \cref{sec:resources} \\
    \midrule
    $\Lift_c(\Ntot)$        & \cite{BDPV08} flat sponge loss   & \cref{sec:loss-terms} \\
    $\KeyGuess_k(Q)$        & Ideal system key guess           & \cref{sec:loss-terms} \\
    $\RootColl_\tau(Q, s)$  & $s$-way root-collision           & \cref{sec:loss-terms} \\
    $\RootPre_\tau(Q)$      & Root-preimage                    & \cref{sec:loss-terms} \\
    $\LeafColl_\tau(Q)$     & Known-key leaf-collision         & \cref{sec:loss-terms} \\
    \bottomrule
  \end{tabular}
  \caption{Resource caps and loss terms.}
  \label{tab:notation}
\end{table}

\subsection{Sponges and Duplexes over \texorpdfstring{\Keccak{}}{Keccak-p}}
\label{sec:sponge}

\paragraph{The \Keccak{} permutation.}
The \Keccak{} family of permutations, defined in
FIPS~202~\cite{FIPS202} as the generalization of the Keccak design's
Keccak-$f$ permutations to a free round count, is parameterized
by a state width $b$ and a round count $n_r$; we write $\Keccak[b, n_r]
\in \Perm(\bits^b)$ for the corresponding permutation. \TW{} uses
\Keccakp{}, the same primitive as
KangarooTwelve~\cite{BDPVVV16}, fixed throughout. In the
public permutation security model of \Cref{sec:games-pp} and the
flat sponge lift of \Cref{app:flat-sponge-lift}, this permutation is
modeled as a uniformly random element of $\Perm(\bits^{1600})$.

\paragraph{The sponge construction.}
Following \cite{BDPV08}, fix a permutation $\pi \in \Perm(\bits^b)$, a
\emph{sponge-compliant} padding rule $\mathsf{pad}$, and a rate $1 \leq
r < b$ with capacity $c = b - r$. A padding rule appends
$\mathsf{pad}[r](|M|)$, a bit string determined by the input length and
the rate, to an input $M$, extending it to a sequence of $r$-bit
blocks; it is sponge-compliant
\cite[Definition~1]{BDPV11} if $M \concat \mathsf{pad}[r](|M|)$ is never
empty and, for all $M \neq M'$ and all $n \geq 0$, $M \concat
\mathsf{pad}[r](|M|) \neq M' \concat \mathsf{pad}[r](|M'|) \concat
0^{n r}$. The sponge function $\Sponge[\pi, \mathsf{pad}, r]
: \bits^* \times \mathbb{N} \to \bits^*$ maps an input $M$ and an output
length $\ell$ to an $\ell$-bit string by absorbing $M \concat
\mathsf{pad}[r](|M|)$ into the all-zero $b$-bit state in $r$-bit blocks
(each separated by an application of $\pi$) and squeezing $\ell$ bits in
$r$-bit blocks from the resulting state. \cite[Theorem~2]{BDPV08} shows that
$\Sponge[\pi, \padten, r]$, where the simple rule $\padten$ appends a
single $1$ bit followed by the minimum number of $0$ bits reaching a
multiple of $r$, is indifferentiable from a random oracle when
$\pi$ is uniformly random. We bound its distinguishing advantage by
$\Lift_c(\Ntot) = \Ntot(\Ntot+1) / 2^{c+1}$, derived in
\Cref{lem:lift-bound} from the product upper bounds of
\cite[Lemmas~4 and~5]{BDPV08}.

\paragraph{The duplex construction.}
The \emph{duplex} construction of \cite{BDPV11} extends the sponge to
a stateful interface. A duplex object $\Duplex[\pi, \mathsf{pad}, r]$
is initialized to the all-zero state. Each call $\duplexing(\sigma,
\ell)$, with $\ell \leq r$, absorbs an input string $\sigma$ of at most
$\rho_{\max}(\mathsf{pad}, r)$ bits via $\pi$ and returns the first
$\ell$ bits of the resulting state. Here $\rho_{\max}(\mathsf{pad}, r)
= \min\{x : x + |\mathsf{pad}[r](x)| > r\} - 1$ is the largest input
length for which the padded input fits in a single $r$-bit block; for
the $\padtenone$ rule, $\rho_{\max}(\padtenone, r) = r - 2$
\cite{BDPV11}. \cite[Lemma~3]{BDPV11}
establishes the duplex-to-sponge equality: for any call
sequence $((\sigma_0, \ell_0), \dots, (\sigma_i, \ell_i))$, the $i$-th
output equals
\[
  \Sponge[\pi, \mathsf{pad}, r]\bigl(
    \sigma_0 \concat \mathsf{pad}_0 \concat \cdots \concat \sigma_i,\,
    \ell_i\bigr),
\]
where $\mathsf{pad}_j = \mathsf{pad}[r](|\sigma_j|)$ is the padding
the duplex appends in its $j$-th call; the flattened input carries
these interior paddings explicitly, and the sponge appends only the
final $\mathsf{pad}_i$. The flattening map is
injective by \cite[Lemma~4]{BDPV11}: the duplex input-string sequence is
recoverable from the flattened sponge input.

\paragraph{\TW{} instantiation.}
\TW{} fixes the padding rule $\padtenone{}$ and the rate $r = 1344$, so
the capacity is $c = 256$ bits. The largest input length per duplex call
is therefore $\rho_{\max}(\padtenone, 1344) = 1342$ bits.

\paragraph{Bridge to the \texorpdfstring{$H$}{H}-input domain.}
The indifferentiability theorem of \cite{BDPV08} is stated for the
simpler $\padten{}$ padding over its unpadded $H$-input domain, while
\TW{} uses $\padtenone{}$ internally. \cite[Theorem~3]{BDPV11} supplies
an injective preprocessing bridge $\bridge[r, r_{\max}]$ between the two
domains. For \TW{}, $r = r_{\max} = 1344$, so the bridge reduces to
appending $1 \concat 0^q$ with $q$ determined by the input length modulo
$r_{\max}$; on the bridge image, $q$ is recovered as the trailing-zero
count. \Cref{sec:ro-model} instantiates this bridge on typed \TW{} duplex
transcript prefixes as the map $\flatmap(\cdot)$ used by the
random oracle proofs throughout the paper.

\subsection{Authenticated Encryption with Associated Data}
\label{sec:ae-defs}

\paragraph{Syntax.}
A \emph{nonce-based authenticated encryption scheme with associated
data} (AEAD scheme) is a pair $\SE = (\Enc, \Dec)$ of deterministic
algorithms with associated key space $\mathcal{K} \subseteq (\bits^8)^*$,
nonce space $\mathcal{U} \subseteq (\bits^8)^*$, associated data space
$\mathcal{A} \subseteq (\bits^8)^*$, message space $\mathcal{M} \subseteq
(\bits^8)^*$, and ciphertext space $\mathcal{C} \subseteq (\bits^8)^*$. The
encryption algorithm $\Enc : \mathcal{K} \times \mathcal{U} \times
\mathcal{A} \times \mathcal{M} \to \mathcal{C}$ produces a ciphertext
$C = \Enc_K(U, A, M)$. The decryption algorithm $\Dec : \mathcal{K}
\times \mathcal{U} \times \mathcal{A} \times \mathcal{C} \to
\mathcal{M} \cup \{\bot\}$ returns either a plaintext or the
rejection symbol $\bot$. The \emph{tag length in bytes} is the constant
$\|\Enc_K(U, A, M)\| - \|M\|$, fixed by the scheme.
We write nonce values as $U$ to keep $N$ available for the construction-cost
resource variable used in the concrete bounds.
In this generic syntax, $C$ denotes the full AEAD ciphertext, including
any tag material. The \TW{} algorithm section (\Cref{sec:tw128}) writes this same value as
$C \concat T$, with $C$ the ciphertext body and $T$ the root tag.

\paragraph{Correctness.}
For every $(K, U, A, M) \in \mathcal{K} \times \mathcal{U} \times
\mathcal{A} \times \mathcal{M}$, $\Dec_K(U, A, \Enc_K(U, A, M)) = M$.
An AEAD scheme is \emph{tidy} if every accepting decryption is consistent
with encryption: whenever $\Dec_K(U, A, C) = M \neq \bot$, one has
$\Enc_K(U, A, M) = C$.

\paragraph{All-in-one AE.}
We use the all-in-one authenticated encryption notion of \cite{RS06},
instantiated for nonce-based AE following \cite{BT16}. The experiment
grants either the real pair $(\Enc_K,\Ver_K)$ or the ideal pair
$(\Rand,\Replay)$, where $\Rand$ returns a fresh uniform byte string of
the ciphertext length on each encryption query and $\Replay$ accepts
exactly the tuples $\Rand$ produced. The advantage
$\Advae_\SE(\mathcal{A})$ is the distinguishing advantage between these
two worlds over nonce-respecting adversaries. This all-in-one notion is
instantiated for \TW{} in \Cref{sec:games-pp}.

\paragraph{Integrity under release of unverified plaintext.}
For the RUP setting of \cite{ABLMMY14}, a conventional decryption algorithm is
separated into a plaintext-computing algorithm and a verification algorithm.
The integrity notion \textsf{INT-RUP} gives the adversary access to encryption,
plaintext computation, and verification, and asks whether verification accepts a
fresh ciphertext and tag not returned by encryption. As in the nonce-IV setting
of \cite{ABLMMY14}, encryption queries are nonce-respecting, while plaintext
computation and verification may be queried on arbitrary nonces. This notion
protects only ciphertext integrity; privacy under release of unverified
plaintext requires plaintext awareness in addition to ordinary encryption
privacy.

\subsection{Multi-User Indistinguishability}
\label{sec:mu-ind}

The multi-user indistinguishability notion of \cite[Fig.~3]{BT16}
extends the same real-vs-ideal oracle pair to an adversary-selected set
of users. The number of users $u$ is adversary-determined: it is the
number of $\mathsf{New}$ queries the adversary makes, i.e., the final
value of the user counter $v$ in the game below. A challenge bit $b
\sample \bits$ selects between the
real construction interface (real ciphertexts, real verification) and
an ideal interface (uniform strings of matching length, verification
that accepts only previously-issued encryption tuples). The adversary
also accesses a public primitive interface: the real permutation
$(\pi,\pi^{-1})$ in the real world and the simulator
$(\Sim^{\ROflat},\Sim^{-1,\ROflat})$ in the ideal world; this
world-dependent primitive interface is our modification---\cite{BT16}
work in the ideal-cipher model and give both worlds the same
ideal-cipher oracles---using the same ideal-side simulator convention
as the single-user AE-style games. The adversary
outputs a guess $b'$ and wins if $b' = b$.

\paragraph{Game and advantage.}
The game $\mathsf{G}^{\mathsf{mu\text{-}ind}}_\SE(\mathcal{A})$ of
\cite[Fig.~3]{BT16} samples $b \sample \bits$ and maintains a user
counter $v$, a set $V$ of $(d, U)$ pairs already used by encryption,
and a set $W$ of recorded encryption tuples; its $\mathsf{New}$,
$\mathsf{Enc}$, and $\mathsf{Vf}$ oracles are instantiated for \TW{} in
\Cref{sec:mu-setting}, together with the added public
primitive interface---the real permutation when $b = 1$ and the
simulator when $b = 0$. We write the user index as $d$ and the
associated data input as $A$, where \cite{BT16} writes $i$ and $H$, and
take the absolute value of their signed advantage so that
\[
  \begin{aligned}
  \Advmuind_\SE(\mathcal{A})
  &= \left|2 \Pr[\mathsf{G}^{\mathsf{mu\text{-}ind}}_\SE(\mathcal{A})
          \text{ returns } 1] - 1\right| \\
  &= \bigl|
      \Pr[\mathcal{A}^{\mathsf{Enc}, \mathsf{Vf}, b=1} \Rightarrow 1]
      - \Pr[\mathcal{A}^{\mathsf{Enc}, \mathsf{Vf}, b=0} \Rightarrow 1]
    \bigr|
  \end{aligned}
\]
the second equality holding because $b$ is uniform. Thus
$\Advmuind_\SE$ is the two-game distinguishing advantage between the
real $(\Enc, \Ver)$ pair and the ideal $(\Rand, \Replay)$ pair
instantiated per user. The single-user case $u = 1$ recovers the
all-in-one AE notion of \Cref{sec:ae-defs}, augmented with the public
primitive interface (the real permutation in the real world and the
simulator in the ideal world).

\paragraph{Per-user nonce uniqueness.}
Nonce uniqueness is enforced per user via the set $V$: the same nonce
may appear under distinct users but not twice under the same user.
This is strictly weaker than the global nonce-uniqueness requirement
that would force each nonce value $U$ to appear under at most one user.

\subsection{Committing Notions and Root Tag Hash Security}
\label{sec:commit-defs}

\cite{BH22} formalize a family of committing security notions for
nonce-based AEAD schemes, parameterized by how much of the
encryption input is committed and by whether commitment is
established through decryption (the D-notions) or through encryption
agreement (the E-notions). We recall their definitions and the
$s$-way generalization, and then frame the root tag itself as a hash function
whose \cite{RS04} security---collision, second-preimage, and preimage
resistance---captures the binding and preimage properties expected of a compact
commitment. The AEAD committing notions are then proved from the same tag
analysis in \Cref{sec:root-tag-hash-security}.

\paragraph{What-is-committed functions.}
For $\ell \in \{1, 3, 4\}$, the \emph{what-is-committed} function
$\WiC_\ell$ projects an encryption input $(K, U, A, M) \in \mathcal{K}
\times \mathcal{U} \times \mathcal{A} \times \mathcal{M}$ to the
portion it should commit:
\[
\begin{aligned}
  \WiC_1(K, U, A, M) &= K, \\
  \WiC_3(K, U, A, M) &= (K, U, A), \\
  \WiC_4(K, U, A, M) &= (K, U, A, M).
\end{aligned}
\]

\paragraph{CMTDs-$\ell$ and CMTs-$\ell$ games.}
For $s \geq 2$, \cite[Fig.~5]{BH22} defines the $s$-way committing
games on tuples $(K_i, U_i, A_i, M_i)_{i=1}^s$ with pairwise distinct
$\WiC_\ell$-images, all entries in the scheme's syntactic domain and
none equal to $\bot$. The $s$-way CMTDs-$\ell$ (decryption) game has
the adversary output a ciphertext $C$ alongside the tuples and returns
$1$ if $\Dec_{K_i}(U_i, A_i, C) = M_i$ for every $i$; the printed
return line of \cite[Fig.~5]{BH22} writes $C_i$, but the game outputs a
single ciphertext $C$. The $s$-way CMTs-$\ell$ (encryption) game has the
adversary output the tuples alone and returns $1$ if all encryptions
$\Enc_{K_i}(U_i, A_i, M_i)$ are equal. Advantages $\Advcmtds{\ell}_\SE(\mathcal{A})$ and
$\Advcmts{\ell}_\SE(\mathcal{A})$ are the probabilities that the
adversary wins the respective game; the parameter $s$ is suppressed
from the notation when clear from context. The two-tuple case $s = 2$
is the original CMTD-$\ell$/CMT-$\ell$ pair of \cite[Fig.~4]{BH22}.
These definitions are instantiated for \TW{} in \Cref{sec:games-pp},
with the public permutation lift deferred to that section.

\begin{definition}[Root tag wrapper]
\label{def:htw}
For every \TW{} encryption input $X = (K,U,A,M)$ in the scheme domain,
define the root tag wrapper by
\[
  \HTW(K,U,A,M) = T
  \quad\text{where}\quad
  \Enc_K(U,A,M) = C \concat T.
\]
We also write $\HTW(X)$ for $\HTW(K,U,A,M)$. Equivalently, \HTW{} runs
\TW{} encryption and discards the ciphertext body.
\end{definition}

\paragraph{The root tag as a hash.}
Many uses of commitment store, transmit, or compare a short tag rather
than the full ciphertext, and ask that the short string alone be
binding---the compactly committing AEAD goal introduced by \cite{GLR17}
for message franking. \TW{} ciphertexts are written $C \concat T$ with
$T$ the $\tauroot$-bit root tag, and we capture this goal directly by
treating \Cref{def:htw} as a hash function. Keyed by the permutation,
this is a hash family, and we ask that it satisfy the hash-function
security notions of \cite{RS04}, over the \TW{} input domain, together
with the local $s$-way multicollision extension used for the committing
theorems. For fixed byte lengths $a$ and $\mu$, write
\[
  \mathcal{X}_{a,\mu}
  =
  \bits^k \times \bits^\nu \times (\bits^8)^a \times (\bits^8)^\mu
\]
for the fixed-length slice of valid inputs, with bit length
$m_{a,\mu} = k + \nu + 8a + 8\mu$.
\begin{description}
  \item[Collision resistance (Coll).] No adversary finds distinct inputs
    $X_1, \dots, X_s$ (any $s \geq 2$) with
    $\HTW(X_1) = \cdots = \HTW(X_s)$. The case $s=2$ is the RS04 Coll
    notion; larger $s$ is the corresponding multicollision extension.
  \item[Preimage resistance (Pre).] For fixed $a,\mu$, given
    $\HTW(X^*)$ for a uniform $X^* \sample \mathcal{X}_{a,\mu}$, no
    adversary finds any valid input $X$ with $\HTW(X) = \HTW(X^*)$.
  \item[Second-preimage resistance (Sec/eSec).] The standard and
    everywhere forms of second-preimage resistance of \cite{RS04}, on
    each fixed input slice $\mathcal{X}_{a,\mu}$, follow from
    collision resistance by \cite[Proposition~6]{RS04}.
  \item[Everywhere preimage resistance (ePre).] For a target tag $T^*$ fixed in
    advance, no adversary finds $X$ with $\HTW(X) = T^*$.
\end{description}
The advantages $\AdvColl_\TW$, $\AdvPreSlice{a}{\mu}_\TW$, and
$\AdvePre_\TW$ are the respective success probabilities, instantiated for
\TW{} in \Cref{sec:games-pp}. The Pre guarantee is proved directly in
\Cref{thm:pre}, while the second-preimage guarantees are inherited from Coll
in \Cref{cor:second-preimage-rs04}. Coll is the binding property of the tag
taken as a compact commitment to $(K, U, A, M)$, Pre is the random-opening
preimage property on fixed input slices, and ePre is the fixed-target
preimage property for query-independent tags.
The \cite{BH22} committing notions are the collision-facing AEAD
projection of the same tag analysis: a winning opening gives distinct
inputs sharing a tag, while the dedicated CMTDs-4 proof uses the
all-bodies-equal game structure to avoid the leaf-collision slack of the
general hash collision theorem (\Cref{sec:root-tag-hash-security}).

\paragraph{Relations.}
\cite[Figs.~4 and~5]{BH22} record the implications among the notions.
For the D-notions we use:
\[
  \mathsf{CMTDs}\text{-}4 \;\Rightarrow\;
  \mathsf{CMTDs}\text{-}\ell \text{ for } \ell \in \{1, 3\}
\]
(both immediate: pairwise distinct keys ($\ell = 1$) or pairwise
distinct $(K, U, A)$ triples ($\ell = 3$) are \emph{a fortiori} pairwise
distinct full tuples, so every CMTDs-$\ell$ winner is a CMTDs-4 winner).
The E-notions are bounded directly in \Cref{cor:committing-combined} by
the same final-root candidate-collision argument, rather than through an
encryption reduction.

\section{The TW128 Construction}
\label{sec:tw128}

\subsection{Tree Structure}
\label{sec:design}

\TW{} is a nonce-based AEAD built as a two-level tree of duplex objects over
\Keccakp{}. The \emph{root} duplex is initialized with $(K, U)$, absorbs the
associated data and the first ciphertext chunk, and emits the final
$\tauroot$-bit root tag $T$. The remaining plaintext is partitioned into
\emph{leaf} chunks $M_1,\dots,M_n$; leaf $j$ is encrypted by an independent
duplex initialized with $(K, U, j)$ and returns a $\tauleaf$-bit leaf tag $T_j$.
For multi-chunk messages, the root aggregates
$T_1 \concat \cdots \concat T_n \concat \langle n \rangle_8$ before emitting
$T$. Thus the authentication tag is also the digest point for the full
$(K, U, A, M)$ context. The construction is depicted in \Cref{fig:tw128-tree};
the exact edge-case conventions are specified in \Cref{sec:enc}.

The tree fixes the transcript but not the implementation's parallelism. Leaves
operate on disjoint states after $(K, U)$ is fixed, so scalar, SIMD, and
multicore implementations can evaluate different numbers of leaves at once
while producing identical ciphertexts and tags. The root state and each leaf
state are bounded $1600$-bit objects, and leaf tags can be aggregated as a
stream, giving constant working memory at any fixed degree of parallelism. This
is the AEAD analog of KangarooTwelve's final-node hashing of per-chunk chaining
values, with leaf tags as the chaining values and the root tag as the final
digest.

\begin{figure}[t]
  \centering
  \begin{tikzpicture}[
    font=\small,
    >=Stealth,
    block/.style={draw, rounded corners=1pt, minimum height=10mm,
                  align=center},
    init/.style={block, minimum width=15mm, fill=black!10,
                 font=\footnotesize},
    msg/.style={block, minimum width=11mm, fill=blue!12},
    cv/.style={block, minimum width=9mm, fill=green!18,
               font=\footnotesize},
    enc/.style={block, minimum width=11mm, fill=black!10,
                font=\scriptsize},
    leaf/.style={block, minimum height=22mm, minimum width=11mm,
                 fill=blue!12},
    arr/.style={->, thin},
    flab/.style={font=\scriptsize, inner sep=1pt},
  ]
    \node[init] (init) {$K, U$};
    \node[msg, right=1mm of init] (ad) {$A$};
    \node[msg, right=1mm of ad] (c0) {$C_0$};
    \node[cv, right=1mm of c0] (t1) {$T_1$};
    \node[cv, right=1mm of t1] (t2) {$T_2$};
    \node[right=1mm of t2, inner sep=1pt] (dots) {$\cdots$};
    \node[cv, right=1mm of dots] (tn) {$T_n$};
    \node[enc, right=1mm of tn] (lenc) {$\langle n \rangle_8$};

    \draw[arr] (lenc.east) -- ++(7mm,0) node[right] {$T$};

    \node[leaf, above=12mm of t1] (leaf1) {$C_1$};
    \node[leaf, above=12mm of t2] (leaf2) {$C_2$};
    \node[leaf, above=12mm of tn] (leafn) {$C_n$};

    \node[flab, above=1mm of leaf1.north] {$K, U, 1$};
    \node[flab, above=1mm of leaf2.north] {$K, U, 2$};
    \node[flab, above=1mm of leafn.north] {$K, U, n$};

    \draw[arr] (leaf1.south) -- (t1.north);
    \draw[arr] (leaf2.south) -- (t2.north);
    \draw[arr] (leafn.south) -- (tn.north);
  \end{tikzpicture}
  \caption{Schematic of \TW{} for $n \geq 1$ leaves. The root duplex
    (bottom row) absorbs the associated data, root ciphertext chunk, leaf tags,
    and leaf count before emitting the root tag $T$. Each independent leaf
    duplex (top row) encrypts one leaf chunk and emits a leaf tag.}
  \label{fig:tw128-tree}
\end{figure}

\subsection{Parameters and Domain}
\label{sec:params}

\Cref{tab:params} lists the user-visible and derived scheme parameters.
The implementation uses the same byte length for the root tag and each
leaf tag, but the proofs of \Cref{sec:aead-security,sec:root-tag-hash-security} keep
$\tauroot$ and $\tauleaf$ as separate symbols.

\begin{table}[ht]
  \centering
  \small
  \begin{tabular}{l l l}
    \toprule
    Parameter & Value & Description \\
    \midrule
    $k$                  & $256$          & Key length (bits) \\
    $\nu$                & $256$          & Nonce length (bits) \\
    $\tauroot$ & $256$          & Root tag length (bits) \\
    $\tauleaf$ & $256$          & Leaf tag length (bits) \\
    $\beta$              & $8183$         & Plaintext chunk size (bytes) \\
    $M_{\max}$           & $8183 \cdot 2^{64}$ bytes & Maximum plaintext length \\
    $A_{\max}$           & $8183 \cdot 2^{64}$ bytes & Maximum associated data length \\
    \midrule
    $b$           & $1600$  & Permutation state width (bits) \\
    $n_r$         & $12$    & Permutation rounds \\
    $r$           & $1344$  & Sponge rate (bits) \\
    $c$           & $256$   & Sponge capacity (bits) \\
    $\rho$        & $167$   & Maximum byte-aligned absorption per duplex call (bytes) \\
    \bottomrule
  \end{tabular}
  \caption{\TW{} parameters.}
  \label{tab:params}
\end{table}

The bound $M_{\max} = 8183 \cdot 2^{64}$ bytes is the largest plaintext whose leaf
index fits in $\langle \cdot \rangle_8$: writing $\leaves(M) =
\max(0, \lceil \|M\| / \beta \rceil - 1)$, every
supported plaintext satisfies $0 \leq \leaves(M) \leq 2^{64} - 1$, so
leaf indices lie in $\{1, \dots, 2^{64} - 1\}$. The associated data
bound $A_{\max}$ is a fixed scheme parameter, taken equal to
$M_{\max}$. Inputs exceeding these bounds are outside the
scheme domain.

The derived per-call absorption bound $\rho$ follows the \cite{BDPV11}
analysis: with $r = 1344$ and $\padtenone$ padding, $\rho_{\max}(\padtenone, r) =
r - 2 = 1342$ bits. The byte-aligned data field shares this allotment
with the four-bit phase suffix of \Cref{sec:domain-sep}, so $\rho$ is
the largest integer satisfying $8\rho + 4 \leq 1342$, giving $\rho =
167$ data bytes (and $8\rho + 4 = 1340 \leq 1342$).

\paragraph{Parameter rationale.}
The capacity, key length, and tag lengths are all set to $256$ bits so the
flat sponge loss, key-guessing terms, tag-guessing terms, and root/leaf
collision terms meet a $128$-bit target at the concrete work cap of
\Cref{sec:concrete-claim}. With $b=1600$, this fixes the rate at
$r=1344$. The 12-round \Keccakp{} permutation matches KangarooTwelve
\cite{BDPVVV16}; \Cref{sec:permutation-cryptanalysis} discusses the
cryptanalytic record for this round count. The chunk size is
$\beta = 49\rho = 8183$ bytes, the largest multiple of the byte-aligned
per-call absorption bound not exceeding $2^{13}$ bytes; this fully utilizes
message duplex calls inside full leaves while keeping the leaf working set
small. The $256$-bit nonce is not required by the nonce-respecting bounds, but
it is essentially free in the root initialization call and supports both random
nonces and narrower counter nonces by zero-padding.

\subsection{Domain Separation}
\label{sec:domain-sep}

Distinct duplex instances and phases are separated at the transcript level by
initialization prefixes and phase suffixes. The initialization strings carry a
6-byte ASCII prefix and, for leaves, an eight-byte little-endian chunk index:
\begin{align*}
  \prfxroot &= \texttt{"TW128R"} \in \bits^{48}, \\
  \prfxleaf &= \texttt{"TW128L"} \in \bits^{48}.
\end{align*}
The root's first absorbed string is $\prfxroot \concat K \concat U$, and leaf
$j$'s first absorbed string is
$\prfxleaf \concat K \concat U \concat \langle j \rangle_8$. Every
byte-aligned absorption also carries one of the four-bit phase suffixes in
\Cref{tab:suffixes}, written least significant bit first: absorbing a byte
block $\sigma$ under suffix \sigph{} feeds
$\sigma \concat \sigph$ to the duplex interface of \Cref{sec:sponge} before
$\padtenone$ padding.

\begin{table}[ht]
  \centering
  \small
  \begin{tabular}{l l l}
    \toprule
    Name & Hex & Use \\
    \midrule
    \siginit & \texttt{0xC} & End of initialization (the sole init call) \\
    \sigadmore   & \texttt{0x9} & Non-final associated data block \\
    \sigadlast   & \texttt{0xD} & Final associated data block \\
    \sigmsgmore  & \texttt{0xA} & Non-final message block \\
    \sigmsglast  & \texttt{0xE} & Final message block \\
    \sigaggmore  & \texttt{0xB} & Non-final aggregation block \\
    \sigagglast  & \texttt{0xF} & Final aggregation block (emits root tag) \\
    \bottomrule
  \end{tabular}
  \caption{Phase suffixes (4 bits each, appended LSB-first).}
  \label{tab:suffixes}
\end{table}

\Cref{tab:transcript-schedule} summarizes the root and leaf transcript
schedule. The algorithms of \Cref{sec:enc,sec:dec} give the normative
byte-level specification; the table records which phase is present and what
output the final call of that phase supplies.

\begin{table}[ht]
  \centering
  \scriptsize
  \setlength{\tabcolsep}{1pt}
  \begin{tabular}{@{}p{0.10\linewidth}p{0.11\linewidth}p{0.13\linewidth}p{0.26\linewidth}p{0.24\linewidth}p{0.09\linewidth}@{}}
    \toprule
    Instance & Phase & Suffixes & Absorbed string & Output supplied & When \\
    \midrule
    Root
      & Init
      & \siginit
      & $\prfxroot \concat K \concat U$
      & $0$ if $A \neq \varepsilon$; otherwise first $M_0$ keystream
      & Always \\
    Root
      & AD
      & \sigadmore, \sigadlast
      & $\rho$-byte block partition of $A$
      & Final call emits first $M_0$ keystream
      & $A \neq \varepsilon$ \\
    Root
      & Message
      & \sigmsgmore, \sigmsglast
      & $\rho$-byte block partition of $C_0$
      & Non-final calls emit next keystream; final call emits $T$ if $n = 0$,
        and $0$ otherwise
      & Always \\
    Leaf $j$
      & Init
      & \siginit
      & $\prfxleaf \concat K \concat U \concat \langle j \rangle_8$
      & First $M_j$ keystream
      & $1 \leq j \leq n$ \\
    Leaf $j$
      & Message
      & \sigmsgmore, \sigmsglast
      & $\rho$-byte block partition of $C_j$
      & Non-final calls emit next keystream; final call emits $T_j$
      & $1 \leq j \leq n$ \\
    Root
      & Aggregation
      & \sigaggmore, \sigagglast
      & $T_1 \concat \cdots \concat T_n \concat \langle n \rangle_8$
      & Final call emits $T$
      & $n \geq 1$ \\
    \bottomrule
  \end{tabular}
  \caption{\TW{} transcript schedule. ``First $M_i$ keystream'' means
    $8\min(\|M_i\|,\rho)$ output bits for the first duplex block of that
    message chunk; subsequent message-phase outputs have the length of the next
    duplex block, as in \Cref{alg:helpers}.}
  \label{tab:transcript-schedule}
\end{table}

The combined effect of the initialization prefixes and the phase
suffixes is that the flattened sponge input of any \TW{} duplex call
uniquely determines the duplex instance, the leaf index when applicable,
the phase, and the duplex-call position within that phase.
\Cref{def:wf-transcripts} (Well-formed \TW{} transcript prefixes),
\Cref{lem:transcript-inj} (Transcript Injectivity), and its structural
corollaries formalize this parsing fact for the proofs of
\Cref{sec:aead-security,sec:root-tag-hash-security} and
\Cref{cor:committing-combined}.

\subsection{Encryption}
\label{sec:enc}

\Cref{alg:enc} specifies $\Enc_K(U, A, M)$ using the helper procedures of
\Cref{alg:helpers}. The helpers sit below the tree layer: the encryption and
decryption algorithms split bodies into $\beta$-byte chunks, while the helpers
split a chosen byte string into $\rho$-byte duplex blocks. \Call{StreamEnc}{}
XORs each plaintext block with the current keystream, absorbs the resulting
ciphertext block under the message suffix, and uses the duplex output as the
next keystream. This is the duplex-level presentation of the overwrite
mechanism of \cite[\S6.2 and Algorithm~5]{BDPV11}. The decryption helper
\Call{StreamDec}{} uses the same schedule with ciphertext blocks as input:
it recovers $m_i = c_i \oplus Z$, absorbs the same $c_i$, and returns the
plaintext string and final output value.

\paragraph{Edge-case convention.}
The helper partition always has at least one block: an empty input is the single
empty final block carrying the final phase suffix. The associated-data phase is
present exactly when $A \neq \varepsilon$; when $A = \varepsilon$, the root
\siginit call emits the first root-chunk keystream. The first keystream request
for a message string $X$ is $8\min(\|X\|, \rho)$ bits. If $n = 0$, the root
message phase requests $\tauroot$ bits on its final \sigmsglast call and no
aggregation phase occurs. If $n \geq 1$, that root message call requests zero
bits, each leaf final message call requests $\tauleaf$ bits, and the final
\sigagglast aggregation call requests $\tauroot$ bits. These conventions apply
identically to decryption.

\begin{algorithm}[t]
  \caption{TW128 helper procedures.}
  \label{alg:helpers}
  \begin{algorithmic}[1]
    \Procedure{Absorb}{$\mathcal{D}$, $s$, $\sigma^-$, $\sigma^+$, $\ell$}
      \State Partition $s$ into duplex blocks $s_1, \dots, s_t$ of exactly
        $\rho$ bytes each, except that the last block may be shorter.
      \State Require $t \geq 1$; the empty case yields the single empty
        duplex block $s_1 = \varepsilon$.
      \For{$i \gets 1$ \textbf{to} $t - 1$}
        \State $\mathcal{D}.\duplexing(s_i \concat \sigma^-, 0)$
      \EndFor
      \State \Return $\mathcal{D}.\duplexing(s_t \concat \sigma^+, \ell)$
    \EndProcedure
    \Statex
    \Procedure{StreamEnc}{$\mathcal{D}$, $m$, $Z$, $\sigma^-$, $\sigma^+$, $\ell$}
      \State Partition $m$ into duplex blocks $m_1, \dots, m_t$ of exactly
        $\rho$ bytes each, except that the last block may be shorter.
      \State Require $t \geq 1$; if $m = \varepsilon$, set $t = 1$ and
        $m_1 = \varepsilon$.
      \For{$i \gets 1$ \textbf{to} $t - 1$}
        \State $c_i \gets m_i \oplus Z$
        \State $Z \gets \mathcal{D}.\duplexing(c_i \concat \sigma^-, 8\|m_{i+1}\|)$
      \EndFor
      \State $c_t \gets m_t \oplus Z$
      \State $Z' \gets \mathcal{D}.\duplexing(c_t \concat \sigma^+, \ell)$
      \State \Return $(c_1 \concat \cdots \concat c_t, Z')$
    \EndProcedure
  \end{algorithmic}
\end{algorithm}

\begin{algorithm}[t]
  \caption{$\Enc_K(U, A, M)$.}
  \label{alg:enc}
  \begin{algorithmic}[1]
    \Require $K \in \bits^k$, $U \in \bits^{\nu}$, $A$ with $\|A\| \leq A_{\max}$, $M$ with $\|M\| \leq M_{\max}$.
    \Ensure $C \concat T$ with $\|C\| = \|M\|$ and $T \in \bits^{\tauroot}$.
    \State $M_0 \gets$ first $\min(\|M\|, \beta)$ bytes of $M$ (possibly $\varepsilon$).
    \State Partition the remaining suffix into chunks $M_1, \dots, M_n$ of exactly $\beta$ bytes each, except that the last chunk may be shorter but nonempty, so $n = \leaves(M)$.
    \State $\mathcal{D}_{\text{root}} \gets$ new $\Duplex$
    \If{$A = \varepsilon$}
      \State $Z \gets \mathcal{D}_{\text{root}}.\duplexing(\prfxroot \concat K \concat U \concat \siginit, 8\min(\|M_0\|, \rho))$
    \Else
      \State $\mathcal{D}_{\text{root}}.\duplexing(\prfxroot \concat K \concat U \concat \siginit, 0)$
      \State $Z \gets \Call{Absorb}{\mathcal{D}_{\text{root}}, A, \sigadmore, \sigadlast, 8\min(\|M_0\|, \rho)}$
    \EndIf
    \If{$n = 0$}
      \State $(C_0,\, T) \gets \Call{StreamEnc}{\mathcal{D}_{\text{root}}, M_0, Z, \sigmsgmore, \sigmsglast, \tauroot}$
      \State \Return $C_0 \concat T$
    \EndIf
    \State $(C_0,\, \_) \gets \Call{StreamEnc}{\mathcal{D}_{\text{root}}, M_0, Z, \sigmsgmore, \sigmsglast, 0}$
    \For{$j \gets 1$ \textbf{to} $n$ \textbf{in parallel}}
      \State $\mathcal{D}_j \gets$ new $\Duplex$
      \State $Z_j \gets \mathcal{D}_j.\duplexing(\prfxleaf \concat K \concat U \concat \langle j \rangle_8 \concat \siginit, 8\min(\|M_j\|, \rho))$
      \State $(C_j, T_j) \gets \Call{StreamEnc}{\mathcal{D}_j, M_j, Z_j, \sigmsgmore, \sigmsglast, \tauleaf}$
    \EndFor
    \State $G \gets T_1 \concat \cdots \concat T_n$.
    \State $T \gets \Call{Absorb}{\mathcal{D}_{\text{root}},\, G \concat \langle n \rangle_8,\, \sigaggmore,\, \sigagglast,\, \tauroot}$
    \State \Return $C_0 \concat C_1 \concat \cdots \concat C_n \concat T$
  \end{algorithmic}
\end{algorithm}

The leaf loop is annotated \textbf{in parallel} to make explicit that
the $n$ leaf computations are mutually independent; \Cref{sec:impl}
discusses the SIMD realizations. The annotation is an implementation
freedom, not a mode parameter: evaluating leaves serially, in SIMD
batches, or across cores leaves the transcript and wire format unchanged.

\paragraph{Streaming aggregation.}
The string $G = T_1 \concat \cdots \concat T_n$ in
\Cref{alg:enc,alg:dec} names the leaf tags only for specification. The
closing \Call{Absorb}{} consumes $G \concat \langle n \rangle_8$ as
$\rho$-byte duplex blocks in leaf order and advances the root state
after each, so an implementation need not hold $G$ in full. It can absorb
leaf tags in order as they become available, or buffer completed
out-of-order leaves until the next in-order tag or SIMD batch can be
absorbed, keeping only a bounded number of tags at a time. With
the bounded $1600$-bit state of every duplex, this keeps the working
memory constant in $\|M\|$ for a fixed degree of parallelism, growing
only with the number of leaves evaluated at once rather than with the
message length; \Cref{sec:impl} describes the batched realization.

\subsection{Decryption and Correctness}
\label{sec:dec}

\Cref{alg:dec} specifies $\Dec_K(U, A, C \concat T)$. Decryption mirrors
encryption at the duplex block level: each ciphertext block is absorbed under
the same phase suffix, and the plaintext block is recovered as
$m_i = c_i \oplus Z^{(i)}$ from the keystream value that encryption would have
used at that position. The edge cases are those of \Cref{sec:enc}. The
procedure absorbs all in-domain ciphertext bytes before comparing the final tag,
so the verification transcript is determined by $(K, U, A, C)$ whether or not the
submitted tag accepts.

\begin{algorithm}[t]
  \caption{$\Dec_K(U, A, C \concat T)$.}
  \label{alg:dec}
  \begin{algorithmic}[1]
    \Require $K, U, A$ as in \Cref{alg:enc}, ciphertext $C$ with $\|C\| \leq M_{\max}$, candidate tag $T \in \bits^{\tauroot}$.
    \Ensure $M$ with $\|M\| = \|C\|$, or $\bot$.
    \State $C_0 \gets$ first $\min(\|C\|, \beta)$ bytes of $C$ (possibly $\varepsilon$).
    \State Partition the remaining suffix into chunks $C_1, \dots, C_n$ of exactly $\beta$ bytes each, except that the last chunk may be shorter but nonempty, so $n = \leaves(C)$.
    \State $\mathcal{D}_{\text{root}} \gets$ new $\Duplex$
    \If{$A = \varepsilon$}
      \State $Z \gets \mathcal{D}_{\text{root}}.\duplexing(\prfxroot \concat K \concat U \concat \siginit, 8\min(\|C_0\|, \rho))$
    \Else
      \State $\mathcal{D}_{\text{root}}.\duplexing(\prfxroot \concat K \concat U \concat \siginit, 0)$
      \State $Z \gets \Call{Absorb}{\mathcal{D}_{\text{root}}, A, \sigadmore, \sigadlast, 8\min(\|C_0\|, \rho)}$
    \EndIf
    \If{$n = 0$}
      \State $(M_0,\, T') \gets \Call{StreamDec}{\mathcal{D}_{\text{root}}, C_0, Z, \sigmsgmore, \sigmsglast, \tauroot}$
      \State \Return $M_0$ if $T' = T$ (compared in constant time), else $\bot$.
    \EndIf
    \State $(M_0,\, \_) \gets \Call{StreamDec}{\mathcal{D}_{\text{root}}, C_0, Z, \sigmsgmore, \sigmsglast, 0}$
    \For{$j \gets 1$ \textbf{to} $n$ \textbf{in parallel}}
      \State $\mathcal{D}_j \gets$ new $\Duplex$
      \State $Z_j \gets \mathcal{D}_j.\duplexing(\prfxleaf \concat K \concat U \concat \langle j \rangle_8 \concat \siginit, 8\min(\|C_j\|, \rho))$
      \State $(M_j, T_j) \gets \Call{StreamDec}{\mathcal{D}_j, C_j, Z_j, \sigmsgmore, \sigmsglast, \tauleaf}$
    \EndFor
    \State $G \gets T_1 \concat \cdots \concat T_n$.
    \State $T' \gets \Call{Absorb}{\mathcal{D}_{\text{root}},\, G \concat \langle n \rangle_8,\, \sigaggmore,\, \sigagglast,\, \tauroot}$
    \If{$T' = T$ (compared in constant time)}
      \State \Return $M_0 \concat M_1 \concat \cdots \concat M_n$
    \Else
      \State \Return $\bot$
    \EndIf
  \end{algorithmic}
\end{algorithm}

\paragraph{Correctness.}
For every $K \in \bits^k$, $U \in \bits^{\nu}$,
$A$ with $\|A\| \leq A_{\max}$, and $M$ with $\|M\| \leq M_{\max}$,
$\Dec_K(U, A, \Enc_K(U, A, M)) = M$. Encryption and decryption absorb the same
ciphertext blocks under the same duplex instances and phase suffixes, so their
duplex states and final root tags agree. Each recovered block satisfies
$c_i \oplus Z^{(i)} = (m_i \oplus Z^{(i)}) \oplus Z^{(i)} = m_i$.

\begin{lemma}[\TW{} tidiness]
\label{lem:tw-tidiness}
For every valid \TW{} decryption input, if
$\Dec_K(U, A, C \concat T) = M \neq \bot$, then
$\Enc_K(U, A, M) = C \concat T$.
\end{lemma}

\begin{proof}
An accepting decryption has $\|M\| = \|C\|$ and therefore the same canonical
chunking as re-encryption of $M$. The decryption transcript absorbs the same
$(K, U, A)$ values and ciphertext blocks that this re-encryption absorbs, under
the same duplex instances and phase suffixes. Its stream outputs recover
$m_i = c_i \oplus Z^{(i)}$, so re-encryption emits the submitted ciphertext blocks
again. The recomputed leaf tags and final root candidate tag are therefore the
same values; since decryption accepts only when the final candidate equals the
submitted tag $T$, re-encryption returns $C \concat T$.
\end{proof}

\section{Security Model}
\label{sec:security-model}

This section instantiates the security claims for \TW{}. The root tag is first
treated as the output of the wrapper $\HTW$: adversaries supply full inputs or
candidate openings, and the games ask whether the tag collides or can be hit as
a preimage. The same construction is also treated as a nonce-based AEAD: the
\textsf{INT-RUP} game gives integrity even when a plaintext-computing interface
is exposed, while all-in-one AE and mu-ind record the baseline
confidentiality and authenticity notions.

\Cref{tab:security-notions} summarizes the proof obligations before the formal
games. The loss terms named there are defined in \Cref{sec:loss-terms}, and the
implemented-parameter instantiations appear in \Cref{sec:concrete-security}.

\paragraph{Proof worlds.}
All final statements are public-permutation claims for \TW{} with \Keccakp{}
modeled as a uniformly random permutation. The proofs first establish random
oracle bounds for the flat duplex interface of \Cref{sec:ro-model}; the flat
sponge lift of \Cref{app:flat-sponge-lift} then transfers those bounds to the
public-permutation games of \Cref{sec:games-pp,sec:mu-setting}.

\paragraph{Scope of the security claims.}
In the AE-style distinguishing games, encryption queries are nonce-respecting
and the verification oracle returns only acceptance or rejection. The
\textsf{INT-RUP} game gives integrity under a plaintext-computing oracle, but it
does not give privacy under release of unverified plaintext or
plaintext-awareness security. None of the AE-style games gives nonce-misuse
resistance. The hash and committing games are adversarial-initialization games
over submitted inputs $(K, U, A, M)$ or ciphertext openings and impose no nonce
uniqueness condition. Across all games, the stated bounds are birthday-type,
single-stage, and leakage-free: beyond-birthday and side-channel security are
outside the model.
\Cref{sec:limitations} discusses the consequences of these exclusions and the
corresponding open problems, and \Cref{sec:permutation-cryptanalysis} reviews
the third-party cryptanalysis supporting the random-permutation modeling of
\Keccakp{}.

\begin{table}[t]
  \centering
  \footnotesize
  \setlength{\tabcolsep}{3pt}
  \begin{tabular}{@{}p{0.18\linewidth}p{0.34\linewidth}p{0.23\linewidth}p{0.17\linewidth}@{}}
    \toprule
    Notion & Adversary goal & Main result & Loss shape \\
    \midrule
    $\HTW$ collision
      & Find distinct full inputs $(K, U, A, M)$ with the same root tag.
      & \Cref{thm:coll}
      & Capacity lift, root collision, leaf collision \\
    $\HTW$ preimage
      & Find an input hitting a fixed target tag.
      & \Cref{thm:epre}
      & Capacity lift, root preimage \\
    $\HTW$ second preimage
      & Find a second preimage in the standard or everywhere
        \cite{RS04} sense.
      & \Cref{cor:second-preimage-rs04}
      & Inherited from collision resistance \\
    CMT / CMTD
      & Produce equal encryptions or one valid ciphertext opening under
        distinct committed inputs.
      & \Cref{cor:cmtds4,cor:committing-combined}
      & Capacity lift, root collision \\
    \textsf{INT-RUP}
      & Produce a fresh accepting ciphertext after encryption and
        plaintext-computing queries.
      & \Cref{thm:int-rup}
      & Capacity lift, key guess, leaf collision, tag guess \\
    AE / mu-ind
      & Distinguish real encryption and verification from random ciphertexts
        and replay-only verification, with nonce-respecting encryption queries.
      & \Cref{thm:mu-ind,cor:ae}
      & Capacity lift, key guess, leaf collision, tag guess \\
    \bottomrule
  \end{tabular}
  \caption{Proof obligations and the dominant loss shapes used for \TW{}.
    The table is only a roadmap; the formal games and resource accounting are
    defined in \Cref{sec:games-pp,sec:mu-setting,sec:resources}.}
  \label{tab:security-notions}
\end{table}

The proof dependency spine is as follows. \Cref{app:lemmas} establishes
transcript decoding, bridge injectivity, output-point separation, and the local
probabilistic accounting lemmas for key tests, leaf tag collisions, and root tag
guesses. \Cref{sec:root-tag-hash-security} uses those structural facts to build
a chronological sequence of final root candidates and apply a single root tag
hash bound; the hash and committing theorems differ only in which collision or
preimage event is reduced to that candidate sequence. \Cref{sec:aead-security}
uses the same facts for the random oracle mu-ind and \textsf{INT-RUP} bounds.
\Cref{app:flat-sponge-lift} transfers both branches to the public permutation
games defined below.

\subsection{Public Permutation Games}
\label{sec:games-pp}

\paragraph{Primitive convention.}
In every public permutation game below, the challenger samples
$\pi \sample \Perm(\bits^{1600})$ and exposes the direct primitive
interface $(\pi, \pi^{-1})$ to the adversary. The real-vs-ideal AE-style
games compare the real construction interface with the ideal interface
specified by the game over this shared primitive, following the
same-oracle convention of \cite{BT16}; the hash and committing games
below are one-world success probabilities over the same sampled
permutation. The simulator $(\Sim^{\ROflat},\Sim^{-1,\ROflat})$ of
\Cref{app:flat-sponge-lift} is a proof device, entering through the
lift from the intermediate random oracle model.

\paragraph{Validity convention.}
All construction interface inputs are checked against \TW{}'s public
syntax and domain before any construction computation or
construction-internal primitive lookup is performed. An encryption input
is valid when $U \in \bits^\nu$, $A$ is a byte string with
$\|A\| \leq A_{\max}$, and $M$ is a byte string with
$\|M\| \leq M_{\max}$. A verification or decryption input is valid when
$U \in \bits^\nu$, $A$ is a byte string with $\|A\| \leq A_{\max}$, and
the full ciphertext parses as $C \concat T$ with $\|C\| \leq M_{\max}$
and $T \in \bits^{\tauroot}$. In the hash and committing games, each
submitted input $(K, U, A, M)$ must be a valid encryption input including
$K \in \bits^k$, and a submitted CMTDs ciphertext must be a valid
verification/decryption ciphertext. Invalid encryption queries
are rejected and not recorded, invalid plaintext-computing queries return
$\bot$, invalid verification queries return $0$, and invalid committing-game
outputs make the game return $0$. Such rejected inputs contribute no
construction cost and create no construction transcripts. Throughout, an
\emph{accepted} encryption query is one that passes this validity check and any
game-specific nonce or user check.

\paragraph{Construction oracles.}
Whenever a game uses a secret key $K$ and permutation $\pi$, $\Enc_K[\pi]$
and $\Dec_K[\pi]$ are the algorithms of \Cref{alg:enc,alg:dec} with that
permutation. The verification oracle is
$\Ver_K[\pi](U, A, C, T) = [\Dec_K[\pi](U, A, C \concat T) \neq \bot]$.
For \textsf{INT-RUP}, $\PDec_K[\pi](U, A, C, T)$ is the plaintext recovered by
the decryption transcript before the final root-tag comparison; it returns
neither the recomputed tag nor the verification bit.

\paragraph{Hash games.}
The hash games for the wrapper $\HTW$ (\Cref{def:htw}) have no
challenger-sampled key. In the $s$-way collision game, all inputs are
adversary-supplied: the adversary outputs $s$ inputs, and the challenger
returns $1$ iff they are pairwise distinct, in the \TW{} domain, and share
a root tag:
\begin{align*}
  &\pi \sample \Perm(\bits^{1600}), \quad
   (K_i, U_i, A_i, M_i)_{i=1}^s \sample \mathcal{A}^{\pi,\, \pi^{-1}}, \\
  &\text{return $0$ unless the inputs are pairwise distinct and in the \TW{} domain}, \\
  &\text{return } \bigl[\HTW(K_1, U_1, A_1, M_1) = \cdots =
     \HTW(K_s, U_s, A_s, M_s)\bigr],
\end{align*}
with advantage $\AdvColl_\TW(\mathcal{A}) = \Pr[\text{the game returns }
1]$; no nonce uniqueness is imposed. The \cite{RS04} standard and
everywhere second-preimage notions are used on fixed input slices; in
\Cref{cor:second-preimage-rs04} they are reduced explicitly to this
two-way collision game.
For fixed byte lengths $a,\mu$, the standard preimage game samples
$X^* \sample \mathcal{X}_{a,\mu}$ independently of the primitive, computes
$T^* = \HTW(X^*)$, gives $T^*$ to
$\mathcal{A}^{\pi,\, \pi^{-1}}$, and returns $1$ iff $\mathcal{A}$ outputs a
valid input $X$ in the \TW{} domain with $\HTW(X) = T^*$; the advantage is
$\AdvPreSlice{a}{\mu}_\TW(\mathcal{A})$. The
everywhere preimage game fixes a target tag
$T^* \in \bits^{\tauroot}$ before the primitive sampling and returns $1$
iff $\mathcal{A}(T^*)$ outputs an input $X$ in the domain with
$\HTW(X) = T^*$, the advantage $\AdvePre_\TW$ taken as the worst case over
query-independent targets.

\paragraph{Committing games.}
The \cite{BH22} committing games are instantiated as in
\Cref{sec:commit-defs}, with no challenger-sampled key. The $s$-way
CMTDs-$\ell$ (decryption) game has the adversary output a ciphertext
$C \concat T$ and $s$ tuples $(K_i, U_i, A_i, M_i)$ with pairwise distinct
$\WiC_\ell$-images, and returns $1$ iff
$\Dec_{K_i}[\pi](U_i, A_i, C \concat T) = M_i$ for every $i$, with the eager
length rejection justified by \Cref{alg:dec}
($\Dec_K(U, A, C \concat T)$ returns $\bot$ or a message of length $\|C\|$);
its advantage is
$\Advcmtds{\ell}_\TW(\mathcal{A})$. The CMTs-$\ell$ (encryption) game has
the adversary output the $s$ tuples alone and returns $1$ iff the
encryptions $\Enc_{K_i}[\pi](U_i, A_i, M_i)$ are all equal; its advantage is
$\Advcmts{\ell}_\TW(\mathcal{A})$.

\paragraph{RUP integrity.}
The \textsf{INT-RUP} game samples $K \sample \bits^k$ and
$\pi \sample \Perm(\bits^{1600})$, exposes $\Enc_K[\pi]$,
$\PDec_K[\pi]$, $\Ver_K[\pi]$, and $(\pi, \pi^{-1})$, and maintains the replay
table $W$ of encryption outputs. Accepted encryption queries are
nonce-respecting: an encryption query is rejected if its nonce has appeared in a
previous accepted encryption query. Plaintext-computing and verification
queries may use arbitrary nonces. On an accepted encryption query $(U, A, M)$ the
game returns $\Enc_K[\pi](U, A, M)$, parsed and recorded as $(U, A, C, T)$ in $W$.
On a valid plaintext-computing query $(U, A, C, T)$ it returns
$\PDec_K[\pi](U, A, C, T)$, and on an invalid query it returns $\bot$. On a valid
verification query $(U, A, C, T)$ it returns $\Ver_K[\pi](U, A, C, T)$; if this bit
is $1$ and $(U, A, C, T) \notin W$ when the query is asked, the game records a
forgery. Invalid verification queries return $0$. The game returns $1$ iff a
forgery is recorded, and
\[
  \AdvintRUP_\TW(\mathcal{A}) =
  \Pr[\mathsf{G}^{\mathsf{INT\text{-}RUP}}_\TW(\mathcal{A}) \Rightarrow 1].
\]

\paragraph{All-in-one AE.}
The all-in-one authenticated encryption game~\cite{RS06} samples
$K \sample \bits^k$ and compares the real pair
$(\Enc_K[\pi], \Ver_K[\pi])$ with the ideal pair $(\Rand, \Replay)$.
On each accepted encryption query $(U, A, M)$, $\Rand$ samples a uniformly
random byte string of length $\|M\| + \tauroot/8$, parses it as
$C \concat T$ with $T \in \bits^{\tauroot}$, records $(U, A, C, T)$, and
returns $C \concat T$; $\Replay$ accepts exactly recorded tuples. Encryption
queries are nonce-respecting, and
\[
  \Advae_\TW(\mathcal{A}) =
  \bigl|
    \Pr[\mathcal{A}^{\Enc_K[\pi],\, \Ver_K[\pi],\, \pi,\, \pi^{-1}} \Rightarrow 1]
    - \Pr[\mathcal{A}^{\Rand,\, \Replay,\, \pi,\, \pi^{-1}} \Rightarrow 1]
  \bigr|.
\]
This all-in-one AE advantage is the single-user case of the multi-user
real-vs-ideal notion of \Cref{sec:mu-ind}; \Cref{cor:ae} derives it from
\Cref{thm:mu-ind}.

\subsection{Multi-User Setting}
\label{sec:mu-setting}

This subsection instantiates the mu-ind game of \cite{BT16}
(\Cref{sec:mu-ind}) for \TW{} in the public permutation model, under
the primitive convention of \Cref{sec:games-pp}. The challenger
samples a challenge bit $b \sample \bits$ and a permutation
$\pi \sample \Perm(\bits^{1600})$, and exposes $(\pi, \pi^{-1})$ as the
primitive interface in both worlds. The
adversary $\mathcal{A}$ creates users on demand via $\mathsf{New}$, and
the game maintains an active-user counter $v$. We write $D$ for an
execution-uniform cap on the number of $\mathsf{New}$ calls. Equivalently,
the game may sample independent keys $K_1,\dots,K_D$ at the start of the
experiment and let the $d$-th $\mathsf{New}$ call activate $K_d$; inactive
indices are never valid users.

The mu-ind construction oracles inherit the validity convention of
\Cref{sec:games-pp}. The encryption oracle on input $(d, U, A, M)$
rejects, without recording, if $(U, A, M)$ is outside the \TW{} encryption
domain, if $d \notin \{1, \dots, v\}$, or if $(d, U) \in V$; otherwise
it sets $C = \Enc_{K_d}[\pi](U, A, M)$ when $b = 1$, and samples a
uniformly random byte string $C$ of length $\|M\| + \tauroot/8$ when
$b = 0$. The game records
$(d, U)$ in $V$ and $(d, U, A, C)$ in $W$, and returns $C$. The
verification oracle on input $(d, U, A, C)$ rejects if $d \notin
\{1, \dots, v\}$ or if $(U, A, C)$ is not a valid \TW{}
verification/decryption input; if $(d, U, A, C) \in W$ it returns
$\mathsf{true}$; if $b = 0$ it returns $\mathsf{false}$; otherwise it
returns $[\Dec_{K_d}[\pi](U, A, C) \neq \bot]$. Nonce uniqueness is
enforced per user by the set $V$.

After querying these oracles, $\mathcal{A}$ outputs $b' \in \bits$.
The mu-ind advantage of $\mathcal{A}$ against \TW{} is
\[
  \Advmuind_\TW(\mathcal{A})
  = |2 \Pr[b' = b] - 1|.
\]
A valid verification query is \emph{non-replay} if its tuple is not in
$W$ when the query is asked; replay hits are answered by the recorded
transcript and are not counted in the verification-query resource below.

Resource caps split per potential user $d \leq D$: $N^{(d)}$,
$\qv^{(d)}$, $\nleaf^{(d)}$ are the per-user counts under
\Cref{sec:resources}, with inactive users contributing zero. The global
totals $N = \sum_{d=1}^D N^{(d)}$, $\qv = \sum_{d=1}^D \qv^{(d)}$,
$\nleaf = \sum_{d=1}^D \nleaf^{(d)}$ appear in the final bound: $N$
through the lift term $\Lift_c(\Ntot)$, and $\nleaf$ and $\qv$ through
the leaf-collision and verification-tag terms. $\Nprim$
remains a single shared count over the primitive interface.

The mu-ind treatment applies only to indistinguishability. The
committing security games already give the adversary control of all
$s$ keys, so the $s$-way CMTDs-$\ell$ statements of
\Cref{sec:games-pp} are themselves the multi-key committing statement
and need no separate $\mathsf{mu\text{-}cmt}$ notion.

\subsection{The TW128 Random Oracle Model}
\label{sec:ro-model}

The proofs of \Cref{sec:aead-security,sec:root-tag-hash-security} and
\Cref{cor:committing-combined} use an intermediate random oracle model. This
model keeps the same construction-level games but replaces every duplex output
by a lookup to a single random oracle through the \cite{BDPV11} bridge of
\Cref{sec:sponge}. The flat sponge lift of \Cref{app:flat-sponge-lift} then
converts the resulting bounds to the public-permutation form.

The model has a single random oracle $\ROflat : \bits^* \to
\bits^\infty$ over the unpadded $H$-input domain of the \cite{BDPV08}
simple-padded sponge interface, with lazy sampling: each distinct
input receives an independent infinite bit string, and repeated queries
return the same string. A direct adversary query is finite-output: on
input $(Y, \ell)$ with finite $\ell$, the oracle returns the first
$\ell$ bits of the lazy-sampled value for $Y$.

Construction calls do not query $\ROflat$ on arbitrary strings. Every root
or leaf duplex output is read from $\ROflat$ at the bridged flattened
transcript prefix for that duplex call, and the caller takes only the
requested output projection. Direct adversary queries to $\ROflat$ are the
only random oracle queries counted as direct oracle queries; construction
lookups are accounted for through the construction cost $N$ of
\Cref{sec:resources}.

For a well-formed \TW{} duplex transcript prefix $X$ in the language of
\Cref{def:wf-transcripts}, at one of the construction transcript lookup
points formalized by \Cref{lem:bridge-decode,tab:keyed-twout}, let
$\flatmap(X) = \bridge[1344, 1344](\mathsf{raw\_flat}(X))$ be the
bridged image of the native flattened input under the
\cite[Theorem~3]{BDPV11} preprocessing map of \Cref{sec:sponge}, and define
\[
  \RF(X) = \ROflat(\flatmap(X)).
\]
The construction interfaces $\Enc_K^{\RO}$, $\Dec_K^{\RO}$,
$\PDec_K^{\RO}$, and $\Ver_K^{\RO}$ are the \TW{} algorithms of
\Cref{sec:tw128} with every root or leaf duplex output replaced by the
corresponding requested projection of $\RF(X)$; the projection length is the
duplexing call's output parameter $\ell$, which may be zero. The bridge and the
typed-restriction convention are formalized in
\Cref{lem:bridge-decode}; the direct-query key-hit predicate ignores
the adversary's requested output length and is formalized as
$\KeyedTWOut$ in \Cref{lem:bridge-decode}.

The random oracle games are obtained from
\Cref{sec:games-pp,sec:mu-setting} by this mechanical substitution: direct
primitive access $(\pi, \pi^{-1})$ becomes finite-output access to $\ROflat$,
real construction or deterministic challenger checks use the corresponding
${}^\RO$ algorithm, and the ideal oracles $(\Rand,\Replay)$ are unchanged.
We use the same advantage names with subscript $\RO$: $\AdvColl_{\RO}$,
$\AdvPreSlice{a}{\mu}_{\RO}$, $\AdvePre_{\RO}$,
$\Advcmtds{\ell}_{\RO}$, $\Advcmts{\ell}_{\RO}$, and
$\AdvintRUP_{\RO}$ are the win probabilities of the corresponding random
oracle games, with CMTs challengers comparing deterministic encryptions via
$\Enc^\RO$. The all-in-one AE advantage $\Adv^{\mathsf{ae}}_{\RO}$ is the
corresponding real-vs-ideal distinguishing advantage with direct $\ROflat$
access in both worlds.

For mu-ind, $\Adv^{\mathsf{mu\text{-}ind}}_{\RO}(\mathcal{B})$ denotes the
\Cref{sec:mu-setting} game with the same $\mathsf{New}$ interface and per-user
nonce discipline: user $d$ uses
$(\Enc_{K_d}^{\RO}, \Ver_{K_d}^{\RO})$ when $b = 1$ and $(\Rand,\Replay)$ when
$b = 0$. Equivalently, in the two-game notation of \Cref{sec:mu-ind},
\[
  \Adv^{\mathsf{mu\text{-}ind}}_{\RO}(\mathcal{B}) =
  \bigl|
    \Pr[\mathcal{B}^{\mathsf{Real}^{\mathsf{mu}}_{\RO},\, \ROflat} \Rightarrow 1]
    - \Pr[\mathcal{B}^{\mathsf{Ideal}^{\mathsf{mu}},\, \ROflat} \Rightarrow 1]
  \bigr|,
\]
where the superscripts name the entire multi-user construction
interface in the corresponding world.
Adversary queries to $\ROflat$ are finite-output and counted in
$\qRO(\mathcal{B})$ of \Cref{sec:resources}. Construction-internal
$\RF(\cdot)$ lookups (those made by $\Enc_K^\RO$, $\Dec_K^\RO$,
$\PDec_K^\RO$, $\Ver_K^\RO$, or a hash or committing challenger) are not.

\subsection{Resources}
\label{sec:resources}

Adversary resources are measured by fixed, execution-uniform caps on
the number and size of oracle queries. An execution-uniform cap is one
that holds on every execution path, including over adversary
randomness, oracle randomness, and adaptive query choices.
Proofs may also introduce local execution-uniform caps, such as
$q_e$ for the number of accepted encryption queries, solely to index a
finite hybrid sequence. Such local caps are part of the adversary's
fixed query bounds, but they are not listed as headline resource parameters
unless they affect the final bound.

\Cref{tab:resource-accounting} gives the compact notation summary used in later
theorems. Lowercase resource quantities are realized counts in a random oracle
experiment, while uppercase quantities are execution-uniform caps used in the
public-permutation bounds. The detailed definitions below fix exactly which
computations contribute to each count, and \Cref{sec:loss-terms} defines the
closed-form loss terms.

\begin{table}[t]
  \centering
  \footnotesize
  \setlength{\tabcolsep}{3pt}
  \begin{tabular}{@{}p{0.18\linewidth}p{0.42\linewidth}p{0.32\linewidth}@{}}
    \toprule
    Symbol & Meaning & Primary use / definition \\
    \midrule
    $N$
      & Construction-internal \Keccakp{} calls induced by valid construction
        oracle queries and deterministic challenger checks.
      & Public-permutation lift and concrete work cap. \\
    $\Nprim$
      & Direct adversary queries to the public primitive interface.
      & Key-guess, root-collision, root-preimage terms; direct-query part of
        the lift. \\
    $\Ntot = N + \Nprim$
      & Converted primitive-query cost used by the flat sponge replacement.
      & Argument of $\Lift_c(\Ntot)$. \\
    $\qRO(\mathcal{B})$, $\QRO$
      & Direct finite-output queries to $\ROflat$ in the intermediate random
        oracle model; construction lookups are excluded.
      & Random oracle bounds; after the lift $\qRO \leq \Nprim$. \\
    $\qv(\mathcal{A})$, $\Qv$
      & Valid verification queries that are not replay-table hits.
      & Verification-tag terms in AE, mu-ind, and \textsf{INT-RUP} bounds. \\
    $\nleaf(\mathcal{B})$, $\Nleaf$
      & Distinct construction-generated final leaf transcripts.
      & Leaf tag collision and tag-guessing terms; always at most $N$. \\
    $D$
      & Execution-uniform cap on users activated in the mu-ind game.
      & Multi-user key-collision and key-guess terms. \\
    \midrule
    $\Lift_c(\Ntot)$
      & \cite{BDPV08} flat sponge differentiating loss.
      & Defined in \cref{sec:loss-terms}. \\
    $\KeyGuess_k(Q)$
      & Ideal-system key-guessing term.
      & Defined in \cref{sec:loss-terms}. \\
    $\RootColl_\tau(Q, s)$
      & $s$-way root-tag collision term.
      & Defined in \cref{sec:loss-terms}. \\
    $\RootPre_\tau(Q)$
      & Root-tag preimage term.
      & Defined in \cref{sec:loss-terms}. \\
    $\LeafColl_\tau(Q)$
      & Known-key leaf-tag collision term.
      & Defined in \cref{sec:loss-terms}. \\
    \bottomrule
  \end{tabular}
  \caption{Resource caps and recurring loss terms. Uppercase resource symbols
    are the caps that appear in public-permutation theorems; lowercase symbols
    are realized counts in random oracle executions or source games.}
  \label{tab:resource-accounting}
  \label{tab:notation}
\end{table}

\begin{description}
\item[$N$:] the construction interface cost, measured as the number of
  \Keccakp{} calls induced after the validity convention has been applied.
  It counts at per-duplex-call granularity, summing each root initialization,
  associated data duplex block, root message duplex block, leaf initialization,
  leaf message duplex block, and aggregation duplex block.

  \emph{AE-style games.} In real games, $N$ counts the valid construction
  computations actually run by the adversary's oracle queries. In ideal games,
  the same accounting applies to the corresponding real \TW{} computation on
  the valid query, even when the ideal oracle answers by table lookup or without
  running the construction. Valid plaintext-computing computations and valid
  AE-style non-replay verification computations are counted whether they accept
  or reject; replay-table verification hits are transcript-consistency checks
  and are not charged to $N$ or to the verification-tag terms.

  \emph{Hash games.} There is no adversary construction oracle. For the $\HTW$
  collision game, $N$ counts the $s$ deterministic encryptions on the submitted
  inputs. For the standard preimage game it counts both challenger encryptions:
  the source encryption computing $T^* = \HTW(X^*)$ and the deterministic
  encryption on the submitted input. For the everywhere-preimage game it counts
  the deterministic encryption on the submitted input. For the Sec and eSec
  games of \Cref{cor:second-preimage-rs04}, $N$ and $\Nprim$ are those of the
  explicit induced two-way collision adversary in that corollary, whose source
  and candidate deterministic encryptions are counted exactly as in the
  $s = 2$ collision game.

  \emph{Committing games.} For CMTDs-$\ell$, $N$ uses conservative
  committing-game accounting: after the syntax, domain, and message-length
  checks, it charges the deterministic decryption checks on the common
  ciphertext whether or not the particular $\WiC_\ell$ distinctness check would
  reject before those checks. Likewise, for the CMTs-$\ell$ source games of
  \Cref{cor:committing-combined}, after syntax and domain checks $N$ charges
  the deterministic encryptions that would be used to compare the $s$
  ciphertexts whether or not the $\WiC_\ell$ distinctness check would reject
  earlier.
\item[$\Nprim$:] the number of direct adversary queries to the
  primitive interface, the pair $(\pi, \pi^{-1})$ in
  public permutation games.
\item[$\Ntot = N + \Nprim$:] the total primitive-query cost
  used by the flat sponge lift. The lift of
  \Cref{app:flat-sponge-lift} is applied at $\Ntot$ because after the
  operational duplex-trace accounting of
  \Cref{lem:operational-duplex-trace}, the construction-internal duplex
  calls counted by $N$ and the adversary's direct primitive-interface
  queries counted by $\Nprim$ form a primitive-query sequence of cost at
  most $N+\Nprim$.
\item[$\qv(\mathcal{A})$, $\Qv$:] the number of verification queries
  with valid inputs that are not replay-table hits, and a cap on it.
  Used in the verification terms of the direct mu-ind, AE, and
  \textsf{INT-RUP} bounds of \Cref{sec:aead-security}. Every counted
  verification query contributes at least
  one verification/decryption
  construction call in the $N$ accounting, so $\qv(\mathcal{A}) \leq N$.
\item[$\qRO(\mathcal{B})$, $\QRO$:] the number of direct
  finite-output queries to $\ROflat$ by a random oracle model adversary
  $\mathcal{B}$, and a cap on it. Repeated direct queries on the same
  $Y$ are counted again. Construction-internal $\RF(\cdot)$ lookups are
  not counted. After the derived-adversary construction, $\qRO$ counts only
  simulator-induced $\ROflat$ calls from the adversary's direct
  primitive-interface queries, not the construction calls already
  counted in $N$ for the flat sponge lift.
\item[$\nleaf(\mathcal{B})$, $\Nleaf$:] the number of distinct
  construction-generated final leaf transcripts in the execution of
  $\mathcal{B}$, and a cap on it. Distinct final leaf transcripts are
  generated by valid encryption computations, by valid plaintext-computing
  computations, by valid AE-style
  verification computations (whether they accept or reject), and by
  hash and committing challenger construction computations that pass the
  early validity checks, including the preimage game's source encryption.
\end{description}

For a byte string $X$,
\[
  \leaves(X) = \max\bigl(0,\,
    \lceil \|X\| / \beta \rceil - 1\bigr)
\]
counts the leaf chunks induced by a body byte string of length $\|X\|$.
Thus $\nleaf$ is bounded by the sum of $\leaves(X)$ over the valid
construction computations that create final leaf transcripts, where
$X$ is the plaintext body for encryption computations and the submitted
ciphertext body for plaintext-computing, verification, and committing
decryption checks.
Each distinct final leaf transcript contains a final leaf duplex call
counted in $N$, so $\nleaf(\mathcal{B}) \leq N$ unconditionally. The
ratio $\nleaf/N$ is largest when most processed bytes lie in leaf
chunks, where each full leaf contributes $1 + \beta/\rho = 50$ duplex
calls to $N$ per final leaf transcript.

All theorems are stated for adversaries whose execution-uniform caps
hold at the stated bounds.

\subsection{Loss Terms}
\label{sec:loss-terms}

Five closed-form expressions recur across the security bounds: $\Lift_c$
throughout, with $\KeyGuess_k$, $\RootColl_\tau$, $\RootPre_\tau$, and
$\LeafColl_\tau$ in some. We define them here, after the resource caps that
parameterize them.

\paragraph{Flat sponge loss.}
The flat sponge differentiating loss is
\[
  \Lift_c(\Ntot) = \Ntot(\Ntot + 1) / 2^{c + 1}.
\]
For $\Ntot < 2^c$ it upper-bounds the \cite{BDPV08} sponge-differentiating
advantage; \Cref{lem:lift-bound} carries out the derivation from the
\cite[Lemmas~4 and~5]{BDPV08} product bounds through the Weierstrass
product inequality. For $\Ntot \geq 2^c$, $\Lift_c(\Ntot) \geq 1$ and the
advantage bound is trivial without any appeal to indifferentiability, and
already at $\Ntot \geq 2^{c/2}$ the bound exceeds $1/2$ and provides no
meaningful security.

\paragraph{Key-guessing term.}
The ideal system key-guessing term of \cite{BDPV11} is
\[
  \KeyGuess_k(Q) = Q / 2^k.
\]
It accounts for direct $\ROflat$ queries whose inputs satisfy the
decoder $\KeyedTWOut$ of \Cref{tab:keyed-twout}. Each such query is an
adaptive test of the decoded candidate key against the hidden key $K$,
independent of the query's requested output length. \Cref{lem:key-tests}
bounds the probability that any of $Q$ direct queries tests the right
key by $Q / 2^k$.

\paragraph{Root-collision term.}
The $s$-way root tag multi-collision term is
\[
  \RootColl_\tau(Q, s) = \binom{Q + s}{s} \big/ 2^{\tau (s - 1)},
\]
with the abbreviation $\RootColl_{\tauroot}(Q, s) =
\RootColl_\tau(Q, s)$ at $\tau = \tauroot$. The $+s$ is candidate-sequence
slack for the $s$ final root transcript inputs generated by the
challenger after the adversary outputs its tuples. These lookups are
construction-internal and are not added to $\qRO(\mathcal{B})$; after the
flat sponge lift of \Cref{app:flat-sponge-lift}, the public-permutation
bounds instantiate $Q$ as $\Nprim$.

\paragraph{Root-preimage term.}
The root tag preimage term is
\[
  \RootPre_\tau(Q) = (Q + 1) \big/ 2^{\tau},
\]
with the abbreviation $\RootPre_{\tauroot}(Q) = \RootPre_\tau(Q)$ at
$\tau = \tauroot$. It is the preimage analog of $\RootColl_\tau$: a
single query-independent target tag replaces the $s$-way internal
collision, one challenger lookup replaces the $s$ final-root lookups,
and the exponent drops from $\tau(s-1)$ to $\tau$. As in the
root-collision term, the flat sponge lift gives $Q = \Nprim$.

\paragraph{Leaf-collision term.}
The known-key leaf tag collision term is
\[
  \LeafColl_\tau(Q) = \binom{Q}{2} \big/ 2^\tau.
\]
It is the pairwise birthday bound on $\tau$-bit leaf tag projections
used by the collision-resistance theorem (\Cref{thm:coll}). In this game
the keys are adversary-supplied, so direct primitive queries can pre-read
candidate leaf tag points; after the flat sponge lift this gives
$Q=\Nprim+\Nleaf$. The hidden-key AE-style bounds instead carry the
linear same-slot term $\nleaf/2^{\tauleaf}$ of
\Cref{lem:leaf-collision}, written inline.

\paragraph{Tightness notes.}
The loss terms above are mostly generic search or birthday terms, and their
leading scales are tight in the usual random oracle sense. The key-guessing
term is matched at constants by direct random guesses for a $k$-bit key. In
the multi-user theorem's free-parameter setting, the term
$D\,\KeyGuess_k(\Nprim)$ is also tight at leading scale: using the per-user
nonce discipline of \Cref{sec:mu-ind}, which permits one nonce value across
users, an adversary obtains one known nonempty encryption from each targeted
user and compares each direct candidate-key query against all collected
root-initialization keystream blocks at once. Each guess then matches some
targeted user with probability $D/2^k$. The concrete joint cap charges the
target encryptions to $N$, lowering the feasible scale as discussed in
\Cref{sec:bound-discussion}.

The root terms are the standard random oracle search terms: an $s$-way
birthday search matches the leading scale of
$\RootColl_\tau(Q,s)$, and fixed-target preimage search matches the per-query
$2^{-\tau}$ rate of $\RootPre_\tau(Q)$. The hidden-key same-slot leaf term
$\nleaf/2^{\tauleaf}$ is also tight at constants for nonce-respecting
verification: after encrypting a two-chunk message under nonce $U$, an
adversary can submit non-replay verification queries that replace the leaf
ciphertext chunk by fresh strings of the same length while keeping the root
tag. Each query recomputes the hidden leaf tag at the same slot $(K,U,1)$,
and acceptance occurs when the fresh leaf tag equals the original hidden leaf
tag, with probability $2^{-\tauleaf}$ per verification-side final leaf
transcript.

The flat sponge lift term
$\Lift_c(\Ntot)=\Ntot(\Ntot+1)/2^{c+1}$ dominates the BDPV differentiating
bound $f_P(\Ntot)$ of \cite[Lemma~5]{BDPV08} via \Cref{lem:lift-bound}. That
differentiating bound is matched at leading order by the standard
inner-collision distinguisher; the remaining conservatism is in composing
indifferentiability into the keyed games, not in an obviously improvable
simulator bound. Real-vs-ideal games pay the same differentiating bound once
more, at $\Nprim$, to replace the ideal-side simulator by the real permutation
(\Cref{lem:sim-perm-proximity}); since $\Nprim \leq \Ntot$, that term never
dominates the $\Ntot$ term.

\section{The Authentication Tag as a Hash Output}
\label{sec:root-tag-hash-security}

Externally, a \TW{} encryption returns a ciphertext body and one authentication
tag. Internally, that tag is the final root-duplex output: leaf tags summarize
leaf chunks, the root aggregates those summaries with the context and root
ciphertext chunk, and the wrapper $\HTW$ returns the tag alone. This section
proves that the visible tag is a hash output for the full opening
$(K, U, A, M)$, not merely an AEAD verification value.

The hash notions are the standard ones for $\HTW$: collision,
preimage, second-preimage, and everywhere preimage resistance. The committing
statements are consequences of that hash security: a successful compact opening
must force distinct openings to agree on the same $\tauroot$-bit tag, the
root-collision event bounded below.

The proof uses three ingredients. Transcript parsing separates final root
inputs, except for the explicitly charged leaf-tag collision event.
\Cref{lem:root-candidate-seq} orders every relevant final-root input into a
chronological candidate sequence, and \Cref{lem:root-tag-hash} applies the
adaptive collision and preimage bound of \Cref{app:lemmas} to that sequence.
The random oracle theorems are short instantiations; the flat sponge lift gives
the public-permutation bounds; the \cite{BH22} committing notions reuse the
same root-collision bound.

\Cref{tab:proof-dependency-map} maps the final claims to the proof ingredients
and loss terms they use. It is only a dependency guide; the formal games and
resource accounting remain those of \Cref{sec:security-model}.

\begin{table}[t]
\centering
\footnotesize
\setlength{\tabcolsep}{3pt}
\begin{tabular}{@{}p{0.22\linewidth}p{0.70\linewidth}@{}}
\toprule
Claim & Dependency summary \\
\midrule
Coll, \Cref{thm:coll}
  & Terms: $\Lift_c$, $\LeafColl$, $\RootColl$. Ingredients: final-root input
    separation under no leaf collision; adaptive root candidate collision
    bound. \\
Pre, \Cref{thm:pre}
  & Terms: $\Lift_c$, $\LeafColl$, $\RootPre$, source-guess term.
    Ingredients: source/output separation under no leaf collision;
    sampled-source candidate hit bound. \\
ePre, \Cref{thm:epre}
  & Terms: $\Lift_c$, $\RootPre$. Ingredient: fixed-target root candidate hit
    bound. \\
CMTDs-4, \Cref{cor:cmtds4}
  & Terms: $\Lift_c$, $\RootColl$. Ingredients: \TW{} tidiness; common
    ciphertext body; opening-triple separation. \\
CMT/CMTD, \Cref{cor:committing-combined}
  & Terms: $\Lift_c$, $\RootColl$. Ingredients: \cite{BH22} hierarchy for
    D-notions; direct encryption-source root collision argument for CMTs. \\
mu-ind and AE, \Cref{thm:mu-ind,cor:ae}
  & Terms: $\Lift_c$, $\Lift_c(\Nprim)$, $\KeyGuess$, leaf tag, root tag.
    Ingredients: adjacent-user hybrid; hidden-key leaf collision and
    verification-first root guessing. \\
\textsf{INT-RUP}, \Cref{thm:int-rup}
  & Terms: $\Lift_c$, $\KeyGuess$, leaf tag, root tag. Ingredients:
    plaintext-computing work charged to transcript resources; hidden-key leaf
    collision and \textsf{INT-RUP} root guessing. \\
\bottomrule
\end{tabular}
\caption{Proof dependency map for the final \TW{} bounds. ``Leaf tag'' denotes
the linear hidden-key leaf term in the AE-facing bounds and $\LeafColl$ where
shown explicitly; ``root tag'' denotes the verification or root-preimage
guessing term for the relevant game.}
\label{tab:proof-dependency-map}
\end{table}

\subsection{Tag Hash and Committing Games}
\label{sec:root-tag-games}

The public-permutation games in this section use the primitive, validity, and
construction-oracle conventions of \Cref{sec:games-pp}. The random oracle
versions are obtained by the mechanical substitution of \Cref{sec:ro-model}.

\paragraph{Hash games.}
The hash games for the wrapper $\HTW$ (\Cref{def:htw}) have no
challenger-sampled key. In the $s$-way collision game, all inputs are
adversary-supplied: the adversary outputs $s$ inputs, and the challenger
returns $1$ iff they are pairwise distinct, in the \TW{} domain, and share
a root tag:
\begin{align*}
  &\pi \sample \Perm(\bits^{1600}), \quad
   (K_i, U_i, A_i, M_i)_{i=1}^s \sample \mathcal{A}^{\pi,\, \pi^{-1}}, \\
  &\text{return $0$ unless the inputs are pairwise distinct and in the \TW{} domain}, \\
  &\text{return } \bigl[\HTW(K_1, U_1, A_1, M_1) = \cdots =
     \HTW(K_s, U_s, A_s, M_s)\bigr],
\end{align*}
with advantage $\AdvColl_\TW(\mathcal{A}) = \Pr[\text{the game returns }
1]$; no nonce uniqueness is imposed. The \cite{RS04} standard and
everywhere second-preimage notions are used on fixed input slices; in
\Cref{cor:second-preimage-rs04} they are reduced explicitly to this
two-way collision game.
For fixed byte lengths $a,\mu$, the standard preimage game samples
$X^* \sample \mathcal{X}_{a,\mu}$ independently of the primitive, computes
$T^* = \HTW(X^*)$, gives $T^*$ to
$\mathcal{A}^{\pi,\, \pi^{-1}}$, and returns $1$ iff $\mathcal{A}$ outputs a
valid input $X$ in the \TW{} domain with $\HTW(X) = T^*$; the advantage is
$\AdvPreSlice{a}{\mu}_\TW(\mathcal{A})$. The
everywhere preimage game fixes a target tag
$T^* \in \bits^{\tauroot}$ before the primitive sampling and returns $1$
iff $\mathcal{A}(T^*)$ outputs an input $X$ in the domain with
$\HTW(X) = T^*$, the advantage $\AdvePre_\TW$ taken as the worst case over
query-independent targets.

\paragraph{Committing games.}
The \cite{BH22} committing games are instantiated as in
\Cref{sec:commit-defs}, with no challenger-sampled key. The $s$-way
CMTDs-$\ell$ (decryption) game has the adversary output a ciphertext
$C \concat T$ and $s$ tuples $(K_i, U_i, A_i, M_i)$ with pairwise distinct
$\WiC_\ell$-images, and returns $1$ iff
$\Dec_{K_i}[\pi](U_i, A_i, C \concat T) = M_i$ for every $i$, with the eager
length rejection justified by \Cref{alg:dec}
($\Dec_K(U, A, C \concat T)$ returns $\bot$ or a message of length $\|C\|$);
its advantage is
$\Advcmtds{\ell}_\TW(\mathcal{A})$. The CMTs-$\ell$ (encryption) game has
the adversary output the $s$ tuples alone and returns $1$ iff the
encryptions $\Enc_{K_i}[\pi](U_i, A_i, M_i)$ are all equal; its advantage is
$\Advcmts{\ell}_\TW(\mathcal{A})$.

\subsection{Separating Final Tag Inputs}

The tag value is read from a final root transcript input. The next two facts
rule out structural collisions at that input, except for the explicitly charged
possibility that two leaf tags collide before the root aggregates them.

\begin{proposition}[Opening Triple Separation]
\label{cor:triple-sep}
If two \TW{} root computations have distinct $(K, U, A)$ triples, then
their bridged final root transcript inputs under $\flatmap(\cdot)$ are
distinct, regardless of whether each computation is an encryption or
a decryption check, and whether or not they act on a common ciphertext.
\end{proposition}

\begin{proof}
Consider two root computations $i$ and $j$ with contexts
$(K_i, U_i, A_i) \neq (K_j, U_j, A_j)$, reaching final root transcripts
$R_i$ and $R_j$. Suppose for contradiction that
$\flatmap(R_i) = \flatmap(R_j)$. Since the
bridge $\flatmap(\cdot)$ is injective on well-formed transcript prefixes
(\Cref{lem:bridge-decode}), the native flattened inputs of the two
final-root prefixes are equal. Then \Cref{lem:transcript-inj}
recovers the full duplex-call sequence and reads off the pair
$(\sigma,\sigph)$ for each call. The seven $4$-bit phase suffixes are
pairwise distinct, so the \siginit, $\sigma_{\mathsf{ad}}^\pm$,
$\sigma_{\mathsf{msg}}^\pm$, and (when present) $\sigma_{\mathsf{agg}}^\pm$
calls in the recovered sequence are unambiguously identified. The
aggregation phase is present exactly when $n \geq 1$, and the two
recovered sequences are equal, so they either both carry an aggregation
phase or both omit it; either way the recovery of $(K, U, A)$ below uses only the
\siginit and AD calls, which precede any message or aggregation block.

Equality of the recovered sequences forces equality of the \siginit
call, whose $\sigma$ is $\prfxroot \concat K \concat U$ with
fixed-width $K$ and $U$, hence $K_i = K_j$ and $U_i = U_j$. It also
forces equality of the recovered AD subsequences. By the AD bullet of
\Cref{lem:transcript-inj} the associated data phase contributes a
$\sigadlast$-tagged block exactly when the associated data is nonempty,
so the recovered sequences agree on the presence of the AD phase. If
exactly one of $A_i, A_j$ were empty the recovered sequences would
differ in the presence of a $\sigadlast$ call, contradicting their
equality; hence either both are empty, giving $A_i = A_j = \varepsilon$,
or both are nonempty, in which case the phase-recovery sub-claim of
\Cref{lem:transcript-inj} inverts the $\rho$-byte-block partition map on
the AD byte string and equal AD $\sigma$-block sequences imply
$A_i = A_j$. In all cases $(K_i, U_i, A_i) = (K_j, U_j, A_j)$,
contradicting the hypothesis of distinct triples.
\end{proof}

Triple Separation handles distinct contexts; when inputs share a context
but differ in their messages, distinctness of the final root transcripts
needs the absence of a leaf tag collision, recorded next.

\begin{proposition}[Final-Root Input Separation]
\label{prop:final-root-sep}
Let $\badleaf^\star$ be the event that two distinct construction-generated
final leaf transcripts receive the same $\tauleaf$-bit $\RF$ projection.
On $\neg\badleaf^\star$, if two \TW{} construction computations reach
final root transcripts with $\flatmap(R) = \flatmap(R')$, then they share
the same context $(K, U, A)$ and the same full ciphertext $C$.
Equivalently, on $\neg\badleaf^\star$ computations with distinct
$(K, U, A, C)$ reach distinct bridged final root transcripts.
\end{proposition}

\begin{proof}
By bridge injectivity (\Cref{lem:bridge-decode}),
$\flatmap(R) = \flatmap(R')$ gives equal native flattened final-root
inputs. \Cref{lem:transcript-inj} then recovers the same root
initialization $\prfxroot \concat K \concat U$---hence the same $K$ and
$U$---the same associated data blocks---hence the same $A$---the same root
ciphertext chunk, and the same leaf count and ordered leaf tag sequence.
Under $\neg\badleaf^\star$, Leaf-Transcript Recovery
(\Cref{lem:leaf-recovery}) identifies the equal leaf tag sequences with
the same final leaf transcripts, hence the same leaf ciphertext bodies;
with the common root ciphertext chunk this gives the same full $C$.
\end{proof}

\subsection{Candidate Bound for Tag Inputs}

After separation, the probabilistic question is local: how likely is a collision
or target hit among the final root inputs whose tag values matter to a game?
The candidate machinery rests on one engine, stated here and proved in full in
\Cref{app:lem-candidate-collisions}: adaptively chosen distinct random oracle
inputs whose values are still unread when chosen behave, for collision and
target-hitting purposes, as fresh uniform draws.

\begin{lemma}[Adaptive Candidate Collision and Preimage]
\label{lem:adaptive-candidate}
Consider any interaction with a lazy random oracle $\ROflat : \bits^* \to
\bits^\infty$ that builds a sequence of distinct \emph{candidate} inputs
$Y_1,\dots,Y_m$. The process may make arbitrary non-candidate oracle
queries and may choose each candidate adaptively from the full prior
transcript and private state sampled before the candidate's unread oracle
value. Whenever a new candidate $Y_j$ is appended, two premises hold:
\emph{predictability}---both the decision to append and the string $Y_j$
are determined by the state preceding the append step---and
\emph{unrevealed value}---no oracle lookup before the append step has
requested a positive number of bits of $\ROflat(Y_j)$; zero-length lookups
at $Y_j$ are permitted. Fix a projection length $\tau$, and suppose the
sequence has at most $L$ candidates ($m \leq L$).
\begin{enumerate}
\item \emph{Collision.} For every $s \geq 2$, the
  probability $\Pr[\mathsf{coll}_{s,\tau}]$ that some $s$ distinct
  candidates have equal $\tau$-bit projections,
  \[
    \bigl\lfloor \ROflat(Y_{i_1}) \bigr\rfloor_\tau = \cdots =
    \bigl\lfloor \ROflat(Y_{i_s}) \bigr\rfloor_\tau
    \quad \text{for some } 1 \leq i_1 < \cdots < i_s \leq m,
  \]
  satisfies
  $\Pr[\mathsf{coll}_{s,\tau}] \leq \binom{L}{s} \cdot 2^{-\tau(s-1)}$.
\item \emph{Preimage.} For every target $t \in \bits^\tau$ chosen
  independently of $\ROflat$, the probability
  $\Pr[\mathsf{pre}_{\tau,t}]$ that some candidate's $\tau$-bit projection
  equals $t$, i.e.\ $\bigl\lfloor \ROflat(Y_j) \bigr\rfloor_\tau = t$ for
  some $1 \leq j \leq m$, satisfies
  $\Pr[\mathsf{pre}_{\tau,t}] \leq L \cdot 2^{-\tau}$.
\end{enumerate}
The preimage case has a \emph{second-preimage} corollary: for a
designated candidate position $j_0$, the probability
$\Pr[\mathsf{sec}_{\tau,j_0}]$ that some candidate $j \neq j_0$ has
$\bigl\lfloor \ROflat(Y_j) \bigr\rfloor_\tau =
 \bigl\lfloor \ROflat(Y_{j_0}) \bigr\rfloor_\tau$, i.e.\ collides with the
designated candidate on its $\tau$-bit projection, satisfies
$\Pr[\mathsf{sec}_{\tau,j_0}] \leq (L-1) \cdot 2^{-\tau}$.
\end{lemma}

\begin{proof}[Proof sketch]
The premises license lazy sampling at the append step: predictability
fixes the candidate string before its value is drawn, and the
unrevealed-value premise states that no bit of that value has been
exposed, so $\ROflat(Y_j)$ is uniform and independent of the entire prior
transcript. Conditioning position by position, each fixed position set
then contributes at most $2^{-\tau(s-1)}$ (each fixed position at most
$2^{-\tau}$), and a union bound over the $\binom{L}{s}$ position sets
(the $L$ positions) gives the bounds. The full argument, including the
second-preimage corollary, appears in
\Cref{app:lem-candidate-collisions}.
\end{proof}

The remaining work is structural: exhibiting, in each game, a candidate
sequence that satisfies these premises and contains the final root
transcripts the win event constrains.

\begin{lemma}[Root Tag Candidate Sequence]
\label{lem:root-candidate-seq}
Consider one of the hash or committing random oracle games of
\Cref{sec:ro-model}. The adversary $\mathcal{B}$ makes direct
$\ROflat$ queries and has no construction oracle. The only
construction-internal positive-length lookups at root tag inputs are the
root tag samplings of $h \geq 1$ deterministic challenger encryption or
decryption checks at game-fixed points. Apart from the standard preimage
game's designated sampled source check, all check inputs and keys are
adversary supplied. In that preimage game the sampled source input,
including its key, is private challenger state sampled before the source
tag is revealed, and the source computation counts as one of the $h$
checks. Let the $i$-th check reach a final root transcript $R_i$ whose
root tag is read by a single $\ROflat$ lookup at $\flatmap(R_i)$. Call an
$\ROflat$ input a \emph{root tag input} if it decodes under
$\KeyedTWOut$ to a root tag output-point class---$\Rmsg^+$ (the final
root message block, which carries the root tag only on the no-leaf path)
or $\Rtag$ (the aggregated root tag); multi-chunk computations also reach
$\Rmsg^+$ with zero requested output before aggregation, so this predicate
over-approximates the root tag candidate set, which only loosens the
count. Every execution
then admits a chronological candidate sequence in which each root tag
input is appended exactly once, immediately before the first event that is
either a direct $\ROflat$ query at it or a positive-length $\ROflat$
lookup reading it, and inputs never reached by such an event are never
appended. This sequence is such that
\begin{enumerate}
\item the candidates are pairwise distinct, and if the challenger reaches
  its checks, each $\flatmap(R_i)$ is in the sequence no later than the
  sampling of the $i$-th root tag answer;
\item it satisfies the predictability and unrevealed-value premises of
  \Cref{lem:adaptive-candidate}; and
\item its length is at most $\qRO(\mathcal{B}) + h$.
\end{enumerate}
\end{lemma}

\begin{proof}
A direct query to a keyed transcript is one of $\mathcal{B}$'s visible
$\ROflat$ queries counted in $\qRO(\mathcal{B})$, not a separate
construction query. The only construction computations in the game are
the $h$ deterministic challenger checks named in the statement.

\emph{Distinctness and coverage.} Each root tag input is appended at its
first triggering event and never again, so the candidates are pairwise
distinct by construction, and the rule is well defined however direct
queries and challenger checks interleave. In particular, in the standard
preimage game the source check's root tag sampling may be the trigger for
$\flatmap(R_1)$; a later direct query at the same address is not a first
trigger and appends nothing. Each $\flatmap(R_i)$ is appended no later
than the sampling of the $i$-th root tag answer, since that sampling is a
positive-length lookup at $\flatmap(R_i)$, giving~(1).

\emph{Predictability.} The decoded root-tag-input predicate depends on
the input string alone through the exact decoder of
\Cref{def:keyed-twout}, which ignores the requested output length; by
Output-Point Separation (\Cref{cor:output-sep}) each input decodes to at
most one output-point class. The append rule also uses the already fixed
next event: a direct query triggers at the address, while a construction
lookup triggers only when it requests positive output length. The string
itself is fixed in the state preceding the trigger: a direct query's
input is chosen by $\mathcal{B}$ before its lookup, and a check's final
root transcript is determined by the check's input and the prior oracle
answers---including any private challenger state sampled before the
candidate's unread value, which \Cref{lem:adaptive-candidate} expressly
allows.

\emph{Unrevealed value.} The trigger is the first direct query or
positive-length lookup at the address, so no earlier direct query at it
exists, and it suffices that no earlier construction lookup has read a
positive number of its bits. An earlier positive-length construction
lookup is either a stream output point, a leaf-tag point $\Ltag(j)$, or a
root-tag sampling. The first two are different output-point classes and
hence, by \Cref{cor:output-sep}, have different addresses from root tag
inputs; the last would itself have been the first trigger at this
address. Zero-requested-output lookups, including the $\Rmsg^+$ lookup a
multi-chunk computation makes before aggregation, are permitted by the
unrevealed-value premise of \Cref{lem:adaptive-candidate}, including
when such a zero-length lookup lands at another check's future no-leaf
final-root address. Together with predictability this gives~(2).

\emph{Length.} Appends are triggered only by direct queries---at most
$\qRO(\mathcal{B})$ of them---or by the $h$ checks' root tag samplings;
zero-length construction lookups never trigger an append. Hence the
sequence has
length at most $\qRO(\mathcal{B}) + h$, giving~(3).
\end{proof}

\begin{lemma}[Root Tag Hash and Target-Hit Bound]
\label{lem:root-tag-hash}
Consider one of the hash or committing \TW{} random oracle games
(\Cref{sec:ro-model}) for an adversary $\mathcal{B}$, satisfying the
hypotheses of \Cref{lem:root-candidate-seq}, with win event
$\mathsf{W}$. By the Root
Tag Candidate Sequence lemma
(\Cref{lem:root-candidate-seq}) the root tag inputs form a chronological
candidate sequence meeting the premises of \Cref{lem:adaptive-candidate} with
$\tau = \tauroot$. Writing
$\qRO = \qRO(\mathcal{B})$:
\begin{enumerate}
\item \emph{Root-collision case.} If $\mathsf{W}$ forces $s$ checks to
  reach pairwise distinct bridged final root transcripts carrying a
  common root tag, then, with $h = s$ checks,
  $\Pr[\mathsf{W}] \leq \RootColl_{\tauroot}(\qRO, s)$;
\item \emph{Fixed-target case.} If $\mathsf{W}$ forces a single check
  ($h = 1$) to carry a root tag equal to a target fixed independently of
  $\ROflat$, then
  $\Pr[\mathsf{W}] \leq \RootPre_{\tauroot}(\qRO)$;
\item \emph{Sampled-source case.} If the game first runs a designated
  source check, reveals its root tag, and $\mathsf{W}$ forces one later
  check to reach a distinct bridged final root transcript carrying that
  designated tag, then, with $h = 2$ checks,
  $\Pr[\mathsf{W}] \leq \RootPre_{\tauroot}(\qRO)$.
\end{enumerate}
\end{lemma}

\begin{proof}
In each case the candidate sequence has the stated length and meets the
premises of \Cref{lem:adaptive-candidate}.
\emph{(1)} The $s$ distinct candidates carrying a common
$\tauroot$-projection are an $s$-way collision; the collision case with
$L = \qRO + s$ gives
$\binom{\qRO + s}{s} \cdot 2^{-\tauroot(s-1)}
= \RootColl_{\tauroot}(\qRO, s)$.
\emph{(2)} The candidate hitting the fixed target is a preimage; the
preimage case with $L = \qRO + 1$ gives
$(\qRO + 1) \cdot 2^{-\tauroot} = \RootPre_{\tauroot}(\qRO)$.
\emph{(3)} Let $j_{\mathsf{src}}$ be the position of the designated
source check. The later check is required to be a distinct bridged final
root transcript, so it is a candidate position
$j \neq j_{\mathsf{src}}$ with the same
$\tauroot$-bit projection. This is an application of the second-preimage
corollary, not the fixed-target case, because the designated source tag is
sampled from $\ROflat$ rather than fixed independently of it.
The candidate sequence has length at most $L = \qRO + 2$, and the
second-preimage corollary of \Cref{lem:adaptive-candidate} gives
$(L-1)2^{-\tauroot} = (\qRO+1)2^{-\tauroot}
= \RootPre_{\tauroot}(\qRO)$.
\end{proof}

\subsection{Random Oracle Tag-Hash Bounds}

Each random oracle bound now matches a hash game to a case of
\Cref{lem:root-tag-hash}. Collision and fixed-slice preimage resistance also
carry a leaf-collision term, because distinct openings may have different
ciphertext bodies before reaching the root.

\begin{theorem}[Random Oracle Collision Resistance of $\HTW$]
\label{thm:coll-ro}
For every $s \geq 2$ and every $s$-way collision
adversary $\mathcal{B}$ in the \TW{} random oracle model,
\[
  \AdvColl_{\RO}(\mathcal{B})
  \;\leq\;
  \LeafColl_{\tauleaf}(\qRO(\mathcal{B})+\nleaf(\mathcal{B}))
  + \RootColl_{\tauroot}(\qRO(\mathcal{B}), s).
\]
\end{theorem}

\begin{proof}
Run the random oracle collision game. Let $\badleaf^\star$ be the event
that two distinct construction-generated final leaf transcripts receive
the same $\tauleaf$-bit $\RF$ projection; by the general count
of \Cref{lem:leaf-collision},
$\Pr[\badleaf^\star] \leq \LeafColl_{\tauleaf}(\qRO(\mathcal{B})+\nleaf(\mathcal{B}))$,
so it remains to bound the win under $\neg\badleaf^\star$. After rejecting
unless the $s$ output inputs are pairwise distinct and in the \TW{}
domain, the challenger performs $s$ deterministic encryptions
$\Enc_{K_i}^{\RO}(U_i, A_i, M_i)$, reaching $R_1,\dots,R_s$, all carrying a
common root tag $T$ on a win. Pairwise distinct tuples $(K_i, U_i, A_i, M_i)$
have pairwise distinct $(K_i, U_i, A_i, C_i)$: two tuples differing in their
context already differ in $(K, U, A)$, while two sharing a context but
differing in the message yield distinct ciphertexts, since for a fixed
context encryption maps the message to the ciphertext bijectively. So on
$\neg\badleaf^\star$ Final-Root Input Separation
(\Cref{prop:final-root-sep}) makes $\flatmap(R_1),\dots,\flatmap(R_s)$
pairwise distinct. The Root Tag Hash and Target-Hit Bound (\Cref{lem:root-tag-hash},
item~1) bounds the win under $\neg\badleaf^\star$ by
$\RootColl_{\tauroot}(\qRO(\mathcal{B}),s)$; adding $\Pr[\badleaf^\star]$
gives the theorem.
\end{proof}

\begin{theorem}[Random Oracle Preimage Resistance of $\HTW$]
\label{thm:pre-ro}
Fix byte lengths $0 \leq a \leq A_{\max}$ and
$0 \leq \mu \leq M_{\max}$, and let
$m_{a,\mu}=k+\nu+8a+8\mu$. For every Pre adversary
$\mathcal{B}$ against $\HTW$ on the fixed input slice
$\mathcal{X}_{a,\mu}$ in the \TW{} random oracle model
(\Cref{sec:ro-model}),
\[
\begin{aligned}
  \AdvPreSlice{a}{\mu}_\RO(\mathcal{B})
  \;\leq\;&
  \LeafColl_{\tauleaf}(\qRO(\mathcal{B})+\nleaf(\mathcal{B})) \\
  &+ \RootPre_{\tauroot}(\qRO(\mathcal{B}))
  + 2^{\tauroot-m_{a,\mu}} .
\end{aligned}
\]
\end{theorem}

\begin{proof}
Run the random oracle Pre game. The challenger samples
$X^* \sample \mathcal{X}_{a,\mu}$, writes
$X^*=(K^*,U^*,A^*,M^*)$, computes $T^*=\HTW(X^*)$, gives $T^*$ to
$\mathcal{B}$, and later checks the output $X=(K,U,A,M)$. Let
$\mathsf{same}$ be the event $X=X^*$. For a
fixed random oracle table and fixed adversary coins, the adversary's output
is a function of the $\tauroot$-bit challenge tag $T^*$, so over the
uniform source $X^*$ it can equal $X^*$ with probability at most
$2^{\tauroot}/2^{m_{a,\mu}}=2^{\tauroot-m_{a,\mu}}$. Hence
$\Pr[\mathsf{same}] \leq 2^{\tauroot-m_{a,\mu}}$.

It remains to bound winning executions with $X \neq X^*$. Let
$\badleaf^\star$ be the event that two distinct construction-generated
final leaf transcripts in the source and output checks receive the same
$\tauleaf$-bit $\RF$ projection. By the general count of
\Cref{lem:leaf-collision}---which charges every direct query and every
construction-generated final leaf transcript regardless of key
provenance, covering the challenger-sampled source key---
\[
  \Pr[\badleaf^\star]
  \leq
  \LeafColl_{\tauleaf}(\qRO(\mathcal{B})+\nleaf(\mathcal{B})).
\]
On $\neg\badleaf^\star$, the distinct tuples $X^*$ and $X$ have distinct
final root transcripts: distinct contexts are separated by
\Cref{cor:triple-sep}, while equal contexts and distinct messages produce
distinct ciphertexts, and then \Cref{prop:final-root-sep} applies. A win
with $X \neq X^*$ therefore forces the later output check to hit the
designated source tag $T^*$ at a distinct bridged final root transcript.
The sampled-source case, item~3, of \Cref{lem:root-tag-hash} bounds this
probability by $\RootPre_{\tauroot}(\qRO(\mathcal{B}))$. Adding the
source-guess and leaf-collision events gives the theorem.
\end{proof}

\begin{theorem}[Random Oracle Everywhere Preimage Resistance of $\HTW$]
\label{thm:epre-ro}
For every query-independent target tag $T^* \in \bits^{\tauroot}$ and
every ePre adversary $\mathcal{B}$ in the \TW{}
random oracle model (\Cref{sec:ro-model}),
\[
  \AdvePre_\RO(\mathcal{B})
  \;\leq\; \RootPre_{\tauroot}(\qRO(\mathcal{B}))
  = \frac{\qRO(\mathcal{B}) + 1}{2^{\tauroot}}.
\]
\end{theorem}

\begin{proof}
Run the random oracle preimage game for the query-independent target
$T^*$, which is independent of $\ROflat$. After rejecting unless the
output input $X = (K, U, A, M)$ is in the \TW{} domain, the challenger
performs one encryption $\Enc_{K}^\RO(U, A, M)$, reaching a final root
transcript $R$. A win means $\HTW(X) = T^*$, i.e.\ the root tag
$\lfloor \ROflat(\flatmap(R)) \rfloor_{\tauroot}$ at $R$ equals the fixed
target $T^*$. The fixed-target case, item~2, of
\Cref{lem:root-tag-hash} gives
$\AdvePre_\RO(\mathcal{B}) \leq \RootPre_{\tauroot}(\qRO(\mathcal{B}))$.
\end{proof}

\subsection{Public Permutation Tag-Hash Bounds}
\label{sec:root-tag-pp}

These are the public-permutation hash claims for the authentication tag. The
flat sponge lift transfers the preceding random oracle bounds to the public
permutation model. For these one-world hash games, the lift contributes
$\Lift_c(\Ntot)$, and the derived random oracle adversary has
$\qRO(\mathcal{B}) \leq \Nprim$; challenger construction
computations---including the preimage game's source encryption---are
construction-internal work counted in $N$.

\begin{theorem}[Collision Resistance of $\HTW$]
\label{thm:coll}
For every $s \geq 2$ and every adversary $\mathcal{A}$ against the $s$-way
multi-collision resistance of $\HTW$ in the sense of \cite{BH22}, with the
resources of \Cref{sec:resources},
\[
  \AdvColl_\TW(\mathcal{A})
  \;\leq\;
  \Lift_c(\Ntot)
  + \LeafColl_{\tauleaf}(\Nprim+\Nleaf)
  + \RootColl_{\tauroot}(\Nprim, s).
\]
\end{theorem}

\begin{proof}
By the Lift Application corollary (\Cref{cor:lift-application}) for the
collision game, the derived random oracle adversary $\mathcal{B}$ has
$\qRO(\mathcal{B}) \leq \Nprim$ and $\nleaf(\mathcal{B}) \leq \Nleaf$. The
random oracle bound of \Cref{thm:coll-ro} is nondecreasing in both
$\qRO(\mathcal{B})$ and $\nleaf(\mathcal{B})$, so the monotone form of
\Cref{cor:lift-application} gives the stated bound. The $s$ colliding
inputs may differ in their messages, hence in their ciphertext bodies; the
$\LeafColl$ term accounts for the only way two distinct inputs sharing a
context $(K, U, A)$ can reach the same final root transcript---a collision
among the leaf tags it absorbs---while distinct contexts already reach
distinct final root transcripts by Opening Triple Separation
(\Cref{cor:triple-sep}).
\end{proof}

\begin{corollary}[Two-Way Collision Resistance of $\HTW$]
\label{cor:coll-rs04}
Specializing \Cref{thm:coll} to $s = 2$ gives the classical two-way
collision resistance of \cite{RS04}: for every collision adversary
$\mathcal{A}$ against $\HTW$ with the resources of \Cref{sec:resources},
\[
  \AdvColl_\TW(\mathcal{A})
  \;\leq\;
  \Lift_c(\Ntot)
  + \LeafColl_{\tauleaf}(\Nprim+\Nleaf)
  + \RootColl_{\tauroot}(\Nprim, 2),
\]
with $\RootColl_{\tauroot}(\Nprim, 2) = \binom{\Nprim + 2}{2} / 2^{\tauroot}$.
\end{corollary}

\begin{proof}
The $s = 2$ case of \Cref{thm:coll}; \cite{BH22} note that $s = 2$
recovers the classical notion of collision resistance, which under the
keyed-by-the-permutation convention of \Cref{sec:commit-defs} is the
Coll notion of \cite{RS04}.
\end{proof}

\begin{corollary}[Second-Preimage Resistance of $\HTW$]
\label{cor:second-preimage-rs04}
Fix byte lengths $0 \leq a \leq A_{\max}$ and
$0 \leq \mu \leq M_{\max}$. For every adversary against $\HTW$, with the
resources of \Cref{sec:resources}, in either the standard
second-preimage (Sec) or everywhere second-preimage (eSec) sense of
\cite{RS04} on the fixed input slice $\mathcal{X}_{a,\mu}$, the
corresponding advantage is bounded, independently of $(a, \mu)$, by the
two-way collision bound of
\Cref{cor:coll-rs04}:
\[
  \Lift_c(\Ntot)
  + \LeafColl_{\tauleaf}(\Nprim+\Nleaf)
  + \RootColl_{\tauroot}(\Nprim, 2).
\]
\end{corollary}

\begin{proof}
We use \cite{RS04} to identify the standard and everywhere second-preimage
notions with the fixed-slice notions
$\mathrm{Sec}[m_{a,\mu}]$ and $\mathrm{eSec}[m_{a,\mu}]$ under the
canonical encoding of $\mathcal{X}_{a,\mu}$; the quantitative reduction
is explicit.

For a standard Sec adversary $\mathcal{A}$, define a two-way collision
adversary $\mathcal{C}$ as follows. It samples
$X_0 \sample \mathcal{X}_{a,\mu}$, runs
$\mathcal{A}^{\pi,\pi^{-1}}(X_0)$, receives $X_1$, and outputs
$(X_0,X_1)$. Whenever $\mathcal{A}$ wins, $X_1$ is a valid input,
$X_1 \neq X_0$, and $\HTW(X_1)=\HTW(X_0)$, so $\mathcal{C}$ wins the
two-way collision game with exactly the same probability.

For eSec, fix a source $X_0 \in \mathcal{X}_{a,\mu}$ and define
$\mathcal{C}_{X_0}$ to run the eSec adversary on that source and output
$(X_0,X_1)$. The same pointwise implication holds for every $X_0$; taking
the maximum over sources gives the eSec bound.

The induced collision adversary makes exactly the direct primitive queries
made by $\mathcal{A}$. Its two deterministic challenger checks are the
source and candidate checks, counted exactly as in the $s=2$ collision
game. Applying \Cref{cor:coll-rs04} to the induced collision adversary
gives the displayed bound.
\end{proof}

\begin{theorem}[Preimage Resistance of $\HTW$]
\label{thm:pre}
Fix byte lengths $0 \leq a \leq A_{\max}$ and
$0 \leq \mu \leq M_{\max}$, and let
$m_{a,\mu}=k+\nu+8a+8\mu$. For every adversary $\mathcal{A}$ against
$\HTW$, with the resources of \Cref{sec:resources}, in the standard
preimage (Pre) sense of \cite{RS04} on the fixed input slice
$\mathcal{X}_{a,\mu}$,
\[
\begin{aligned}
  \AdvPreSlice{a}{\mu}_\TW(\mathcal{A})
  \;\leq\;&
  \Lift_c(\Ntot)
  + \LeafColl_{\tauleaf}(\Nprim+\Nleaf) \\
  &+ \RootPre_{\tauroot}(\Nprim)
  + 2^{\tauroot - m_{a,\mu}}.
\end{aligned}
\]
\end{theorem}

\begin{proof}
By the Lift Application corollary (\Cref{cor:lift-application}) for the
preimage game, the derived random oracle adversary $\mathcal{B}$ has
$\qRO(\mathcal{B}) \leq \Nprim$ and
$\nleaf(\mathcal{B}) \leq \Nleaf$; per \Cref{sec:resources}, these caps
cover both challenger encryptions' final leaf transcripts, the source
encryption's included; the encryptions' duplex calls are counted in $N$
and enter through $\Lift_c(\Ntot)$. The random oracle bound of
\Cref{thm:pre-ro} is nondecreasing in both resources, so the monotone form
of \Cref{cor:lift-application} gives the stated bound.
\end{proof}

\begin{theorem}[Everywhere Preimage Resistance of $\HTW$]
\label{thm:epre}
For every target tag $T^* \in \bits^{\tauroot}$ fixed in advance and every
ePre adversary $\mathcal{A}$ against $\HTW$ with the resources of
\Cref{sec:resources},
\[
  \AdvePre_\TW(\mathcal{A})
  \;\leq\;
  \Lift_c(\Ntot) + \RootPre_{\tauroot}(\Nprim)
  \;=\;
  \Lift_c(\Ntot) + \frac{\Nprim + 1}{2^{\tauroot}}.
\]
\end{theorem}

\begin{proof}
By the Lift Application corollary (\Cref{cor:lift-application}) for the
everywhere-preimage game---whose target $T^*$ is fixed before the
primitive
sampling---the derived random oracle adversary $\mathcal{B}$ has
$\qRO(\mathcal{B}) \leq \Nprim$, and the random oracle bound of
\Cref{thm:epre-ro} is nondecreasing in $\qRO(\mathcal{B})$, so the monotone
form of \Cref{cor:lift-application} gives the stated bound. Compared with the
$s$-way root-collision term $\RootColl_{\tauroot}(\cdot, s)$ of
\Cref{sec:loss-terms}, the single fixed target replaces the exponent
$\tauroot(s-1)$ by $\tauroot$.
\end{proof}

\begin{corollary}[Tag Hash Security and Compact Commitment]
\label{cor:tag-commitment}
The \TW{} authentication tag, represented by the wrapper $\HTW$, is a secure
hash of the full encryption context $(K, U, A, M)$ in the sense of
\cite{RS04}: collision-resistant
(\Cref{thm:coll}), preimage-resistant on each fixed input slice
(\Cref{thm:pre}), second-preimage-resistant
(\Cref{cor:second-preimage-rs04}), and everywhere preimage-resistant
(\Cref{thm:epre}). Under the implemented
parameters and work cap $\Ntot \leq 2^{63}$, these advantages are at most
$2^{-128}$ (\Cref{cor:concrete-security}). These hash bounds make the
single $\tauroot$-bit tag a binding compact commitment to $(K, U, A, M)$, with
standard and everywhere preimage resistance, native to the mode. The commitment
is to an opening supplied for the tag: it includes the key and does not make the
tag an independently verifiable public digest of $(U,A,M)$.
\end{corollary}

\begin{proof}
Collision resistance is \Cref{thm:coll}; standard preimage resistance on
fixed input slices is \Cref{thm:pre}; second-preimage resistance follows from
collision resistance by \Cref{cor:second-preimage-rs04}; and everywhere
preimage resistance is \Cref{thm:epre}. These bounds quantify the
adversary's success at the same $\tauroot$-bit projection of the final root
transcript, and \Cref{cor:concrete-security} evaluates them at the
implemented parameters. Binding is the compact-commitment property supplied
by collision and second-preimage resistance. Standard preimage resistance is
the property that a tag produced from a uniformly sampled fixed-length
opening does not reveal an efficiently findable opening; everywhere preimage
resistance is the corresponding fixed-target property for tags chosen
independently of the primitive.
\end{proof}

\subsection{Committing Consequences of Tag Hash Security}

The committing consequences use the same tag-collision analysis under the
opening constraints of the \cite{BH22} games. CMTDs-4 is slightly tighter than
general collision resistance: all openings share one ciphertext body, so
distinct openings separate through their opening triples and no leaf-collision
term is needed.

\begin{theorem}[Random Oracle CMTDs-4]
\label{thm:cmtds4-ro}
For every $s \geq 2$ and every $s$-way CMTDs-4
adversary $\mathcal{B}$ in the \TW{} random oracle model
(\Cref{sec:ro-model}),
\[
  \Advcmtds{4}_\RO(\mathcal{B})
  \;\leq\; \RootColl_{\tauroot}(\qRO(\mathcal{B}), s)
  = \frac{\binom{\qRO(\mathcal{B}) + s}{s}}{2^{\tauroot (s - 1)}}.
\]
\end{theorem}

\begin{proof}
Run the random oracle CMTDs-4 game. After rejecting on any failed syntax,
message-length, or pairwise-full-input check, the challenger performs $s$
deterministic decryption checks $\Dec_{K_i}^\RO(U_i, A_i, C \concat T)$ on
the one common ciphertext $C \concat T$, reaching $R_1,\dots,R_s$. On a winning
execution the tuples $(K_i, U_i, A_i, M_i)$ are pairwise distinct, and since
all checks decrypt the common ciphertext, determinism of $\Dec$ makes the
triples $(K_i, U_i, A_i)$ pairwise distinct---equal triples would return
equal messages, hence equal tuples. By Opening Triple Separation
(\Cref{cor:triple-sep}) the candidates $\flatmap(R_1),\dots,\flatmap(R_s)$
are pairwise distinct (leaf tag collisions cannot merge them), and every
accepting check carries the common root tag $T$. The Root Tag Hash and Target-Hit Bound
(\Cref{lem:root-tag-hash}, item~1) gives
$\Advcmtds{4}_\RO(\mathcal{B}) \leq \RootColl_{\tauroot}(\qRO(\mathcal{B}), s)$.
\end{proof}

\begin{corollary}[CMTDs-4 Committing Security]
\label{cor:cmtds4}
For every $s \geq 2$ and every $s$-way CMTDs-4
adversary $\mathcal{A}$ against \TW{} with the resources of
\Cref{sec:resources},
\[
  \Advcmtds{4}_\TW(\mathcal{A})
  \;\leq\;
  \Lift_c(\Ntot) + \RootColl_{\tauroot}(\Nprim, s)
  \;=\;
  \Lift_c(\Ntot) + \frac{\binom{\Nprim + s}{s}}{2^{\tauroot (s - 1)}}.
\]
\end{corollary}

\begin{proof}
A CMTDs-4 winner gives one ciphertext $C \concat T$ and $s$ distinct inputs
$(K_i, U_i, A_i, M_i)$ such that
$\Dec_{K_i}(U_i, A_i, C \concat T) = M_i$ for every $i$. By \TW{} tidiness
(\Cref{lem:tw-tidiness}), each opened input re-encrypts to the same ciphertext:
$\Enc_{K_i}(U_i, A_i, M_i) = C \concat T$. In particular the opened inputs share
the root tag $T$. For the displayed no-$\LeafColl$ bound, use the dedicated
root-collision derivation of \Cref{thm:cmtds4-ro}, lifted by
\Cref{cor:lift-application}: in this all-bodies-equal case, the single shared
body $C$ forces the triples $(K_i, U_i, A_i)$ pairwise distinct; by Opening
Triple Separation (\Cref{cor:triple-sep}), a leaf tag collision cannot merge
two openings and no $\LeafColl$ term arises.
\end{proof}

\paragraph{Remaining committing notions.}
The hierarchy and $s$-way extension of the \cite{BH22} committing notions
are given in \cite[\S3, Figs.~4--5, and App.~A]{BH22}. We use those
relations for the D-notions and a direct source-game argument for the
CMTs-$\ell$ notions to preserve the resource split of
\Cref{sec:resources}.

\begin{corollary}[Committing-Notion Root-Collision Bound]
\label{cor:committing-combined}
For every $s \geq 2$, every
\[
  X \in \bigl\{\mathsf{CMTDs\text{-}1},\, \mathsf{CMTDs\text{-}3},\,
              \mathsf{CMTDs\text{-}4},\, \mathsf{CMTs\text{-}1},\,
              \mathsf{CMTs\text{-}3},\, \mathsf{CMTs\text{-}4}\bigr\},
\]
and every $s$-way $X$-adversary $\mathcal{A}$ against \TW{} with the
resources of \Cref{sec:resources},
\[
  \Adv^X_\TW(\mathcal{A})
  \;\leq\;
  \Lift_c(\Ntot) + \RootColl_{\tauroot}(\Nprim, s)
  \;=\;
  \Lift_c(\Ntot) + \frac{\binom{\Nprim + s}{s}}{2^{\tauroot(s-1)}}.
\]
\end{corollary}

\begin{proof}
CMTDs-4 is bounded directly by \Cref{cor:cmtds4}. For the remaining
D-notions, the $s$-way hierarchy of \cite[\S3, Fig.~5 and
App.~A]{BH22} gives, for $\ell \in \{1,3\}$,
$\Advcmtds{\ell}_\TW(\mathcal{A}) \leq
\Advcmtds{4}_\TW(\mathcal{A})$ for the same output distribution. By the
conservative committing-game accounting of \Cref{sec:resources}, the
CMTDs-$\ell$ and CMTDs-4 resource caps charge the same deterministic
decryption checks after syntax, domain, and message-length validation,
independently of their distinctness gates.

For CMTs-$\ell$, the challenger performs $s$ deterministic encryption
checks, which supply the candidate final-root transcripts
(\Cref{lem:root-candidate-seq} builds the candidate sequence for
encryption and decryption checks alike). On a winning execution all
$s$ encryptions are equal, hence of equal length and with one common root
tag $T$. Pairwise $\WiC_1$-distinctness gives pairwise distinct keys,
pairwise $\WiC_3$-distinctness gives pairwise distinct $(K, U, A)$ triples,
and pairwise $\WiC_4$-distinctness gives pairwise distinct full inputs; in
the last case, equal triples and equal ciphertexts would decrypt
deterministically to equal messages, contradicting full-input
distinctness. Thus each CMTs-$\ell$ winner reaches pairwise distinct
opening triples. Opening Triple Separation (\Cref{cor:triple-sep}) makes
the $s$ final-root candidates distinct while their $\tauroot$-bit
projections all equal $T$.
The CMTs-$\ell$ encryption checks are the construction-internal work
counted in $N$ by \Cref{sec:resources}.

By the Lift Application corollary (\Cref{cor:lift-application}) for the
CMTs-$\ell$ game, the derived random oracle adversary $\mathcal{B}$ has
$\qRO(\mathcal{B}) \leq \Nprim$. The Root Tag Hash and Target-Hit Bound
(\Cref{lem:root-tag-hash}, item~1) then bounds its win by
$\RootColl_{\tauroot}(\qRO(\mathcal{B}), s)$ in the random oracle model,
which the lift transfers to the public permutation setting at $\Nprim$;
the encryption checks are construction-internal work counted in $N$, not
in $\Nprim$.
\end{proof}

\section{Baseline Authenticated Encryption Bounds}
\label{sec:aead-security}

\Cref{sec:root-tag-hash-security} proves the tag-centric properties of \TW{}.
Here we discharge the AEAD-facing obligations: multi-user indistinguishability,
its single-user all-in-one AE corollary, and \textsf{INT-RUP} integrity. These
bounds use the same hidden-key, leaf-collision, and tag-guessing accounting as
the tag analysis, but their role is narrower: they show that the mode remains a
conventional nonce-respecting AEAD.

We first prove mu-ind in the flat random oracle model, lift the result to the
public permutation model using \Cref{app:flat-sponge-lift}, and derive
all-in-one AE as the $D = 1$ case. We then prove \textsf{INT-RUP} in the
separated-interface model of \cite{ABLMMY14}, where plaintext-computing queries
increase the construction work counts but do not change the non-replay
verification count. All concrete instantiations are summarized in
\Cref{sec:concrete-security}.

\subsection{AEAD Public Permutation Games}
\label{sec:aead-games}

The public-permutation games in this section use the primitive, validity, and
construction-oracle conventions of \Cref{sec:games-pp}. The random oracle
versions are obtained by the mechanical substitution of \Cref{sec:ro-model}.

\paragraph{All-in-one AE.}
The all-in-one authenticated encryption game~\cite{RS06} samples
$K \sample \bits^k$ and compares the real pair
$(\Enc_K[\pi], \Ver_K[\pi])$ with the ideal pair $(\Rand, \Replay)$.
On each accepted encryption query $(U, A, M)$, $\Rand$ samples a uniformly
random byte string of length $\|M\| + \tauroot/8$, parses it as
$C \concat T$ with $T \in \bits^{\tauroot}$, records $(U, A, C, T)$, and
returns $C \concat T$; $\Replay$ accepts exactly recorded tuples. Encryption
queries are nonce-respecting, and
\[
  \Advae_\TW(\mathcal{A}) =
  \bigl|
    \Pr[\mathcal{A}^{\Enc_K[\pi],\, \Ver_K[\pi],\, \pi,\, \pi^{-1}} \Rightarrow 1]
    - \Pr[\mathcal{A}^{\Rand,\, \Replay,\, \pi,\, \pi^{-1}} \Rightarrow 1]
  \bigr|.
\]
This all-in-one AE advantage is the single-user case of the multi-user
real-vs-ideal notion of \Cref{sec:mu-ind}; \Cref{cor:ae} derives it from
\Cref{thm:mu-ind}.

\paragraph{RUP integrity.}
The \textsf{INT-RUP} game samples $K \sample \bits^k$ and
$\pi \sample \Perm(\bits^{1600})$, exposes $\Enc_K[\pi]$,
$\PDec_K[\pi]$, $\Ver_K[\pi]$, and $(\pi, \pi^{-1})$, and maintains the replay
table $W$ of encryption outputs. Accepted encryption queries are
nonce-respecting: an encryption query is rejected if its nonce has appeared in a
previous accepted encryption query. Plaintext-computing and verification
queries may use arbitrary nonces. On an accepted encryption query $(U, A, M)$ the
game returns $\Enc_K[\pi](U, A, M)$, parsed and recorded as $(U, A, C, T)$ in $W$.
On a valid plaintext-computing query $(U, A, C, T)$ it returns
$\PDec_K[\pi](U, A, C, T)$, and on an invalid query it returns $\bot$. On a valid
verification query $(U, A, C, T)$ it returns $\Ver_K[\pi](U, A, C, T)$; if this bit
is $1$ and $(U, A, C, T) \notin W$ when the query is asked, the game records a
forgery. Invalid verification queries return $0$. The game returns $1$ iff a
forgery is recorded, and
\[
  \AdvintRUP_\TW(\mathcal{A}) =
  \Pr[\mathsf{G}^{\mathsf{INT\text{-}RUP}}_\TW(\mathcal{A}) \Rightarrow 1].
\]

\subsection{Multi-User Setting}
\label{sec:mu-setting}

This subsection instantiates the mu-ind game of \cite{BT16}
(\Cref{sec:mu-ind}) for \TW{} in the public permutation model, under
the primitive convention of \Cref{sec:games-pp}. The challenger
samples a challenge bit $b \sample \bits$ and a permutation
$\pi \sample \Perm(\bits^{1600})$, and exposes $(\pi, \pi^{-1})$ as the
primitive interface in both worlds. The
adversary $\mathcal{A}$ creates users on demand via $\mathsf{New}$, and
the game maintains an active-user counter $v$. We write $D$ for an
execution-uniform cap on the number of $\mathsf{New}$ calls. Equivalently,
the game may sample independent keys $K_1,\dots,K_D$ at the start of the
experiment and let the $d$-th $\mathsf{New}$ call activate $K_d$; inactive
indices are never valid users.

The mu-ind construction oracles inherit the validity convention of
\Cref{sec:games-pp}. The encryption oracle on input $(d, U, A, M)$
rejects, without recording, if $(U, A, M)$ is outside the \TW{} encryption
domain, if $d \notin \{1, \dots, v\}$, or if $(d, U) \in V$; otherwise
it sets $C = \Enc_{K_d}[\pi](U, A, M)$ when $b = 1$, and samples a
uniformly random byte string $C$ of length $\|M\| + \tauroot/8$ when
$b = 0$. The game records
$(d, U)$ in $V$ and $(d, U, A, C)$ in $W$, and returns $C$. The
verification oracle on input $(d, U, A, C)$ rejects if $d \notin
\{1, \dots, v\}$ or if $(U, A, C)$ is not a valid \TW{}
verification/decryption input; if $(d, U, A, C) \in W$ it returns
$\mathsf{true}$; if $b = 0$ it returns $\mathsf{false}$; otherwise it
returns $[\Dec_{K_d}[\pi](U, A, C) \neq \bot]$. Nonce uniqueness is
enforced per user by the set $V$.

After querying these oracles, $\mathcal{A}$ outputs $b' \in \bits$.
The mu-ind advantage of $\mathcal{A}$ against \TW{} is
\[
  \Advmuind_\TW(\mathcal{A})
  = |2 \Pr[b' = b] - 1|.
\]
A valid verification query is \emph{non-replay} if its tuple is not in
$W$ when the query is asked; replay hits are answered by the recorded
transcript and are not counted in the verification-query resource below.

Resource caps split per potential user $d \leq D$: $N^{(d)}$,
$\qv^{(d)}$, $\nleaf^{(d)}$ are the per-user counts under
\Cref{sec:resources}, with inactive users contributing zero. The global
totals $N = \sum_{d=1}^D N^{(d)}$, $\qv = \sum_{d=1}^D \qv^{(d)}$,
$\nleaf = \sum_{d=1}^D \nleaf^{(d)}$ appear in the final bound: $N$
through the lift term $\Lift_c(\Ntot)$, and $\nleaf$ and $\qv$ through
the leaf-collision and verification-tag terms. $\Nprim$
remains a single shared count over the primitive interface.

The mu-ind treatment applies only to indistinguishability. The
committing security games already give the adversary control of all
$s$ keys, so the $s$-way CMTDs-$\ell$ statements of
\Cref{sec:root-tag-games} are themselves the multi-key committing statement
and need no separate $\mathsf{mu\text{-}cmt}$ notion.

All random oracle adversaries in this section are flat-RO adversaries
(\Cref{sec:ro-model}) with the displayed resources, and the flat sponge lift
transfers each random oracle bound using the derived-adversary cap
$\qRO(\mathcal{B}) \leq \Nprim$.

\begin{theorem}[Random Oracle Multi-User Indistinguishability]
\label{thm:mu-ind-ro}
For every flat-RO mu-ind adversary $\mathcal{B}$ with
execution-uniform user cap $D$ and resource counts of
\Cref{sec:mu-setting},
\[
  \Adv^{\mathsf{mu\text{-}ind}}_{\RO}(\mathcal{B})
  \;\leq\;
  \frac{D(D - 1)}{2^{k + 1}}
  + D \cdot \KeyGuess_k(\qRO(\mathcal{B}))
  + \frac{\nleaf(\mathcal{B})}{2^{\tauleaf}}
  + \frac{\qv(\mathcal{B})}{2^{\tauroot}}.
\]
Here $\qv$ is the non-replay verification-query count of
\Cref{sec:resources}.
\end{theorem}

\begin{proof}
Use the pre-sampled-key formulation of \Cref{sec:mu-setting}: sample
keys $K_1,\dots,K_D$ at the start of the experiment, and let the $d$-th
$\mathsf{New}$ call activate $K_d$, unused keys being ignored.
Let $\mathsf{coll}$ denote the event that two of the pre-sampled keys
$K_1,\dots,K_D$ collide;
$\Pr[\mathsf{coll}] \leq \binom{D}{2}/2^k = D(D-1)/2^{k+1}$ by union
bound. In the random oracle model, use a standard hybrid over the $D$
users: $H_d$ answers users $1,\dots,d$ with $(\Rand,\Replay)$ and
users $d+1,\dots,D$ with the real construction interface, so $H_0$ is
the real mu-ind game and $H_D$ is the ideal one. Let $H_d^\star$ be
$H_d$ with a fixed output bit on $\mathsf{coll}$. Since
$\mathsf{coll}$ has the same probability in all hybrids,
\[
  \Adv^{\mathsf{mu\text{-}ind}}_{\RO}(\mathcal{B})
  \leq
  \Pr[\mathsf{coll}]
  +
  \bigl|\Pr[\mathcal{B}^{H_0^\star} \Rightarrow 1]
       - \Pr[\mathcal{B}^{H_D^\star} \Rightarrow 1]\bigr|.
\]
By the triangle inequality, it remains to bound each adjacent stopped
hybrid. For each $d$, \Cref{lem:mu-adjacent-hybrid} gives
\[
  \bigl|\Pr[\mathcal{B}^{H_{d-1}^\star} \Rightarrow 1]
       - \Pr[\mathcal{B}^{H_d^\star} \Rightarrow 1]\bigr|
  \leq
  \KeyGuess_k(\qRO(\mathcal{B}))
  + \frac{\nleaf^{(d)}}{2^{\tauleaf}}
  + \frac{\qv^{(d)}}{2^{\tauroot}}.
\]
\Cref{lem:mu-adjacent-hybrid} is the only construction-specific step in
the baseline mu-ind proof. It couples the adjacent user-$d$ hybrids through
the all-reject continuation of \Cref{def:adjacent-hybrid-games}. The
coupling agrees until a direct random-oracle query tests the hidden key
$K_d$ or a non-replay verification comparison succeeds. The resulting
first-occurrence decomposition has exactly the three terms displayed above:
the foreign-key test is bounded by
$\KeyGuess_k(\qRO(\mathcal{B}))$; a same-slot collision between final leaf
tags for user $d$ is bounded by
$\nleaf^{(d)}/2^{\tauleaf}$ via \Cref{lem:leaf-collision}; and a successful
verification-first root comparison is bounded by
$\qv^{(d)}/2^{\tauroot}$ via \Cref{lem:ver-first-root-guess}. Encryption
transcript marginal uniformity and environment neutrality
(\Cref{lem:enc-transcript-marginal-uniformity,lem:env-neutrality}) are the
facts that make this coupling independent of the other users.

Summing across $d = 1,\dots,D$ gives
$D \KeyGuess_k(\qRO(\mathcal{B}))$ for the key-test contribution. The
leaf tag contribution sums exactly:
\[
  \sum_d \frac{\nleaf^{(d)}}{2^{\tauleaf}}
  = \frac{\nleaf(\mathcal{B})}{2^{\tauleaf}},
\]
since $\sum_d \nleaf^{(d)} = \nleaf(\mathcal{B})$. The
verification-tag contribution satisfies
\[
  \sum_d \qv^{(d)} / 2^{\tauroot}
  = \qv(\mathcal{B}) / 2^{\tauroot}.
\]
Combining these estimates with the $\mathsf{coll}$ penalty gives the
stated RO bound.
\end{proof}

\begin{theorem}[Multi-User Indistinguishability]
\label{thm:mu-ind}
For every mu-ind adversary $\mathcal{A}$ against \TW{} with
execution-uniform user cap $D$ and the global resources of
\Cref{sec:mu-setting},
\[
  \Advmuind_\TW(\mathcal{A})
  \;\leq\;
  \Lift_c(\Ntot)
  + \Lift_c(\Nprim)
  + \frac{D(D - 1)}{2^{k + 1}}
  + D \cdot \KeyGuess_k(\Nprim)
  + \frac{\Nleaf}{2^{\tauleaf}}
  + \frac{\Qv}{2^{\tauroot}}.
\]
\end{theorem}

\begin{proof}
By the Lift Application corollary (\Cref{cor:lift-application}) for the
mu-ind game---a real-vs-ideal game, so the ideal-side proximity term
applies---there is a random oracle mu-ind adversary
$\mathcal{B} = \mathcal{B}_{\mathrm{derived}}(\mathcal{A})$ with
$\qRO(\mathcal{B}) \leq \Nprim$, $\nleaf(\mathcal{B}) \leq \Nleaf$,
$\qv(\mathcal{B}) \leq \Qv$, and the per-user construction-side caps of
\Cref{sec:mu-setting} such that
\[
  \Advmuind_\TW(\mathcal{A})
  \;\leq\;
  \Lift_c(\Ntot) + \Adv^{\mathsf{mu\text{-}ind}}_{\RO}(\mathcal{B})
  + \Lift_c(\Nprim).
\]
\Cref{thm:mu-ind-ro} bounds the random oracle term. Substituting
$\qRO(\mathcal{B}) \leq \Nprim$,
$\nleaf(\mathcal{B}) \leq \Nleaf$, and
$\qv(\mathcal{B}) \leq \Qv$ gives the stated bound.
\end{proof}

\begin{corollary}[All-in-One AE Security]
\label{cor:ae}
For every nonce-respecting AE adversary $\mathcal{A}$ against \TW{}
with the resources of \Cref{sec:resources},
\[
  \Advae_\TW(\mathcal{A})
  \;\leq\;
  \Lift_c(\Ntot)
  + \Lift_c(\Nprim)
  + \KeyGuess_k(\Nprim)
  + \frac{\Nleaf}{2^{\tauleaf}}
  + \frac{\Qv}{2^{\tauroot}}.
\]
\end{corollary}

\begin{proof}
Build a mu-ind adversary $\mathcal{A}^{\mu}$ that makes one
$\mathsf{New}$ query, runs $\mathcal{A}$, and forwards each
encryption and verification query of $\mathcal{A}$ to user $1$. The
primitive interface is forwarded unchanged. The real and ideal views of
$\mathcal{A}$ in the all-in-one AE experiment are exactly the real and
ideal views of $\mathcal{A}^{\mu}$ in the mu-ind experiment with
$D = 1$, and the resource counts are unchanged. On a replayed
encryption tuple the two real verification oracles agree: the mu-ind
oracle short-circuits through its replay table, while the AE oracle
re-runs decryption, which accepts by \TW{} decryption correctness
(\Cref{sec:dec}); all other queries are answered by identical
code. Applying
\Cref{thm:mu-ind} with $D = 1$ eliminates the key-collision term and
gives the stated bound.
\end{proof}

The \textsf{INT-RUP} bound is AEAD-facing, but it is closer to the tag-hash
analysis than to the mu-ind baseline: plaintext-computing queries expose
unverified stream plaintext, and freshness must still reduce to hidden key
tests, leaf-tag collisions, and root-tag guessing. The separated interface
changes the accounting by adding plaintext-computing work to $N$ and
$\nleaf$, while leaving the non-replay verification count $\qv$ unchanged.

\begin{theorem}[Random Oracle \textsf{INT-RUP} Integrity]
\label{thm:int-rup-ro}
For every flat-RO \textsf{INT-RUP} adversary $\mathcal{B}$ against \TW{}
with nonce-respecting encryption queries and the resources of
\Cref{sec:resources},
\[
  \AdvintRUP_{\RO}(\mathcal{B})
  \;\leq\;
  \KeyGuess_k(\qRO(\mathcal{B}))
  + \frac{\nleaf(\mathcal{B})}{2^{\tauleaf}}
  + \frac{\qv(\mathcal{B})}{2^{\tauroot}}.
\]
Plaintext-computing queries contribute to $N$ and $\nleaf$, but not to
$\qv$.
\end{theorem}

\begin{proof}
The proof uses the same hidden-tag accounting as the adjacent-hybrid bound, with
the plaintext-computing interface included in the construction transcript.
Work in the flat-RO \textsf{INT-RUP} game of \Cref{sec:ro-model}. Let
$\forge$ be the event that a valid non-replay verification query accepts.
Let $\badkey$ be the event that a direct $\ROflat$ query decodes under
$\KeyedTWOut$ to the hidden key $K$ and a counted output-point class, and let
$\badleaf$ be the event that some construction-generated final leaf
transcript receives the same $\tauleaf$-bit leaf tag projection as a
distinct final leaf transcript that shares its key, nonce, and leaf index
and is reached by an encryption computation (the same-slot event of
\Cref{lem:leaf-collision}). Final leaf
transcripts generated by encryption, plaintext-computing, and verification
computations are all included in $\nleaf$.

The key-test term is unchanged by the plaintext-computing oracle. By
\Cref{lem:ver-key-uniformity}, stated for the \textsf{INT-RUP} extension,
the hidden key remains uniform before $\badkey$ outside the candidate keys
already tested by direct queries, and \Cref{lem:key-tests} gives
\[
  \Pr[\badkey] \leq \KeyGuess_k(\qRO(\mathcal{B})).
\]

For the leaf term, no direct query has pre-read a hidden-key final leaf tag
point before $\badkey$: such a query would decode under $\KeyedTWOut$ to $K$
and the corresponding $\Ltag(j)$ class. Coupling the real
\textsf{INT-RUP} execution to the all-reject continuation used in
\Cref{lem:int-rup-root-guess} preserves the construction computations and
leaf-tag table accesses up to the first accepting non-replay submission.
The same-slot case of \Cref{lem:leaf-collision}, applied to the enlarged set
of construction computations, gives
\[
  \Pr[\badleaf \text{ before } \badkey \cup \forge]
  \leq
  \frac{\nleaf(\mathcal{B})}{2^{\tauleaf}}.
\]

It remains to bound accepting verification before $\badkey \cup \badleaf$.
The \textsf{INT-RUP} root-guessing lemma (\Cref{lem:int-rup-root-guess})
gives
\[
  \Pr[\forge \text{ before } \badkey \cup \badleaf]
  \leq \qv(\mathcal{B}) / 2^{\tauroot}.
\]
Combining the three bounds through the first-occurrence decomposition
of $\Pr[\forge]$ proves the theorem.
\end{proof}

\begin{theorem}[\textsf{INT-RUP} Integrity]
\label{thm:int-rup}
For every \textsf{INT-RUP} adversary $\mathcal{A}$ against \TW{} with
nonce-respecting encryption queries and the resources of \Cref{sec:resources},
\[
  \AdvintRUP_\TW(\mathcal{A})
  \;\leq\;
  \Lift_c(\Ntot)
  + \KeyGuess_k(\Nprim)
  + \frac{\Nleaf}{2^{\tauleaf}}
  + \frac{\Qv}{2^{\tauroot}}.
\]
\end{theorem}

\begin{proof}
By the Lift Application corollary (\Cref{cor:lift-application}) for the
\textsf{INT-RUP} game, there is a random oracle \textsf{INT-RUP} adversary
$\mathcal{B} = \mathcal{B}_{\mathrm{derived}}(\mathcal{A})$ with
$\qRO(\mathcal{B}) \leq \Nprim$, $\nleaf(\mathcal{B}) \leq \Nleaf$, and
$\qv(\mathcal{B}) \leq \Qv$ such that
\[
  \AdvintRUP_\TW(\mathcal{A})
  \;\leq\;
  \Lift_c(\Ntot) + \AdvintRUP_{\RO}(\mathcal{B}).
\]
\Cref{thm:int-rup-ro} bounds the random oracle term, and substituting the
derived-adversary resource caps gives the stated bound.
\end{proof}

\section{Concrete Security Summary}
\label{sec:concrete-security}

\Cref{sec:aead-security,sec:root-tag-hash-security} separate the security
analysis into AEAD bounds and root tag hash and committing bounds. This section
instantiates those bounds at the
implemented \TW{} parameters and records the single concrete work cap used
throughout the paper.

\subsection{Concrete TW128 Bounds}
\label{sec:concrete}

The implemented parameters of \Cref{sec:params} fix
\[
  k = c = \nu = \tauroot = \tauleaf = 256.
\]
Substituting into
\Cref{cor:ae,thm:mu-ind,thm:int-rup,thm:coll,thm:pre,cor:second-preimage-rs04,cor:cmtds4,thm:epre}
gives the following closed-form bounds.

\paragraph{All-in-one AE.}
\[
  \Advae_\TW(\mathcal{A})
  \;\leq\;
  \frac{\Ntot(\Ntot + 1)}{2^{257}}
  + \frac{\Nprim(\Nprim + 1)}{2^{257}}
  + \frac{\Nprim + \Qv + \Nleaf}{2^{256}}.
\]

\paragraph{Multi-user indistinguishability.}
\[
  \Advmuind_\TW(\mathcal{A})
  \;\leq\;
  \frac{\Ntot(\Ntot + 1)}{2^{257}}
  + \frac{\Nprim(\Nprim + 1)}{2^{257}}
  + \frac{D \Nprim + \Qv + \Nleaf}{2^{256}}
  + \frac{D(D - 1)}{2^{257}}.
\]

\paragraph{\textsf{INT-RUP} integrity.}
\[
  \AdvintRUP_\TW(\mathcal{A})
  \;\leq\;
  \frac{\Ntot(\Ntot + 1)}{2^{257}}
  + \frac{\Nprim + \Qv + \Nleaf}{2^{256}}.
\]

\paragraph{$s$-way collision resistance of $\HTW$.}
For every $s \geq 2$,
\[
  \AdvColl_\TW(\mathcal{A})
  \;\leq\;
  \frac{\Ntot(\Ntot + 1)}{2^{257}}
  + \frac{\binom{\Nprim+\Nleaf}{2}}{2^{256}}
  + \frac{\binom{\Nprim + s}{s}}{2^{256 (s - 1)}}.
\]
For $s=2$, this is the ordinary two-input collision bound:
\[
\begin{aligned}
  \AdvColl_\TW(\mathcal{A})
\leq\;&
  \frac{\Ntot(\Ntot + 1)}{2^{257}}
  + \frac{(\Nprim+\Nleaf)(\Nprim+\Nleaf-1)}{2^{257}} \\
& + \frac{(\Nprim + 1)(\Nprim + 2)}{2^{257}} .
\end{aligned}
\]

\paragraph{Second-preimage resistance of $\HTW$.}
By \Cref{cor:second-preimage-rs04}, the standard and everywhere
\cite{RS04} second-preimage advantages of $\HTW$ are bounded by the same
$s=2$ collision expression above.

\paragraph{Preimage resistance of $\HTW$.}
For fixed byte lengths $a,\mu$, let
$m_{a,\mu}=k+\nu+8a+8\mu$. By \Cref{thm:pre},
\[
\begin{aligned}
  \AdvPreSlice{a}{\mu}_\TW(\mathcal{A})
\leq\;&
  \frac{\Ntot(\Ntot + 1)}{2^{257}}
  + \frac{(\Nprim+\Nleaf)(\Nprim+\Nleaf-1)}{2^{257}} \\
& + \frac{\Nprim + 1}{2^{256}}
  + 2^{256 - m_{a,\mu}} .
\end{aligned}
\]
Under the implemented parameters $k = \nu = 256$, the slice term is
$2^{-256-8a-8\mu}$.

\paragraph{$s$-way CMTDs-4 committing security.}
For every $s \geq 2$, the all-bodies-equal special case drops the
known-key leaf tag collision term:
\[
  \Advcmtds{4}_\TW(\mathcal{A})
  \;\leq\;
  \frac{\Ntot(\Ntot + 1)}{2^{257}}
  + \frac{\binom{\Nprim + s}{s}}{2^{256 (s - 1)}}.
\]

\begin{remark}[Two-Opening CMTDs-4 Specialization]
\label{rem:cmtd4-s2}
Setting $s = 2$ recovers the ordinary two-opening
CMTDs-4 statement of \cite{BH22}:
\[
  \Advcmtds{4}_\TW(\mathcal{A})
  \;\leq\;
  \frac{\Ntot(\Ntot + 1)}{2^{257}}
  + \frac{\binom{\Nprim + 2}{2}}{2^{256}}
  \;=\;
  \frac{\Ntot(\Ntot + 1) + (\Nprim + 1)(\Nprim + 2)}{2^{257}}.
\]
\end{remark}

\paragraph{Everywhere preimage resistance of $\HTW$.}
\[
  \AdvePre_\TW(\mathcal{A})
  \;\leq\;
  \frac{\Ntot(\Ntot + 1)}{2^{257}}
  + \frac{\Nprim + 1}{2^{256}}.
\]

\subsection{Concrete Security Claim}
\label{sec:concrete-claim}

\begin{corollary}[Concrete \TW{} Security]
\label{cor:concrete-security}
Under the implemented parameters $k = c = \nu = \tauroot = \tauleaf = 256$, every
adversary with work cap $\Ntot \leq 2^{63}$---and, in the multi-user case,
user cap $D \leq 2^{63}$---has advantage at most $2^{-128}$ in each of
\TW{}'s notions: multi-user indistinguishability (\Cref{thm:mu-ind}),
all-in-one AE (\Cref{cor:ae}), \textsf{INT-RUP} integrity
(\Cref{thm:int-rup}), the hash security of
$\HTW$---$s$-way collision resistance for every $s \geq 2$ in the sense
of \cite{BH22}, with the \cite{RS04} notions of standard and everywhere
second-preimage resistance inherited from collision resistance on every
fixed input slice, standard preimage resistance on every fixed input
slice, and everywhere preimage resistance
(\Cref{thm:coll,cor:second-preimage-rs04,thm:pre,thm:epre})---and the \cite{BH22}
CMTDs-4 committing bound it yields (\Cref{cor:cmtds4}).
\end{corollary}

\begin{proof}
The resource consequences of \Cref{sec:resources} give
$\Nprim \leq \Ntot$, $\Qv \leq \Ntot$, $\Nleaf \leq \Ntot$, and
$\Nprim+\Nleaf \leq \Ntot$. The user cap $D$ remains independent, since
created-but-idle users need not contribute to $N$.

At $\Ntot \leq 2^{63}$, each flat-sponge term---at $\Ntot$ or at
$\Nprim$---and each birthday-scale root or leaf term at $s=2$ is at most
$2^{-131} + 2^{-192}$, while each
linear root-preimage, key-guessing, verification, or leaf tag term with no
factor $D$ is at most a $2^{-193}$-scale term.
The fixed-slice Pre degradation is at most $2^{-256}$, since
$m_{a,\mu} \geq k+\nu = 512$.
Thus the largest non-multi-user bound is the $s=2$ collision bound, which is below
$4 \cdot 2^{-131} = 2^{-129}$; larger $s$ only decreases the
root-collision term. The all-in-one AE bound carries its two flat-sponge
terms as its only birthday-scale terms and lies below $2^{-129}$, and the
\textsf{INT-RUP} bound carries $\Lift_c(\Ntot)$ alone. In the multi-user
case, using $D \leq 2^{63}$,
\[
  \Advmuind_\TW(\mathcal{A})
  < \bigl(2^{-130} + 2^{-193}\bigr) + 2^{-130} + 2^{-192} + 2^{-131}
  < 2^{-128}:
\]
the two flat-sponge terms sum to at most
$2 \cdot (2^{-131} + 2^{-194}) = 2^{-130} + 2^{-193}$, the key-guessing
term satisfies $D \Nprim / 2^{256} \leq 2^{126}/2^{256} = 2^{-130}$, the
remaining linear terms satisfy
$(\Qv + \Nleaf)/2^{256} \leq 2^{64}/2^{256} = 2^{-192}$, and
$D(D-1)/2^{257} < 2^{-131}$; the total is below
$5 \cdot 2^{-131} + 2^{-192} + 2^{-193} < 2^{-128}$. These estimates
instantiate every closed form in \Cref{sec:concrete};
\Cref{sec:bound-discussion} examines the margin they leave.
\end{proof}

\begin{table}[t]
\centering
\caption{Concrete \TW{} security at the work cap $\Ntot \leq 2^{63}$
(with $D \leq 2^{63}$ for mu-ind), under
$k = c = \nu = \tauroot = \tauleaf = 256$ and the resource consequences
$\Nprim, \Qv, \Nleaf \leq \Ntot$ and $\Nprim + \Nleaf \leq \Ntot$ of
\Cref{sec:resources}. The Coll and
CMTDs-4 rows hold for every $s \geq 2$; larger $s$ only shrinks the
root-collision term.}
\label{tab:concrete-security}
\begin{tabular}{@{}lcc@{}}
\toprule
Notion & Dominant scale & At the cap \\
\midrule
All-in-one AE & $\Lift_c$ & $\leq 2^{-128}$ \\
Multi-user ind. & $D\,\KeyGuess_k(\Nprim)$ & $\leq 2^{-128}$ \\
\textsf{INT-RUP} & $\Lift_c$ & $\leq 2^{-128}$ \\
$\HTW$ collision (Coll) & $\Lift_c$, $\LeafColl$, $\RootColl$ & $\leq 2^{-128}$ \\
$\HTW$ second-preimage (Sec/eSec) & inherited from Coll & $\leq 2^{-128}$ \\
$\HTW$ preimage (Pre) & $\Lift_c$, $\LeafColl$ & $\leq 2^{-128}$ \\
$\HTW$ everywhere preimage (ePre) & $\Lift_c$ & $\leq 2^{-128}$ \\
CMTDs-4 \cite{BH22} & $\Lift_c$ and $\RootColl$ & $\leq 2^{-128}$ \\
\bottomrule
\end{tabular}
\end{table}

\subsection{Discussion of the Bounds}
\label{sec:bound-discussion}

The bounds of \Cref{sec:concrete} aggregate the loss terms of
\Cref{sec:loss-terms}, whose sizes at the work cap
\Cref{tab:concrete-security} records and whose dominance
\Cref{sec:concrete-claim} proves.

\paragraph{At the concrete cap.}
At $\Ntot \leq 2^{63}$, the hash bounds other than ePre and the
CMTDs-4 bound have flat-sponge
and root or leaf birthday terms on the same $2^{-131}$ scale; the
ePre and \textsf{INT-RUP} bounds retain the flat sponge loss
as their only birthday-scale term, and the single-user AE bound carries
its two flat-sponge terms, at $\Ntot$ and $\Nprim$, as its only
birthday-scale terms. Thus the flat sponge loss is not always the largest
summand: under the loose per-term cap, the $s=2$ root-collision term
$\binom{\Nprim+2}{2}/2^{256}$ is on the same scale as $\Lift_c(\Ntot)$.
The multi-user bound is the exception: its
hybrid-induced
term $D\,\KeyGuess_k(\Nprim)$ can reach $2^{-130}$ at the extreme cap
$D=\Nprim=2^{63}$. That corner is a property of the bound, not of any
admissible attack: the matching attack of the Tightness paragraph below
needs one accepted nonempty encryption per targeted user, and even the
empty-AD one-block case costs two duplex calls (root initialization and
closing root message). With $t$ targeted users and $q=\Nprim$ direct
guesses, the joint cap imposes $2t+q \leq 2^{63}$, so the feasible
key-guessing attack scale is about $2^{-133}$, reached near
$t=2^{61}$ and $q=2^{62}$. The remaining factor is hybrid slack from
charging created-but-idle users. Expressed in bytes, $\Ntot \leq 2^{63}$ corresponds to
at most $\rho \cdot 2^{63} = 167 \cdot 2^{63} < 2^{71}$ bytes of
plaintext, associated data, ciphertext, and aggregation through the
duplex (\Cref{sec:params}).

\paragraph{Margin at the cap.}
\Cref{cor:concrete-security} clears its target with a stated remainder:
at the joint corner $D = \Nprim = 2^{63}$ the multi-user estimate has
exact value $5 \cdot 2^{-131} + 2^{-194}$, since the corner forces
$N=0$ and hence $\Qv=\Nleaf=0$; this is just over five eighths of the
$2^{-128}$ budget. The largest of the remaining bounds, the $s = 2$
collision bound, stays below $4 \cdot 2^{-131}$
(\Cref{sec:concrete-claim}). The closeness at the corner is a proof
artifact rather than a consequence of the design. The corner itself is
unoccupiable: it forces $N = 0$, as the previous paragraph records, and
an adversary without construction computations receives identical
answers in both worlds---the primitive interface is the same real
permutation (\Cref{sec:games-pp})---so its advantage there is exactly
$0$. At resource splits an adversary can occupy, the multi-user
key-guessing witness reaches about the $2^{-133}$ scale, while the
root-collision and flat-sponge tightness witnesses can still reach the
$2^{-131}$ scale. The intervening distance to the proof bound is
consumed by the two flat-sponge composition terms and the hybrid slack
recorded above, the overhead of the indifferentiability route. The margin also widens away
from the corner: at $D \leq 2^{32}$, the $D$-indexed terms drop below
$2^{-160}$ and the multi-user bound falls to its two flat-sponge terms,
below $2^{-130} + 2^{-159}$, two bits under the target. A direct
keyed-duplex analysis would be expected to replace the quadratic
$\Lift_c$ terms by finer construction--primitive cross terms, much as
the keyed-duplex bounds synthesized by \cite{Men23} improve on the
sponge-indifferentiability bound underlying the original SpongeWrap
analysis of \cite{BDPV11}; \Cref{sec:limitations} discusses that route
and the gaps it would need to close.

\paragraph{Reading the work cap.}
\Cref{cor:concrete-security} is a success-probability statement at a fixed
cap, not a statement that security is exhausted there. The closed forms of
\Cref{sec:concrete} remain explicit functions of the resources beyond
$2^{63}$, and their dominant terms are birthday-type in $\Ntot$ on the
$256$-bit capacity, so the guarantees degrade quadratically rather than
abruptly: at $\Ntot = 2^{100}$, for example, the $s = 2$ collision bound is
still below $2^{-55}$. The bounds become vacuous only as $\Ntot$ approaches
$2^{c/2} = 2^{128}$, where $\Lift_c(\Ntot)$ exceeds $1/2$
(\Cref{sec:loss-terms}). The cap $\Ntot \leq 2^{63}$ is therefore a
data-scale choice: a single work bound at which every notion's advantage
lies below the $2^{-128}$ target simultaneously, not the largest workload
the mode tolerates.

\paragraph{Tightness.}
\begin{itemize}
\item \emph{$\KeyGuess_k(Q) = Q / 2^k$.} Tight at constants: a direct
  random-guess attack on a $k$-bit key matches the per-query success
  probability $1 / 2^k$. The multi-user bound's dominant term
  $D\,\KeyGuess_k(\Nprim) = D\Nprim / 2^k$ is likewise tight at the
  leading constant in the theorem's free-parameter setting, using the
  per-user nonce discipline of \Cref{sec:mu-ind}, which permits one
  nonce value across users. The adversary queries every targeted user
  on the same nonce $U$, empty associated data, and a known nonempty
  message---construction work charged to $N$, which does not enter the
  term---so each user's first ciphertext block reveals the keystream
  emitted by that user's root initialization call
  (\Cref{tab:transcript-schedule}). Each direct guess is then a single
  forward primitive query at the root initialization transcript
  $\prfxroot \concat K' \concat U$ for a fresh candidate key $K'$,
  whose output is compared against all $D$ collected keystream blocks
  at once; it matches some targeted key with probability $D / 2^k$, so
  $\Nprim$ guesses recover a user key---and distinguish---with
  probability $\approx D\Nprim / 2^k$, matching the bound. Under the
  joint cap the construction work for those target encryptions lowers
  the feasible scale to about $2^{-133}$, as noted in the concrete-cap
  paragraph above.
\item \emph{$\RootColl_\tau(Q, s) = \binom{Q + s}{s} / 2^{\tau (s -
  1)}$.} Tight up to lower-order terms: it is the standard $s$-way
  multi-collision bound on a $\tau$-bit projection of $\ROflat$, and
  a generic birthday-style adversary against the random oracle
  matches the leading-order term. It is the root term of the collision
  bound ($\HTW$ Coll, \Cref{thm:coll}). Its $\Nprim$ counts every
  primitive query, though by Output-Point Separation
  (\Cref{cor:output-sep}) only those whose induced $\ROflat$ call lands at
  a root tag output point ($\Rtag$, or $\Rmsg^+$ for single-chunk
  ciphertexts) can enter the candidate sequence; a class-filtered count
  would be no larger, and potentially smaller, but the displayed bound keeps the
  $\Nprim$ form for uniformity with \Cref{cor:ae,thm:mu-ind}, the birthday
  threshold at $\tauroot = 256$ being unaffected.
\item \emph{$\RootPre_\tau(Q) = (Q + 1) / 2^\tau$.} Tight at constants: a
  generic preimage search against a fixed $\tau$-bit target matches the
  per-query $2^{-\tau}$ rate. It is the root term of the standard and
  everywhere preimage bounds, where a single target tag yields a linear term.
\item \emph{Leaf tag term $\Nleaf / 2^{\tauleaf}$.} Tight at constants
  by a nonce-respecting attack: encrypt a single two-chunk message
  under nonce $U$, obtaining ciphertext with one leaf chunk and root
  tag $T$, then submit non-replay verification queries that each
  replace the leaf chunk by fresh bytes of the same length, keeping
  $T$. Each query recomputes the leaf tag at the hidden slot
  $(K, U, 1)$ and accepts whenever that fresh tag equals the
  encryption's hidden leaf tag $T_1$, since the aggregation
  string---and hence the recomputed root tag---is then unchanged. Each
  verification-side final leaf transcript therefore forges with
  probability $2^{-\tauleaf}$, matching the bound's per-transcript
  rate.
\item \emph{Known-key leaf collision $\LeafColl_\tau(Q)$.} Genuinely
  birthday-tight: when keys are adversary-supplied, direct queries can
  target leaf tag points, and a birthday search over chosen leaf bodies
  yields two distinct $\HTW$ inputs sharing a leaf tag and hence
  colliding openings; the collision proof uses $Q=\Nprim+\Nleaf$.
\item \emph{Flat sponge lift.} The term
  $\Lift_c(\Ntot) = \Ntot(\Ntot + 1) / 2^{c + 1}$ dominates the
  random-permutation differentiating bound $f_P(\Ntot)$ of
  \cite[Lemma~5]{BDPV08} via \Cref{lem:lift-bound}. The
  differentiating bound $f_P(\Ntot)$ is itself matched at leading order
  by the standard inner-collision differentiator, which succeeds
  against every simulator because the ideal-world construction
  interface is the random oracle regardless of simulator strategy; any
  residual conservatism in $\Lift_c(\Ntot)$ lies in composing
  indifferentiability into the specific keyed games, not in an
  improvable simulator bound, and is not specific to \TW{}. The
  real-vs-ideal games pay the same differentiating bound once more, at
  $\Nprim$, to exchange the ideal-side simulator for the real
  permutation (\Cref{lem:sim-perm-proximity}); since
  $\Nprim \leq \Ntot$, that term never dominates the $\Ntot$ term.
\end{itemize}

Cross-scheme bound-level comparisons against prior sponge AEADs and
the committing AEAD line of \cite{BH22} are deferred to
\Cref{sec:compare-table,sec:compare-sponge,sec:compare-commit}.

\section{Performance Evaluation}
\label{sec:performance}

\TW{} fixes the root/leaf transcript but leaves leaf scheduling to the
implementation. This section evaluates that separation: the same ciphertext and
tag are produced by scalar and vectorized implementations, while complete leaf
chunks can be batched to recover practical long-message throughput. The
comparison baseline is hardware-accelerated AES-128-GCM on the same hosts, so
the measurements should be read as an implementation sanity check and
same host reference point, not as a claim that \TW{} is faster than AES-GCM.

\subsection{Optimized Implementations}
\label{sec:impl}

We analyze a Go implementation of \TW{}\footnote{The implementation is available
at \url{https://github.com/codahale/treewrap}.} with optimized assembly paths
for the performance-critical permutation and absorb/squeeze loops. We
compare three vectorized assembly paths: \texttt{amd64} AVX-512,
\texttt{amd64} AVX2, and \texttt{arm64} NEON using the Armv8.2 SHA-3
extension. The framing, domain separation, padding, tree bookkeeping, and
dispatch remain in shared Go code. A pure Go path uses a scalar permutation
and a scalar eight-instance loop, and serves both as a reference path for
cross-checking the assembly cores and as a fallback on other architectures.

The tree structure of \TW{} exposes two distinct permutation workloads. The
root duplex---which absorbs the associated data, processes chunk~$0$, and
absorbs the leaf tag aggregation transcript---is inherently sequential, since
each call consumes the previous call's output. It is driven by a
\emph{single-duplex} ($\times 1$) core, except for its chunk-$0$ message phase.
The leaf duplexes operate on disjoint state and disjoint inputs
(\Cref{sec:tw128}), so complete leaf chunks are batched and run through an
\emph{eight-way} ($\times 8$) core. Leftover runs of two to seven complete
chunks use narrower remainder kernels, while a single leftover complete chunk,
a partial trailing chunk, and the count-aggregation tail fall back to the
single-duplex core.

Chunk~$0$ itself runs on the batched kernels rather than serially, a technique
we refer to as \emph{chunk~$0$ fusion}. The root duplex's chunk-$0$ message phase
has the same kernel-visible schedule as a leaf chunk and differs only in its
initial state. The implementation therefore seeds lane~$0$ of the first batched
call with the root state, runs chunk~$0$ alongside the first leaf chunks, then
writes lane~$0$'s output state back to the root duplex before the root absorbs
the leaf tags from that same call. This avoids a serial pass over chunk~$0$ and
is the reason the cycles-per-byte curve drops as soon as complete leaves are
present (\Cref{sec:throughput}). \Cref{app:kernels} gives the lane-major layout,
remainder-kernel strategies, lane~$0$ writeback details, and instruction-level
kernel structure.

\paragraph{Architecture-specific cores.}
On \texttt{amd64} the optimized implementation selects an AVX-512 or AVX2 core
once at startup from the \texttt{AVX512VL} CPUID feature bit. On \texttt{arm64}
the optimized implementation uses NEON with the Armv8.2 SHA-3 extension
(\texttt{FEAT\_SHA3}, available on Apple Silicon and Neoverse V1/V2 and N2).
Both architecture-specific paths provide a single-duplex core and batched leaf
kernels; the assembly stays below the mode boundary, so framing, domain
separation, padding, and tree bookkeeping remain shared with the reference path.

\paragraph{Data-independent execution.}
The scalar reference and all three vectorized implementations avoid
secret-dependent branches and secret-dependent memory or vector indices for the
key, nonce, associated-data contents, and plaintext or ciphertext contents. The
permutation and the XOR-based absorb/squeeze loops are straight-line with
respect to data values. The only indexed table is the round-constant array,
addressed by the public round number; dispatch, batched loads, stores, gathers,
scatters, and inactive-lane handling are controlled by public lengths, chunk
counts, and fixed chunk strides. \Cref{app:kernels} records the
architecture-specific addressing details.

\subsection{Benchmark Methodology}
\label{sec:bench-method}

We measure on two hosts, one per architecture. The \texttt{amd64} host is a
Google Compute Engine \texttt{c4-standard-4} instance with an Intel Xeon Platinum
8581C (Emerald Rapids, $2.30$\,GHz base, AVX-512 including \texttt{AVX512VL},
\texttt{AVX512\_VBMI2}, \texttt{GFNI}, and \texttt{VAES}), two physical cores
with simultaneous multithreading enabled, $48$\,KiB L1d and $32$\,KiB L1i per
core, $2$\,MiB L2 per core, and a $260$\,MiB shared L3, running Debian~12 (kernel
6.1) under KVM. The guest does not expose a \texttt{cpufreq} interface, so the
frequency governor and turbo state are not under our control; we mitigate the
resulting variance statistically (below). The \texttt{arm64} host is an Apple M4
Pro ($10$ performance cores and $4$ efficiency cores; per performance core,
$128$\,KiB L1i, $64$\,KiB L1d, and a $16$\,MiB shared L2) running macOS~26.5.1
(Darwin~25.5), with \texttt{FEAT\_SHA3} present. macOS exposes no userspace
frequency control.

Measurements are collected by a standalone harness (\texttt{cmd/cpb}) that locks
its goroutine to an OS thread and, for each algorithm, operation, and message
length, first calibrates an iteration count that fills at least $100$\,ms, then
records $100$ samples. It reports the median cycles per byte and median
wall-clock throughput together with the interquartile range as a dispersion
measure (\Cref{tab:dispersion}). Both metrics are read from the same timed
loop, so the cycles-per-byte and throughput figures for a given measurement
come from one interleaved run rather than two separately scheduled ones.

Each timed call uses a fixed all-zero key and nonce, empty associated data, and
preallocated source and destination buffers. Open measurements decrypt and
verify a ciphertext prepared before timing; seal and open reset the destination
slice length inside the timed loop but do not allocate per call. The reported
figures are therefore warm-cache, steady-state rates for message processing
under a fixed AEAD object. The single-threaded measurement leaves the SMT
sibling of the \texttt{amd64} host idle. All reported figures are
single-threaded: they exercise the SIMD axis of \TW{}'s parallelism, and
multicore scaling remains unmeasured. We report a sweep of seven message
lengths from $1$\,B to $16$\,MiB.

The cycle counters differ by architecture, and this affects how the
cycles-per-byte figures may be read. On \texttt{amd64} the harness reads
\texttt{RDTSC}; on this part the time-stamp counter is invariant and ticks at the
$2.30$\,GHz reference frequency, so the reported figures are \emph{reference}
cycles per byte and are unaffected by turbo (and convert directly to wall-clock
time). On \texttt{arm64} the harness reads \texttt{CNTVCT\_EL0} and scales it by
the ratio of the peak core frequency to \texttt{CNTFRQ\_EL0}; since Apple exposes
no frequency API, the peak is auto-detected by timing a dependent integer-add
chain and taking the top observed P-state, yielding \emph{core} cycles per byte.
Consequently the cycles-per-byte figures are comparable within an architecture
and against the same-host AES-128-GCM baseline, but not directly across
architectures; absolute throughput (\Cref{tab:throughput}) is the cross-platform
metric.

The baseline is Go's AES-128-GCM implementation
(\texttt{crypto/\allowbreak aes} with \texttt{crypto/\allowbreak cipher}).
It uses AES-NI and \texttt{PCLMULQDQ} on \texttt{amd64}, and the AES
and \texttt{PMULL} instructions on \texttt{arm64}.
The AES block cipher, GCM wrapper, and \TW{} AEAD object are constructed once
before timing, so the measurements exclude key schedule and GHASH key-power
precomputation. Each timed \TW{} call still performs the root duplex
initialization inside the call. Both directions are measured for each scheme;
the cycles-per-byte table reports both, and the throughput table reports the
encryption direction.

\subsection{Throughput}
\label{sec:throughput}

\Cref{tab:cpb} reports cycles per byte across the message-length sweep and
\Cref{tab:throughput} the corresponding wall-clock throughput in the bulk
regime. The main pattern in the \TW{} rows is the descent produced by the
tree-parallel leaf core. A message of at most one chunk has no complete leaf
chunk for chunk~$0$ to share a kernel with, so it runs entirely on the
single-duplex core, and the $8$\,KiB column sits at that core's rate. Past one
chunk, chunk~$0$ fusion (\Cref{sec:impl}) keeps every complete chunk on the
batched kernels, and the descent tracks lane occupancy.

On the \texttt{amd64} AVX-512 path, \TW{} encryption falls
$1.81 \to 0.73 \to 0.39$ cycles per byte across the $8$, $32$, and $64$\,KiB
columns and then holds at $0.38$--$0.39$ through $16$\,MiB. The AVX2 path has
the same shape at lower throughput: it runs the eight instances as two groups
of four rather than in a single $8$-wide register, leaving the AVX-512 core
about $2.9\times$ faster in bulk. The \texttt{arm64} curve flattens earlier
because the eight-way core processes its batch as four two-instance pairs; a
$32$\,KiB message, chunk~$0$ and three leaves as two full pairs, is already
close to the long-message rate. In wall-clock terms \TW{} reaches
$6.00$\,GB/s on the \texttt{amd64} AVX-512 host and $4.93$\,GB/s on the M4 Pro
for $16$\,MiB messages.

\begin{table}
\centering
\caption{Median cycles per byte ($100$ samples). The \texttt{amd64} figures are
\texttt{RDTSC} reference cycles at the $2.30$\,GHz invariant time-stamp counter;
the \texttt{arm64} figures are core cycles at the auto-detected peak performance-core
frequency. Cycle counts are not directly comparable across architectures
(\Cref{sec:bench-method}).}
\label{tab:cpb}
\begin{tabular}{l rrrrrrr}
\toprule
& \multicolumn{7}{c}{Message length} \\
\cmidrule(l){2-8}
& $1$\,B & $64$\,B & $8$\,KiB & $32$\,KiB & $64$\,KiB & $1$\,MiB & $16$\,MiB \\
\midrule
\multicolumn{8}{l}{\emph{Intel Xeon 8581C (Emerald Rapids), AVX-512}} \\
\quad \TW{} encrypt      & 571 & 9.04 & 1.81 & 0.73 & 0.39 & 0.38 & 0.38 \\
\quad \TW{} decrypt      & 605 & 9.55 & 1.82 & 0.73 & 0.39 & 0.38 & 0.38 \\
\quad AES-128-GCM seal   & 119 & 2.74 & 0.30 & 0.29 & 0.29 & 0.29 & 0.29 \\
\quad AES-128-GCM open   & 99 & 2.17 & 0.29 & 0.28 & 0.28 & 0.28 & 0.29 \\
\midrule
\multicolumn{8}{l}{\emph{Intel Xeon 8581C (Emerald Rapids), AVX2}} \\
\quad \TW{} encrypt      & 647 & 10.18 & 2.21 & 1.15 & 1.10 & 1.09 & 1.10 \\
\quad \TW{} decrypt      & 682 & 10.70 & 2.21 & 1.06 & 1.10 & 1.08 & 1.06 \\
\midrule
\multicolumn{8}{l}{\emph{Apple M4 Pro, NEON\,+\,\texttt{FEAT\_SHA3}}} \\
\quad \TW{} encrypt      & 620 & 9.88 & 1.97 & 0.92 & 0.89 & 0.89 & 0.91 \\
\quad \TW{} decrypt      & 671 & 10.65 & 2.04 & 0.90 & 0.87 & 0.87 & 0.89 \\
\quad AES-128-GCM seal   & 113 & 2.84 & 0.44 & 0.43 & 0.43 & 0.44 & 0.46 \\
\quad AES-128-GCM open   & 123 & 2.56 & 0.42 & 0.41 & 0.41 & 0.42 & 0.43 \\
\bottomrule
\end{tabular}
\end{table}

\begin{table}
\centering
\caption{Measured wall-clock throughput (GB/s, $\mathrm{GB} = 10^9$ bytes) from
the same \texttt{cmd/cpb} run as \Cref{tab:cpb}, encryption direction. The
\texttt{amd64} AES-128-GCM row is the AVX-512 host and is unaffected by the
\TW{} core selection.}
\label{tab:throughput}
\begin{tabular}{l rrrrr}
\toprule
& $8$\,KiB & $32$\,KiB & $64$\,KiB & $1$\,MiB & $16$\,MiB \\
\midrule
\multicolumn{6}{l}{\emph{Intel Xeon 8581C (Emerald Rapids)}} \\
\quad \TW{} (AVX-512)  & 1.27 & 3.17 & 5.93 & 6.10 & 6.00 \\
\quad \TW{} (AVX2)     & 1.04 & 2.00 & 2.10 & 2.10 & 2.08 \\
\quad AES-128-GCM      & 7.75 & 8.00 & 8.05 & 8.00 & 7.87 \\
\midrule
\multicolumn{6}{l}{\emph{Apple M4 Pro}} \\
\quad \TW{} (NEON+SHA3) & 2.27 & 4.84 & 5.04 & 5.01 & 4.93 \\
\quad AES-128-GCM       & 10.04 & 10.27 & 10.33 & 10.07 & 9.71 \\
\bottomrule
\end{tabular}
\end{table}

\paragraph{Dispersion.}
Each cell of \Cref{tab:cpb,tab:throughput} is the median of $100$ samples;
\Cref{tab:dispersion} summarizes their spread as the interquartile range
(IQR) relative to the median. The medians are stable: across the sweep the
relative IQR of the \TW{} cycles-per-byte measurements is at most $0.7\%$
on the \texttt{amd64} AVX-512 path, $0.9\%$ on AVX2, and $1.9\%$ on the M4
Pro, with per-path medians of $0.3\%$, $0.4\%$, and $1.2\%$. The full
min--max range is wider at isolated points---up to $4.9\%$, $12.0\%$, and
$16.7\%$ on the three paths---because turbo, governor, and scheduler states
are not under our control (\Cref{sec:bench-method}); these are tail
excursions that the IQR excludes. We therefore report medians throughout
and treat the IQR as the dispersion measure.

\begin{table}
\centering
\caption{Cycles-per-byte measurement dispersion for \TW{} across the
seven-length sweep (seal and open, $100$ samples per point), as the
interquartile range relative to the median. The final column is the widest
single-point min--max range; it is excluded from the IQR as a tail effect.}
\label{tab:dispersion}
\begin{tabular}{l rrr}
\toprule
& \multicolumn{2}{c}{Relative IQR} & Max \\
\cmidrule(lr){2-3}
Path & median & max & range \\
\midrule
Xeon 8581C, AVX-512 & $0.3\%$ & $0.7\%$ & $\phantom{0}4.9\%$ \\
Xeon 8581C, AVX2    & $0.4\%$ & $0.9\%$ & $12.0\%$ \\
Apple M4 Pro        & $1.2\%$ & $1.9\%$ & $16.7\%$ \\
\bottomrule
\end{tabular}
\end{table}

\subsection{Latency}
\label{sec:latency}

For short messages in this benchmark, the cost is set by the serial root
duplex and is independent of the parallel leaf machinery. With empty associated
data, a message of at most $\rho = 167$ bytes occupies a single rate block and
never enters leaf mode, so it costs exactly two permutations: one to initialize
the root duplex and one to close the message block and extract the tag. The
$1$\,B and $64$\,B columns of \Cref{tab:cpb} bear this out: per-message latency
is nearly flat across them at roughly $571$ reference cycles on \texttt{amd64}
($571$ at $1$\,B versus $9.04 \times 64 \approx 578$ at $64$\,B) and $620$
core cycles on \texttt{arm64} ($620$ versus $9.88 \times 64 \approx 633$),
i.e.\ about $285$ and $310$ cycles per \Keccakp{} respectively. Because such
messages bypass the leaf core entirely, leaf-level parallelism imposes no
latency penalty on them.

The latency/throughput trade-off appears instead in the intermediate regime,
which chunk~$0$ fusion has narrowed to roughly the one-to-eight-chunk band. The
worst per-byte cost in the sweep is the single-chunk regime at $8$\,KiB
($1.8$--$2.2$ cycles per byte across the three paths), where the message is the
chunk-$0$ phase alone on the single-duplex core. Once the message has complete
leaf chunks for chunk~$0$ to share a kernel with, the cost falls quickly: by
$32$\,KiB \TW{} is within a factor of two of its AVX-512 floor and close to its
\texttt{arm64} floor, and the floor itself is reached once the eight-way core
fills at $64$\,KiB. The root duplex also bounds the tail of a long
message: the final tag cannot be produced until chunk~$0$ has been processed and
every leaf tag has been absorbed into the aggregation transcript. That absorption
is cheap---each rate block of the root absorbs just over five $32$-byte leaf tags, so a
$16$\,MiB message ($\approx 2000$ leaves) adds on the order of $400$ root
permutations against roughly $10^5$ permutations of leaf work, under $0.5\%$
overhead---so the serial root does not materially cap bulk throughput.

\subsection{AES-128-GCM Baseline}
\label{sec:compare-perf}

We compare against the one baseline measured under the same harness on the same
hosts, the Go standard library's hardware-accelerated AES-128-GCM
(\Cref{tab:cpb,tab:throughput}). The harness times \TW{} and AES-128-GCM
back-to-back at each length and reads both metrics from the same loop, so the
cycles-per-byte and throughput comparisons hold under identical conditions and
agree. In the bulk regime AES-128-GCM is faster on both architectures, as
expected for a primitive with dedicated AES and carryless-multiply instructions:
at $16$\,MiB it runs at $0.29$ cycles per byte on \texttt{amd64} and $0.46$ on
\texttt{arm64}, against $0.38$ and $0.91$ for \TW{}---about $1.3\times$ and
$2.0\times$ faster. The wall-clock throughputs show the same margins: \TW{}
reaches $6.00$ and $4.93$\,GB/s against AES-128-GCM's $7.87$ and
$9.71$\,GB/s. The wider \texttt{arm64} gap reflects the strength of Apple's AES
and \texttt{PMULL} units. \TW{} attains these rates with a permutation that has
no dedicated hardware support, relying only on the SHA-3 extension on
\texttt{arm64} and on general SIMD on \texttt{amd64}.

AES-128-GCM is also faster for short messages in this steady-state measurement:
at $1$\,B it costs $119$ and $113$ cycles on the two hosts against \TW{}'s
$571$ and $620$. This is the expected result when AES-GCM key setup is
amortized outside the timed call. In the sub-bulk regime AES-128-GCM is several
times faster at one chunk ($0.30$ versus $1.81$ cycles per byte at $8$\,KiB on
\texttt{amd64}), but the gap narrows once the fused kernels engage: about
$2.5\times$ at $32$\,KiB and the bulk ratio from $64$\,KiB on. \TW{}'s intended
value relative to AES-128-GCM lies in the tag-centric hash and committing
properties analyzed in \Cref{sec:root-tag-hash-security}, with the baseline
AEAD bounds of \Cref{sec:aead-security}, rather than in raw speed.

\section{Discussion}
\label{sec:discussion}

\subsection{Comparison Table}
\label{sec:compare-table}

\Cref{tab:scheme-compare} situates \TW{} against representative authenticated
encryption schemes along the axes that distinguish it: the underlying primitive,
whether the mode is parallelizable, the key, nonce, and tag sizes, associated data
support, nonce-misuse resistance, and the published committing result or attack
most relevant to the scheme. The rows are representative rather than
exhaustive: sponge/duplex relatives, a Farfalle-based parallel design, AES
baselines, a misuse-resistant AES baseline, and a generic committing transform.
The comparison is structural, not a total security ranking; \TW{}'s own
concrete bounds and the regimes in which each term dominates are treated in
\Cref{sec:bound-discussion}, and its measured throughput against
hardware-accelerated AES-128-GCM in \Cref{sec:compare-perf}. The committing
column should be read within the named notions only, and a negative entry there
does not assert a confidentiality or integrity break of the underlying AEAD. It
records \TW{}'s $128$-bit \textsf{CMTDs-4} bound derived from tag hash
security, Ascon's roughly $64$-bit committing bound \cite{KSW23}, published
low-cost attacks where known \cite{KSW23,CR22,ADGKLS22}, and the
hash-collision level inherited by \textsf{CTX}.
The table does not separately track compact (tag-only) commitment, since
that property has not been systematically analyzed for the comparison
schemes.

\begin{table}[t]
\centering
\footnotesize
\caption{Structural comparison of \TW{} with representative AEADs. Sizes are in
bytes; AES-GCM-SIV also admits a $32$-byte key, and Kravatte-SAE's key and
nonce lengths are parameters (keys of up to $40$ bytes, arbitrary-length
nonces). ``Misuse'' is nonce-misuse
resistance in the sense of \cite{RS06}; ``Commit'' records a published
committing result or attack in the relevant committing notion, with ``broken''
denoting a published low-cost committing attack, ``hash CR'' denoting the
collision-resistance level of the hash used by the transform, and ``open''
marking a scheme with no published committing analysis. A dash means the entry
depends on the base scheme. These entries are notion-specific summaries, not
broad AE security claims. \TW{}'s concrete bounds appear in
\Cref{sec:bound-discussion}.}
\label{tab:scheme-compare}
\begin{tabular}{@{}llccccc@{}}
\toprule
Scheme & Primitive & Par. & K/U/T & AD & Misuse & Commit \\
\midrule
\TW{}                        & \Keccakp   & yes & $32/32/32$ & yes & no  & CMTDs-4 \\
Ascon-128a \cite{DEMS21}     & Ascon-$p$ ($320$ b) & no  & $16/16/16$ & yes & no  & $\approx64$ b \\
Xoodyak \cite{DHPVV20}       & Xoodoo ($384$ b)   & no  & $16/16/16$ & yes & no  & $2^{17}$ atk. \\
Kravatte-SAE \cite{BDHPVV17} & \Keccak[1600, 6]   & yes & $\leq 40$/var./$16$ & yes & no & open \\
AES-GCM \cite{Dwo07}         & AES                & yes & $16/12/16$ & yes & no  & broken \\
AES-GCM-SIV \cite{GLL17}     & AES                & yes & $16/12/16$ & yes & yes & broken \\
nAE $+$ \textsf{CTX} \cite{CR22} & any nAE        & ---  & ---        & yes & --- & hash CR \\
\bottomrule
\end{tabular}
\end{table}

\subsection{Comparison with Prior Sponge AEADs}
\label{sec:compare-sponge}

\TW{} belongs to the single-permutation duplex AEAD family introduced in
\cite{BDPV11}. Representative lightweight members of that family, Ascon
\cite{DEMS21} and Xoodyak \cite{DHPVV20}, target $128$-bit security on small
permutations ($320$ and $384$ bits) with a strictly serial duplex: each call
depends on the previous one, so there is no mode-level parallelism to exploit.
\TW{} instead pays for a larger permutation, \Keccakp, and a $256$-bit
capacity---the same round count as KangarooTwelve \cite{BDPVVV16}---in exchange
for a tree that processes leaf chunks in parallel and for the headroom that
meets the $2^{-128}$ target at the stated work cap. The cost of this choice is
visible in \Cref{sec:performance}: on short, single-leaf inputs \TW{} is slower
per byte than a compact serial duplex would be, and its advantage appears only
once a message spans enough leaves to fill the vector units.

\TW{} reuses KangarooTwelve's final-node/leaf split but not its Sakura coding
\cite{BDPV14}. KangarooTwelve uses Sakura to make the final node parseable as a
tree-hash transcript \cite{BDPVVV16}. In \TW{}, the analogous chaining values
are typed leaf tags: each leaf duplex is initialized with its leaf prefix and
index, and the root absorbs leaf tags in index order followed by the leaf count
(\Cref{sec:tw128,sec:domain-sep}). Sakura's structural role is therefore
handled by mode initialization and domain separation rather than by a tree
coding. The transcript separation lemmas make this explicit
(\Cref{lem:transcript-inj,cor:output-sep}); the hash collision bound sees the
leaf encoding only through the explicit leaf tag collision term of
\Cref{thm:coll}.

Keyak \cite{BDPVV16}, whose Motorist mode is already parallelizable, is the
closest prior duplex AEAD in spirit; \Cref{sec:related} contrasts its
wire-visible, encryptor-and-decryptor-shared degree of parallelism with \TW{}'s
run-time leaf choice.
Strobe \cite{Ham17} is a serial sponge framework aimed at whole protocols rather
than at a standalone AEAD, and like Ascon and Xoodyak it does not parallelize.

Farfalle \cite{BDHPVV17} leaves the duplex setting altogether: it is a
parallel pseudorandom function whose compression layer sums
masked-and-permuted input blocks into an accumulator and whose expansion
layer generates output from a rolling state, with authenticated encryption
supplied by the Farfalle-SAE and Farfalle-SIV modes on top
(\Cref{sec:related}). Its Kravatte instance runs \Keccak[1600, 6], and the
parallelism available to an implementation is arbitrary, as in \TW{}. The two
designs price their permutations for different adversary models. Farfalle's
six-round compression layer is analyzed in a PRF setting where inputs are
randomized by secret masks and an adversary does not see both the input and
output of the same permutation call \cite{BDHPVV17}. \TW{}'s tag-hash claims
are adversarial-initialization claims: in the hash and committing games, the
adversary may choose the key as part of the input being hashed. \TW{} therefore
keeps the strict duplex object, uses the hash-grade round count discussed in
\Cref{sec:permutation-cryptanalysis}, and takes its parallelism from the tree
topology rather than from a parallel PRF compression layer.

Committing security is not unique to \TW{} in this family: a dedicated analysis
of the sponge-based NIST lightweight finalists proves Ascon committing while
breaking others \cite{KSW23}, and some constructions, such as Strobe and the
Farfalle modes, have not been examined in that setting at all. What
distinguishes \TW{} is that the authentication tag is analyzed directly as a
hash output for the full input. Collision, second-preimage, preimage, and
everywhere preimage bounds for that tag then yield the compact commitment
statements, including the $128$-bit full-input \textsf{CMTDs-4} bound, as
mode-native consequences without leaving the single-permutation duplex setting.

\subsection{Comparison with Committing AEAD Constructions}
\label{sec:compare-commit}

The dominant route to committing AEAD is a generic transform over a
non-committing base scheme. \cite{BH22} give the \textsf{UtC}, \textsf{RtC}, and
\textsf{HtE} transforms, which add key- or context-commitment to an arbitrary
AEAD, and \cite{CR22} give \textsf{CTX}, which makes any nonce-based AE scheme
commit to its full context at the cost of a single hash over a short string and
with no ciphertext expansion. These transforms are inexpensive and broadly
applicable; \textsf{CTX} in particular is nearly free. It is nevertheless a
committing-AE transform, not a tag-hash transform: it binds the transformed
ciphertext, but does not make the tag a standalone hash output for the full
input. \TW{} differs not by undercutting their cost but by integration. The
same permutation provides confidentiality, integrity, and the tag hash output;
no separate commitment primitive, key-derivation step, or auxiliary hash is
invoked. The multi-input full-commitment bound follows from the same root tag
collision bound used for hash security (\Cref{sec:root-tag-hash-security}), and
\Cref{thm:coll} applies even when the ciphertext bodies used as openings
differ, up to the additional leaf tag birthday term. This places \TW{} alongside
the dedicated committing primitives---the compactly committing AEAD of
\cite{GLR17} and the encryptment of \cite{DGRW18}---which likewise build
commitment in rather than bolt it on, but which were designed for message
franking rather than as general nonce-based AEADs. The contrast with transformed
conventional schemes is sharpened by the attacks of \Cref{sec:related}: base
schemes such as GCM and OCB are not committing on their own \cite{CR22}, and a
transform is the only way to make them so.

Closest to \TW{} here is the recent work of \cite{DHMVV24}, which builds
nonce-based AE directly on the SHA-3 functions \textsf{SHAKE} and
\textsf{TurboSHAKE}---the latter using the same round-reduced \Keccak{}
permutation as \TW{}---and likewise achieves \textsf{CMT-4} natively, from the
collision resistance of the underlying function rather than from a transform. The
two designs make different structural choices: \cite{DHMVV24} targets
session-based communication, chaining a stream of messages through intermediate
tags and generalizing the keyed duplex into a \emph{duplex cipher}, whereas \TW{}
is a single two-level tree whose independent leaves let the implementation
choose parallelism at runtime with constant memory (\Cref{sec:related,sec:impl}).
Because a session
chain processes each message serially, \TW{}'s principal advantage is throughput
on large messages (\Cref{sec:performance}): its parallel leaves fill vector units
a serial duplex leaves idle, at the cost of the session support \cite{DHMVV24}
provides and \TW{} does not. Both demonstrate
that committing AE needs no add-on transform once it is built on a
well-scrutinized \Keccak{} permutation.

\subsection{Cryptanalysis of the Underlying Permutation}
\label{sec:permutation-cryptanalysis}

Every bound in this paper models \Keccakp{} as a uniformly random permutation
(\Cref{sec:security-model}). That modeling choice is supported by
cryptanalytic scrutiny rather than by proof, so this subsection summarizes the
third-party record and the safety margin it indicates. The Keccak round
function has been unchanged since 2008, and \TW{} uses the permutation at the
round count introduced by KangarooTwelve \cite{BDPVVV16} and later exposed
directly as TurboSHAKE \cite{BDHPVVV23}, whose designers maintain a survey of
third-party cryptanalysis. TurboSHAKE128 shares \TW{}'s rate and capacity
($r = 1344$, $c = 256$) as well as its permutation, so attacks on it and on
round-reduced SHAKE128 transfer to \TW{}'s geometry directly.

In the unkeyed setting, the deepest collision attacks on any SHA-3 or SHAKE
instance reach $6$ of \TW{}'s $12$ rounds: a SAT-aided attack on $6$-round
SHAKE128 with time $2^{123.5}$ \cite{GLST22}, and internal-differential
attacks reaching $6$ rounds of SHAKE256 and $5$ rounds of SHA3-384
\cite{ZHL24}. Preimage attacks reach $5$ rounds \cite{ZHL25}. Against the
collision- and preimage-style attacks most relevant to the root tag hash
claims of \Cref{sec:root-tag-hash-security}, the permutation therefore
retains a margin of $6$ of $12$ rounds---the assessment under which the
KangarooTwelve round count was chosen and later reaffirmed
\cite{BDPVVV16,BDHPVVV23}. The structural distinguisher reaching the most
rounds of the permutation itself, the SymSum property on $9$ rounds reported
in the survey of \cite{BDHPVVV23}, has, like the zero-sum distinguishers
before it, not been extended to an attack on any sponge or duplex mode.

The keyed setting is closer to \TW{}'s AE games: a hidden key, with
adversary-controlled nonces and messages entering through the duplex rate.
The attacks reaching the most rounds here are cube and conditional cube
attacks. Dinur et al.~break up to $9$ rounds of Keccak-based stream ciphers
and MACs faster than generic search, with keystream prediction on $8$ rounds
at time and data $2^{128}$ and on $9$ rounds (of a $512$-bit-key variant) at
$2^{256}$ \cite{DMPSS15}. With MILP-optimized conditional cubes,
nonce-respecting key recovery reaches $8$ of Lake Keyak's $12$ rounds at time
$2^{71}$ for $128$-bit keys, $9$ rounds at $2^{137}$ for $256$-bit keys, and
$9$ of KMAC256's $24$ rounds at $2^{147}$ \cite{SGLS18}. Lake Keyak is a
keyed duplex over the same $1600$-bit permutation whose initialization, like
\TW{}'s root and leaf initializations, absorbs the key and nonce ahead of
nonce-varied cube queries, so these results are the closest published analog
to \TW{}'s keyed setting. They leave a margin of $3$ to $4$ of $12$ rounds
against attacks whose complexities start at $2^{71}$.

In sum, the known attack frontier stands at $6$ of $12$ rounds for collisions
and preimages and at $8$ to $9$ rounds for keyed-mode key recovery and
structural distinguishers. \TW{} uses \Keccakp{} at the same round count and
with the same rate/capacity split as KangarooTwelve and TurboSHAKE128, so its
margin assessment tracks theirs, and any future advance against the
permutation applies to all three alike. What \TW{} adds at the mode level and
these analyses do not cover is the tree structure itself---in particular the
family of leaf initializations sharing $(K, U)$ and differing only in the
leaf index---and we consider dedicated cryptanalysis of that structure a
natural target for future work.

\subsection{Limitations and Open Problems}
\label{sec:limitations}

\Cref{sec:security-model} states the formal scope of the security claims; this
subsection discusses the consequences of the excluded settings and the
corresponding open problems.

\paragraph{Nonce misuse.}
As scoped in \Cref{sec:security-model}, \TW{}'s AE and mu-ind bounds require
nonce-respecting encryption queries, and the mode is not misuse-resistant: a
repeated nonce reuses leaf keystream and breaks confidentiality, as for any
online duplex stream cipher. Schemes such as AES-GCM-SIV \cite{GLL17}, in the
deterministic/misuse-resistant lineage of \cite{RS06}, tolerate nonce
repetition at the price of a second pass. A misuse-resistant variant of \TW{}
is left open.

\paragraph{Unverified plaintext release.}
The \textsf{INT-RUP} theorem of \Cref{thm:int-rup} covers integrity when an
adversary can obtain plaintext-computing oracle outputs before verification, in
the release-of-unverified-plaintext setting of \cite{ABLMMY14}. This is an
integrity-only statement. It does not give plaintext awareness or privacy under
release of unverified plaintext: \TW{} is online and single-pass, and exposing
unverified stream plaintext can reveal keystream information. Application
behavior on unverified plaintext is likewise outside the provable model.

\paragraph{Conservative lift.}
Per-term tightness is itemized in \Cref{sec:loss-terms}; the flat-sponge terms
are conservative because our proof follows the unkeyed sponge/hash line. This
route matches the role of the authentication tag: \TW{}'s committing statements
are consequences of tag hash security, and sponge hash security is supported by
a mature literature. The price is the indifferentiability lift of
\Cref{app:flat-sponge-lift}: one $\Lift_c(\Ntot)$ term on the real side, and in
real-vs-ideal games a second proximity term at $\Nprim$
(\Cref{cor:lift-application}).

A direct keyed-duplex treatment would have to close two gaps. First, existing
keyed-duplex bounds are hidden-key AE bounds: the game samples the key, and the
keyed initialization material is not part of the adversary's output. The
root tag claims here are adversarial-initialization statements. In the hash and
committing games, the adversary may choose the key, nonce, associated data,
message, or competing opening. Thus the proof must handle adversarially chosen
initialization transcripts, not only hidden-key encryption and verification
queries.

Second, the available keyed-duplex AE interface is not a direct substitute for a
byte-string AEAD with arbitrary final-block length. The modern bounds
synthesized by \cite{Men23} use a duplex call phased as permutation, squeeze,
then absorb or overwrite (\cite[Sec.~3.4, Table~1]{Men23}). In the
authenticated-encryption application, this interface is instantiated as
\textsf{MonkeySpongeWrap} \cite[Alg.~8]{Men23}. As written, however, the
overwrite interface does not reconstruct the missing padding bits of a truncated
final ciphertext block on decryption: encryption absorbs the padded final
plaintext block after squeezing and then truncates the resulting ciphertext to
the message length, while decryption receives only this truncated ciphertext,
pads it independently, and uses the overwrite branch to set the rate. The
decryption transcript can therefore diverge before tag recomputation.

\TW{} instead stays with the BDPV11 duplex phasing: a call absorbs the actual
input string, padded by the duplex rule, and then returns the requested output.
At the stream layer this supports arbitrary length messages by splitting them
into bounded duplex blocks, with the final short ciphertext block absorbed as it
is rather than reconstructed through an overwrite branch (\Cref{sec:tw128}). A
direct keyed-duplex theorem for \TW{} would therefore need both this BDPV11
phasing and adversarial-initialization hash and committing games. We retain the
indifferentiability lift because it covers these claims uniformly and already
meets the $128$-bit target
(\Cref{sec:bound-discussion}).

\paragraph{Beyond-birthday security and leakage.}
All bounds are birthday-type on the $256$-bit capacity, delivering the
$128$-bit target but no more; a beyond-birthday analysis is open. The model is
also single-stage and leakage-free: side-channel resistance is addressed only at
the implementation level, through the data-independent execution properties of
\Cref{sec:impl}, and is not part of the provable-security claims.

\section{Conclusion}
\label{sec:conclusion}

\TW{} starts from the premise that an AEAD tag should be the compact digest
applications often treat it as: a hash output for the full encryption context,
with committing security following from that hash story rather than added as an
external layer. The construction adapts the KangarooTwelve tree-hashing pattern
into a two-level tree of keyed duplexes over a single \Keccakp{} permutation.
The chunk decomposition is canonical, while the number of leaf duplexes
evaluated in parallel is an implementation choice rather than a mode parameter.
This gives scalar, SIMD, and multicore implementations the same ciphertexts and
tags, with constant working memory and scalable parallelism.

The main security result is tag hash security for the wrapper that returns the
final authentication tag on input $(K, U, A, M)$. In the public permutation
model, via an intermediate random oracle and the flat sponge lift, we prove
collision resistance, standard and everywhere second-preimage resistance,
preimage resistance on fixed input slices, and everywhere preimage resistance.
Those bounds yield the AEAD-facing committing statements, including full-input
\textsf{CMTDs-4}. The same root structure gives \textsf{INT-RUP}
integrity under release of unverified plaintext. As baseline AEAD properties, we
prove concrete multi-user indistinguishability and derive all-in-one AE as its
single-user corollary. At the implemented parameters, all stated notions meet a
$128$-bit target.

The implementation results show the cost and payoff of putting that tag hash
inside a tree. Short inputs pay for the larger permutation and the root/leaf
structure, but long inputs expose independent leaves while preserving one
canonical wire format. Vectorized \texttt{amd64} and \texttt{arm64}
implementations reach bulk throughputs of $48.0$\,Gbit/s on the AVX-512 path,
$39.4$\,Gbit/s on the \texttt{arm64} path, and $16.6$\,Gbit/s on the AVX2
path, with the AVX-512 and \texttt{arm64} paths at $51$--$76\%$ of
hardware-accelerated AES-128-GCM throughput.

The remaining questions are those left open in \Cref{sec:limitations}: nonce
misuse would require a different construction, privacy under release of
unverified plaintext is not provided by the \textsf{INT-RUP} theorem, and sharper
bounds would need either a direct keyed-duplex treatment of the BDPV11-style
phasing with adversarial initialization or a beyond-birthday argument. More
broadly, \TW{} supports the thesis that a single \Keccak{} permutation with
substantial cryptanalytic history can provide authenticated encryption whose tag
has hash security strong enough to drive commitment, while retaining scalable
hashing-style performance.


\section*{Acknowledgments}
A hearty thank-you to my wife, who remains reasonably mystified at my choice
of hobbies.

\bibliographystyle{alpha}
\bibliography{refs}

\appendix
\section{Supporting Lemmas}
\label{app:lemmas}

The proof core separates structural decoding from probabilistic
accounting. The well-formed transcript language of
\Cref{def:wf-transcripts} fixes which flattened queries are valid \TW{}
construction prefixes.
\Cref{lem:duplex-flat,lem:transcript-inj,lem:bridge-decode} then show
how such a query is flattened, parsed as a typed transcript, and
assigned an output point; the exact keyed decoder of
\Cref{def:keyed-twout} then decodes direct queries to candidate keys,
when any.
\Cref{lem:key-tests} turns those decoded keys into an adaptive key-test
bound. The remaining proof interfaces expose the events their consumers
must control: encryption-side freshness, verification hidden-key exposure
and root tag guessing, leaf tag collisions, and the committing
candidate-sequence collision event. Thus
the mu-ind and committing proofs only have to derive the relevant local
event premise before applying the corresponding accounting argument.

\paragraph{How to read this appendix.}
\Cref{fig:lemma-dag} records the proof dependency map. It groups local
interfaces when the exact proof citations would obscure the spine. The
Bridge \& Decode node is \Cref{lem:bridge-decode} together with the
exact keyed decoder of \Cref{def:keyed-twout}. The Root/Output
Separation node groups Output-Point Separation (\Cref{cor:output-sep}),
Opening Triple Separation (\Cref{cor:triple-sep}), and Final-Root Input
Separation (\Cref{prop:final-root-sep}); the use of Leaf-Transcript
Recovery (\Cref{lem:leaf-recovery}) is folded into the final-root
separation fact. The Key-Indexed Renaming node (\Cref{lem:key-renaming})
supplies the key-field disjointness and renaming used by hidden-key
encryption and verification accounting. The Adaptive Candidate node
(\Cref{lem:adaptive-candidate}) is the generic probabilistic engine: it
bounds the Leaf Tag Collision lemma (\Cref{lem:leaf-collision}) and,
through the Root Tag Candidate Sequence and Root Tag Hash and Target-Hit Bound
(\Cref{lem:root-candidate-seq,lem:root-tag-hash}) of \Cref{sec:cmtds3},
the root tag hash and committing proofs. The mu-ind proof reuses the
same structural lemmas through the explicit hybrid game sequence of
\Cref{def:adjacent-hybrid-games}, adding freshness, key-test,
leaf-collision, and verification-first root guessing accounting.

\begin{figure}[H]
  \centering
  \resizebox{\linewidth}{!}{%
  \begin{tikzpicture}[
    every node/.style={font=\footnotesize},
    lemma/.style={draw, rounded corners=2pt, align=center, inner sep=3pt,
                  minimum width=2.6cm, minimum height=0.7cm, fill=white},
    group/.style={draw, rounded corners=2pt, align=center, inner sep=3pt,
                  minimum width=2.8cm, minimum height=0.7cm, fill=white},
    consumer/.style={draw, dashed, rounded corners=2pt, align=center,
                     inner sep=3pt, minimum width=2.3cm, minimum height=0.7cm,
                     fill=white},
    arr/.style={-{Latex[length=1.6mm]}, semithick, draw=black!70},
  ]
    % Column 1: foundational lemmas.
    \node[lemma] (df) at (0,  2.6) {Duplex\\ Flattening};
    \node[lemma] (ba) at (0, -2.6) {Bridge \&\\ Decode};

    % Column 2: structural interfaces.
    \node[lemma] (ti) at (3.8,  2.6) {Transcript\\ Injectivity};
    \node[lemma] (kr) at (3.8,  0.8) {Key-Indexed\\ Renaming};
    \node[lemma] (kt) at (3.8, -1.0) {Adaptive\\ Key Tests};
    \node[group] (rs) at (3.8, -2.8) {Root/Output\\ Separation};

    % Column 3: local accounting lemmas.
    \node[lemma] (tf) at (7.6,  2.8) {Transcript\\ Freshness};
    \node[lemma] (lc) at (7.6,  1.2) {Leaf Tag\\ Collision};
    \node[lemma] (rg) at (7.6, -0.4) {Verification-First\\ Root Guessing};
    \node[lemma] (cc) at (7.6, -2.0) {Adaptive\\ Candidate};

    % Column 4: grouped proof spines.
    \node[group] (mi) at (11.4,  1.7) {Mu-ind\\ Accounting};
    \node[lemma] (rc) at (11.4, -0.5) {Root Tag\\ Candidate Seq.};
    \node[lemma] (rh) at (11.4, -2.0) {Root Tag\\ Hash/Target};

    % Column 5: consumers.
    \node[consumer] (c6) at (15.2,  1.7) {Mu-ind\\proof};
    \node[consumer] (c8) at (15.2, -1.3) {Hash +\\committing proofs};
    \node[consumer] (cb) at (15.2, -3.0) {App.~\ref{app:flat-sponge-lift}};

    \begin{scope}[on background layer]
      % Foundational -> structural interfaces.
      \draw[arr] (df) -- (ti);
      \draw[arr] (ti) -- (kr);
      \draw[arr] (ba.east) to[out=35, in=205] (kr.west);
      \draw[arr] (ba.east) to[out=10, in=200] (kt.west);
      \draw[arr] (ti.east) to[out=-35, in=145] (rs.west);
      \draw[arr] (ba) -- (rs);

      % Local accounting from the structural interfaces.
      \draw[arr] (ti.east) to[out=20, in=180] (tf.west);
      \draw[arr] (ti.east) to[out=0, in=180] (lc.west);
      \draw[arr] (ti.east) to[out=-20, in=170] (rg.west);
      \draw[arr] (cc) -- (lc);
      \draw[arr] (rs.east) to[out=25, in=190] (rc.west);
      \draw[arr] (ba.east) to[out=-5, in=200] (rc.west);
      \draw[arr] (rc) -- (rh);
      \draw[arr] (cc.east) to[out=-5, in=180] (rh.west);

      % Mu-ind accounting and consumer.
      \draw[arr] (tf.east) to[out=0, in=170] (mi.west);
      \draw[arr] (lc.east) to[out=0, in=185] (mi.west);
      \draw[arr] (rg.east) to[out=20, in=205] (mi.west);
      \draw[arr] (kr.east) to[out=25, in=195] (mi.west);
      \draw[arr] (kt.east) to[out=35, in=215] (mi.west);
      \draw[arr] (mi) -- (c6);

      % Root tag hash and committing consumers.
      \draw[arr] (lc.east) to[out=-20, in=160] (c8.west);
      \draw[arr] (rs.east) to[out=-5, in=185] (c8.west);
      \draw[arr] (rh.east) -- (c8.west);

      % The flat sponge lift consumes bridge decoding directly.
      \draw[arr] (ba.east) to[out=-20, in=180] (cb.west);
    \end{scope}
  \end{tikzpicture}
  }
  \caption{Proof dependency map for the supporting components in this
    appendix and their consumers. Solid boxes are lemmas or grouped local
    proof interfaces; dashed boxes are proof components that consume them.}
  \label{fig:lemma-dag}
\end{figure}

\subsection{Duplex Flattening}
\label{app:lem-duplex-flat}

\begin{lemma}[Duplex Flattening]
\label{lem:duplex-flat}
Every \TW{} root or leaf duplex output is equal to a \cite{BDPV11} sponge
output on the corresponding flattened duplex-call transcript, and the
ciphertext-absorbing message phase admits a flattened duplex transcript
that is well defined for both encryption and decryption.
\end{lemma}

\begin{proof}
Let $\mathsf{D}$ be a duplex object using permutation $\pi$, rate $r$,
and padding $\padtenone$. For a sequence of duplex calls $Z_i =
\mathsf{D}.\duplexing(\sigma_i, \ell_i)$ with each $\sigma_i$ short
enough to fit in one duplex block after padding, \cite[Lemma~3]{BDPV11}
gives
\[
  Z_i = \Sponge[\pi, \padtenone, r]
    \bigl(\sigma_0 \concat \mathsf{pad}_0 \concat \sigma_1
          \concat \mathsf{pad}_1 \concat \cdots \concat \sigma_i,\, \ell_i\bigr),
\]
where $\mathsf{pad}_j = \padtenone[r](|\sigma_j|)$, and Lemma~4 says
that, for fixed $\padtenone$ and $r$, the map from the duplex input-string
sequence to the flattened sponge input is injective.

In \TW{}, each duplex call absorbs a byte-aligned block $\sigma$
followed by a $4$-bit phase suffix $\sigph \in \{\siginit,
\sigadmore, \dots, \sigagglast\}$
(\Cref{sec:domain-sep}), so the duplex-call input is
$\sigma \concat \sigph$ followed by $\padtenone$. With
$r = 1344$ and $\rho_{\max}(\padtenone, r) = r - 2 = 1342$ bits, the
maximum byte-aligned input is $\rho = 167$ bytes
(\Cref{sec:params}), giving a worst-case duplex-call input of $8 \cdot
167 + 4 = 1340 \leq 1342$ bits. The two initialization calls are
shorter still: $|\prfxroot \concat K \concat U| + 4 = 8
\cdot (6 + 32 + 32) + 4 = 564$ bits and $|\prfxleaf \concat
K \concat U \concat \langle j \rangle_8| + 4 = 8 \cdot (6 + 32 + 32 +
8) + 4 = 628$ bits. \cite[Lemmas~3 and~4]{BDPV11} therefore apply
independently to every \TW{} duplex object.

The root duplex is initialized by absorbing $\prfxroot
\concat K \concat U$ under \siginit (\Cref{alg:enc}), then
absorbs the associated data blocks under \sigadmore /
\sigadlast, the root ciphertext blocks in the message phase under
\sigmsgmore / \sigmsglast, and the leaf tag aggregation
blocks under \sigaggmore / \sigagglast. By the
duplex-to-sponge equality, every root duplex output is the sponge
function evaluated on the root transcript prefix accumulated up to the
relevant output point. The same holds for each leaf duplex,
initialized by $\prfxleaf \concat K \concat U \concat
\langle j \rangle_8$ and absorbing only its own ciphertext blocks in the
message phase. In decryption and verification, \Cref{alg:dec} absorbs the
submitted ciphertext blocks before the root tag comparison, so the same
flattened message phase transcript is defined whether the comparison
accepts or rejects.
\end{proof}

\subsection{Well-formed Transcript Prefixes}
\label{app:wf-transcripts}

For a duplex-call prefix $((\sigma_0, \sigph[0]), \dots, (\sigma_i,
\sigph[i]))$, write
\[
  \Flat\bigl((\sigma_0, \sigph[0]), \dots, (\sigma_i, \sigph[i])\bigr)
  = \bigconcat_{j=0}^{i-1}\!
    \bigl(\sigma_j \concat \sigph[j] \concat \mathsf{pad}_j\bigr)
    \concat \sigma_i \concat \sigph[i],
\]
where $\mathsf{pad}_j = \padtenone[r](|\sigma_j \concat \sigph[j]|)$,
the padding of the $j$-th full duplex-call input.
The current call's padding $\mathsf{pad}_i$ is not part of $\Flat$;
it is applied internally by the sponge function. Equivalently, after
the $i$-th call the duplex state is the sponge state after absorbing
$\Flat(\cdot) \concat \mathsf{pad}_i$.

\begin{definition}[Well-Formed \TW{} Transcript Prefixes]
\label{def:wf-transcripts}
A typed \TW{} transcript prefix is \emph{well formed} if it is a prefix,
ending at a construction transcript lookup point, of one of the following
two call languages.
\begin{itemize}
\item A root transcript starts with
  $(\prfxroot \concat K \concat U,\siginit)$, with fixed-width fields
  $K \in \bits^k$ and $U \in \bits^\nu$. It then absorbs the
  associated data byte string (only when the associated data is
  nonempty; the phase is otherwise elided, \Cref{sec:enc}), the root
  ciphertext byte string, and the aggregation byte string (only when
  $n \geq 1$; when $n = 0$ the aggregation phase is elided,
  \Cref{sec:enc}, and the root tag is emitted by the final
  $\sigmsglast$ message call), in that order. Each present phase byte
  string is partitioned by the block rule
  of \Cref{alg:helpers}: the empty string gives one empty final block,
  and a nonempty string gives zero or more $\rho$-byte non-final blocks
  followed by one final block. The phases use respectively
  $\sigadmore/\sigadlast$, $\sigmsgmore/\sigmsglast$, and
  $\sigaggmore/\sigagglast$.
\item A leaf-$j$ transcript starts with
  $(\prfxleaf \concat K \concat U \concat \langle j \rangle_8,\siginit)$,
  where $1 \leq j < 2^{64}$ and the fields have fixed width. It then
  absorbs the leaf ciphertext byte string under
  $\sigmsgmore/\sigmsglast$, using the same block rule.
\end{itemize}
The message byte string in this language is always the ciphertext byte
string absorbed by the construction: encryption obtains it by XORing the
plaintext with the current keystream, while decryption and verification
use the submitted ciphertext. Thus a verification transcript is well
formed for every in-domain ciphertext before the tag comparison is made;
acceptance is not part of the transcript language.

The root aggregation byte string has the form
$T_1 \concat \cdots \concat T_n \concat \langle n \rangle_8$, where each
$T_j$ has fixed byte length $\tauleaf/8$. The trailing fixed-width field
therefore gives a unique parse of the leaf tag sequence. The
output-point class of a well-formed counted prefix is the class in
\Cref{tab:keyed-twout} determined by its duplex instance and final-call
suffix.
\end{definition}

\begin{remark}[Well-Formedness versus Realizability]
\label{rem:wf-vs-realizable}
The language of \Cref{def:wf-transcripts} is deliberately larger than
the set of transcript prefixes that \TW{} computations realize. For
example, it does not require an aggregation phase to follow a root
message phase of exactly $\beta$ bytes, although every real computation
with $n \geq 1$ absorbs a full $\beta$-byte root chunk (\Cref{alg:enc}).
The inclusion runs in the safe direction everywhere: every
construction-generated prefix is well formed, so the structural lemmas
of this appendix and the separation propositions of
\Cref{sec:root-tag-hash-security} apply to construction transcripts as
restrictions of statements proved on the larger language, while the
direct-query decoder of \Cref{def:keyed-twout} and the root tag
candidate predicate of \Cref{lem:root-candidate-seq} accept the larger
language and therefore only over-count adversary queries, which loosens
no bound.
\end{remark}

Name the tag-emitting transcript families:
\begin{align*}
  \Rtag   &= \text{root transcript prefix whose sponge output gives the final root tag}, \\
  \Ltag(j)   &= \text{leaf-$j$ transcript prefix whose sponge output gives the leaf tag}.
\end{align*}
The final root tag is emitted by the $\sigagglast$ aggregation call (the
$\Rtag$ family) when $n \geq 1$, and by the closing $\sigmsglast$
message call (an $\Rmsg^+$ prefix, in the notation of
\Cref{cor:output-sep}) when $n = 0$ and the aggregation phase is elided.
Since $n = \leaves(C)$ is fixed by $\|C\|$, every computation on a given
ciphertext emits its tag at the output point selected by $\|C\|$.
Intermediate construction transcript lookup points (init, non-final
AD, non-final message, and non-final aggregation) are also typed
prefixes covered by \Cref{lem:bridge-decode}; the exact keyed
direct-query decoder names them in \Cref{tab:keyed-twout}.

\begin{definition}[Root Tag Output Point and Final Root Transcript Class]
\label{def:root-tag-output-point}
For any in-domain encryption or verification computation on ciphertext
$C$, the \emph{root tag output point} is the construction lookup that
emits the root tag: $\Rtag$ when $\leaves(C) \geq 1$, and the closing
root message point $\Rmsg^+$ when $\leaves(C) = 0$ and aggregation is
elided. The \emph{final root transcript class} of such a computation is
the set of construction computations that realize the same bridged final
root input at this root tag output point. A class is
\emph{encryption-first} if its first construction generation is by an
encryption query, and \emph{verification-first} if its first construction
generation is by a non-replay verification query.
\end{definition}

\subsection{Transcript Injectivity}
\label{app:lem-transcript-inj}

\begin{lemma}[Transcript Injectivity]
\label{lem:transcript-inj}
Well-formed \TW{} flattened duplex-call transcript prefixes are
injectively encoded and separated by duplex instance (root vs leaf),
leaf index $j$, phase suffix, and duplex-call order. For final-root
prefixes, equality of flattened inputs recovers the same root
initialization, AD blocks, root ciphertext blocks, aggregation input,
and leaf tag sequence.
\end{lemma}

\begin{proof}
By \Cref{lem:duplex-flat}, every well-formed output point can be written
as a \cite{BDPV11} flattened sponge input for the corresponding
duplex-call prefix, namely $\Flat$ as defined above. The deterministic
current-call $\padtenone$ is applied internally by the sponge function
and is not part of the input to which \cite[Lemma~4]{BDPV11} is applied.
That lemma recovers the duplex-call sequence from the flattened sponge
input; in \TW{}, each call input is
$\sigma \concat \sigph$ with byte-aligned $\sigma$ and
fixed $4$-bit suffix $\sigph$, so recovering the call input recovers the pair
$(\sigma, \sigph)$.

\Cref{def:wf-transcripts} fixes the root and leaf initialization fields,
the admissible phase order, and the final-call suffixes. The distinct
prefixes $\prfxroot$ and $\prfxleaf$ separate root from leaf transcripts;
the fixed-width leaf field $\langle j \rangle_8$ separates leaves; and
the seven suffix values of \Cref{tab:suffixes} separate phases and
non-final/final calls. Thus equality of flattened inputs forces equality
of the recovered instance type, key, nonce, leaf index when present,
phase suffixes, and duplex-call order.

\paragraph{Phase-recovery sub-claim.}
For each present phase $P \in \{\mathsf{ad}, \mathsf{msg}, \mathsf{agg}\}$
(the $\mathsf{ad}$ phase is present exactly when the associated data is
nonempty; see \Cref{sec:enc}),
\TW{} absorbs the phase's input byte string $s_P$ by partitioning into
$\rho$-byte blocks (with the empty case yielding a single empty block,
per \Cref{sec:enc}) and feeding each block to the duplex with suffix
$\sigma_P^-$ (non-final) or $\sigma_P^+$ (final).
\Cref{lem:duplex-flat} and \cite[Lemma~4]{BDPV11} recover the
corresponding $(\sigma, \sigph)$ sequence from the flattened input.
Because the partition map is deterministic and uniquely invertible on
its image (the empty case yielding the single empty block carrying
$\sigma_P^+$), concatenating the recovered $\sigma$ blocks
reconstructs $s_P$ exactly. The map from $s_P$ to the phase's $(\sigma,
\sigph)$ sequence is therefore injective.

Applied phase by phase to a well-formed final-root transcript:
\begin{itemize}
\item AD. The associated data phase contributes a $\sigadlast$-tagged
  block exactly when the associated data is nonempty; an empty
  associated data elides the phase and contributes no block. When the
  phase is present the sub-claim recovers the AD byte string. Thus
  distinct associated data strings $A_i \neq A_j$ are separated either
  by the presence versus absence of the AD phase (when exactly one is
  empty) or, when both are nonempty, by their distinct AD
  $\sigma$-block sequences.
\item MSG. Distinct absorbed root or leaf ciphertext byte strings
  produce distinct MSG $\sigma$-block sequences; the message language of
  \Cref{def:wf-transcripts} is ciphertext-absorbing for encryption and
  decryption alike.
\item AGG. The aggregation phase is present exactly when $n \geq 1$
  (\Cref{sec:enc}). When it is present, the sub-claim applied to the
  aggregation byte string $T_1 \concat \cdots \concat T_n \concat
  \langle n \rangle_8$ recovers that byte string. The final eight bytes
  give $\langle n \rangle_8$, and each leaf tag has fixed byte length
  $\tauleaf / 8$, so the byte string has a unique parse as a leaf tag
  sequence. Two root aggregation encodings are therefore equal exactly
  when their leaf tag sequences are equal. When
  $n = 0$ the aggregation phase is absent and the final-root transcript
  ends at the recovered $\sigmsglast$ root message call; the MSG bullet
  recovers the root ciphertext byte string, and the absence of any
  $\sigaggmore/\sigagglast$ call distinguishes such a prefix from every
  $n \geq 1$ final-root prefix.
\end{itemize}

Therefore the recovered final-root transcript determines the root
initialization, AD blocks, root ciphertext blocks, and, when
$n \geq 1$, the aggregation byte string and leaf tag sequence.
\end{proof}

\begin{lemma}[Leaf-Transcript Recovery]
\label{lem:leaf-recovery}
Let two construction-generated final-root prefixes have equal flattened
inputs, and suppose that, for each position $j$ of their common leaf tag
sequence, the two computations' position-$j$ final leaf transcripts are
either equal or receive distinct leaf tags. Then equality of flattened
inputs recovers the same leaf transcript sequence, and therefore the
same full absorbed ciphertext. Both computations' position-$j$ final
leaf transcripts are leaf-$j$ transcripts under the common recovered
$(K, U)$, so the hypothesis holds whenever no two distinct
construction-generated final leaf transcripts with the same key, nonce,
and leaf index receive the same leaf tag; when one of the two prefixes
is generated by an encryption computation, it suffices that no
construction-generated final leaf transcript receives the same leaf tag
as a distinct encryption-generated final leaf transcript with the same
key, nonce, and leaf index.
\end{lemma}

\begin{proof}
By Transcript Injectivity (\Cref{lem:transcript-inj}), the two equal
flattened final-root inputs recover the same root initialization, AD
blocks, root ciphertext blocks, and, when $n \geq 1$, the same
aggregation byte string and leaf tag sequence; the trailing
$\langle n \rangle_8$ field fixes the common leaf count $n$. When
$n = 0$ no leaf is present and the recovered root ciphertext blocks
already constitute the full absorbed ciphertext. When $n \geq 1$, the
aggregation byte string parses uniquely as the leaf tag sequence
$T_1 \concat \cdots \concat T_n$; the two leaf tag sequences are equal,
so for each $j$ the two computations' position-$j$ final leaf
transcripts---leaf-$j$ transcripts carrying the common recovered
$(K, U)$ (\Cref{def:wf-transcripts}, \Cref{sec:enc})---receive the same
leaf tag, and by hypothesis are equal, hence so are
their absorbed leaf ciphertext bodies. Together with the common root
ciphertext blocks this recovers the same full absorbed ciphertext.
\end{proof}

\subsection{Bridge and Typed Decoding}
\label{app:lem-bridge-decode}

\begin{lemma}[Bridge Injectivity]
\label{lem:bridge-decode}
The random oracle games used in
\Cref{sec:aead-security,sec:root-tag-hash-security} and \Cref{cor:committing-combined} are
ordinary random oracle games on a single random oracle
$\ROflat : \bits^* \to \bits^\infty$ over the unpadded
\cite{BDPV08} $H$-input domain (\Cref{sec:ro-model}), with \TW{} root and
leaf duplex outputs answered by the typed restriction
\[
  \RF(X) = \ROflat(\flatmap(X)),
  \qquad \flatmap(X) = \bridge[r, r_{\max}]\bigl(\mathsf{raw\_flat}(X)\bigr),
\]
for $r = r_{\max} = 1344$. Here $\mathsf{raw\_flat}(X)$ is the native
flattened input $\Flat(\cdot)$ of \Cref{app:wf-transcripts} applied to
the duplex-call sequence of $X$, and $\bridge[r, r_{\max}]$ is the
\cite{BDPV11} Theorem~3 preprocessing map. The induced map
$X \mapsto \flatmap(X)$ is injective on well-formed typed \TW{} duplex
transcript prefixes.
\end{lemma}

\begin{proof}
The bridge $\bridge[r, r_{\max}]$ is described by the proof of
\cite[Theorem~3]{BDPV11} in three steps: (1)~pad the input with
$\padtenone[r]$; (2)~for each $r$-bit block, append $r_{\max} - r$
zero bits; (3)~strip the trailing $\padten$ from the result. The
$\padtenone$ padding in step~(1) appends a~$1$, then
$q = (-|\mathsf{raw\_flat}(X)| - 2) \bmod r$ zero bits, then a closing
$1$; the $\padten$-unpadding in step~(3) strips the $r_{\max} - r$
trailing zero bits introduced by step~(2) together with the closing
$1$ of step~(1). For \TW{} with $r = r_{\max}$, step~(2) appends an
empty string and step~(3) strips zero trailing zeros and the closing
$1$, so the net effect of steps~(1) and~(3) is to map a native flattened
input $W = \mathsf{raw\_flat}(X)$ to $W \concat 1 \concat 0^q$. Any
string in this image has a unique inverse: $q$ is the number of trailing
zero bits, the preceding bit is the inserted $1$, and deleting this
suffix recovers $W$.

Injectivity of $\flatmap$ on well-formed typed \TW{} prefixes then
follows from \Cref{lem:transcript-inj}: distinct well-formed typed
prefixes have distinct
$\mathsf{raw\_flat}(\cdot)$ images, and the bridge is injective on
that image. Direct adversary queries on inputs outside the bridged
\TW{} well-formed transcript language still count in $\qRO(\mathcal{B})$
(\Cref{sec:resources}) but cannot trigger
$\badkey$.
\end{proof}

The bridge injectivity established here underlies Opening Triple
Separation (\Cref{cor:triple-sep}, \Cref{sec:cmtds3}): distinct
$(K, U, A)$ triples yield distinct bridged final root transcripts.

\begin{corollary}[Output-Point Separation]
\label{cor:output-sep}
Every well-formed \TW{} typed transcript prefix at a counted
construction transcript lookup point belongs to exactly one of the
following output-point classes, identified by the recovered duplex
instance and the $4$-bit suffix of its final call:
\begin{itemize}
\item \Rinit (root duplex, suffix \siginit),
\item \Linit($j$) (leaf-$j$ duplex, suffix \siginit),
\item $\Rad^-$ / $\Rad^+$ (root duplex, suffix \sigadmore or \sigadlast),
\item $\Rmsg^-$ / $\Rmsg^+$ (root duplex, suffix \sigmsgmore or \sigmsglast),
\item $\Lmsg^-(j)$ (leaf-$j$ duplex, suffix \sigmsgmore),
\item $\Ltag(j)$ (leaf-$j$ duplex, suffix \sigmsglast),
\item \Ragg (suffix \sigaggmore),
\item $\Rtag$ (root duplex, suffix \sigagglast).
\end{itemize}
For any two such prefixes whose bridged flattened inputs are equal, the
separation of \Cref{lem:transcript-inj} recovers the same duplex
instance, leaf index $j$ when present, $(K, U)$ pair, output-point
class, and duplex-call sequence. Equivalently, any two distinct counted
transcript prefixes already differ in one of these recovered fields, and
hence in their flattened inputs.
\end{corollary}

\begin{proof}
\Cref{lem:bridge-decode} recovers equality of native flattened inputs
from equality of bridged inputs, after which \Cref{lem:transcript-inj}
recovers the duplex-call sequence and its per-call
$(\sigma, \sigph)$ pairs. The first call's $\sigma$ is
$\prfxroot \concat K \concat U$ for root instances and
$\prfxleaf \concat K \concat U \concat \langle j \rangle_8$ for leaf
instances, with fixed-width fields, so it fixes the instance type (by
$\prfxroot$ vs $\prfxleaf$), the $(K, U)$ pair, and the leaf index $j$;
within each instance type the phase suffixes are pairwise distinct, so the
recovered instance type together with the final-call suffix fixes the
output-point class, and each prefix therefore lies in exactly one class.
Conversely, two distinct well-formed prefixes in the same class on the
same instance differ in some recovered call's $(\sigma, \sigph)$ pair
or in the
sequence length, hence in their flattened inputs.
\end{proof}

\begin{definition}[Exact Keyed Decoder]
\label{def:keyed-twout}
For a direct query to $\ROflat$ with input $Y$ and finite length
$\ell$, define $\KeyedTWOut(Y)$ as a partial map from $\bits^*$ to
$\bits^k \times \mathcal{C}_{\mathsf{out}}$. The map ignores
$\ell$. It returns $(K,\mathcal{C})$ iff
there exists a typed \TW{} duplex-call prefix $X$ such that
$Y = \flatmap(X)$, the raw flattened prefix of $X$ parses as a
well-formed \TW{} transcript prefix in the class $\mathcal{C}$ listed in
\Cref{tab:keyed-twout}, and the first call of $X$ contains the
candidate key $K$ in the fixed-width root field
$\prfxroot \concat K \concat U$ or leaf field
$\prfxleaf \concat K \concat U \concat \langle j \rangle_8$. If no
such prefix exists, $\KeyedTWOut(Y) = \bot$. In particular, inputs with
the right final suffix but a malformed phase order, noncanonical block
partition, malformed aggregation string, or out-of-range leaf index are
outside the well-formed language and decode to $\bot$. A non-$\bot$
decoding always fixes the candidate key $K$ and nonce $U$ in the
fixed-width initialization field. It fixes the associated data $A$ only
for root prefixes whose recovered phase order determines the complete AD
string: a final AD prefix, or any later root message or aggregation
prefix, with absence of AD calls determining $A=\varepsilon$. We use the
key projection $(K,\mathcal{C})$ throughout. When a game hides a single
context component $Z \in \{K,U,A\}$, write $\KeyedTWOut_Z(Y)$ for the
decoded value of that component when it is determined, and $\bot$
otherwise.
\end{definition}

By \Cref{lem:bridge-decode,cor:output-sep}, $\KeyedTWOut$ is well
defined: a bridged input cannot parse as two distinct counted
well-formed prefixes, and an accepted input fixes a unique candidate key
and output-point class. Inputs outside this well-formed transcript
language decode to $\bot$.

\begin{table}[ht]
  \centering
  \small
  \begin{tabular}{l l l}
    \toprule
    Class $\mathcal{C}$ & Instance & Final call / role \\
    \midrule
    \Rinit       & root     & \siginit, root initialization \\
    $\Rad^-$     & root     & \sigadmore, non-final AD block \\
    $\Rad^+$     & root     & \sigadlast, final AD block \\
    $\Rmsg^-$    & root     & \sigmsgmore, non-final root message block \\
    $\Rmsg^+$    & root     & \sigmsglast, final root message block \\
    \Ragg        & root     & \sigaggmore, non-final aggregation block \\
    $\Rtag$      & root     & \sigagglast, final aggregation / root tag \\
    \Linit($j$)  & leaf $j$ & \siginit, leaf initialization \\
    $\Lmsg^-(j)$ & leaf $j$ & \sigmsgmore, non-final leaf message block \\
    $\Ltag(j)$   & leaf $j$ & \sigmsglast, final leaf message block / leaf tag \\
    \bottomrule
  \end{tabular}
  \caption{Counted keyed transcript classes for $\KeyedTWOut$.}
  \label{tab:keyed-twout}
\end{table}

The table is a transcript-lookup table, not an observation-length
table. A class is counted even when the construction's requested
output length at that call is zero, including root initialization,
non-final AD and aggregation calls, a final AD call whose following
root chunk requests no keystream bytes, and the final root message
call. A direct query triggers the same decoder regardless of $\ell$:
short finite-output queries and zero-length queries $(Y,0)$ still
count as direct queries to the address $Y$. Leaves have no AD or
aggregation classes; a leaf transcript enters the table only through
leaf initialization, non-final leaf message blocks, or the final leaf
tag point. When the associated data is empty the root transcript
likewise has no $\Rad^-$ or $\Rad^+$ class (\Cref{sec:enc}); the
\Rinit class then supplies the first root keystream block, exactly as
\Linit($j$) supplies the first leaf keystream block, and the decoder
still ignores $\ell$. Symmetrically, when $n = 0$ the root transcript
has no \Ragg or $\Rtag$ class (\Cref{sec:enc}); the final root message
call $\Rmsg^+$ then supplies the root tag, exactly as $\Ltag(j)$
supplies a leaf tag, and the decoder again ignores $\ell$. An $n \geq 1$
computation also reaches an $\Rmsg^+$ point when closing chunk $0$, but
there the construction requests zero output and continues into the
aggregation phase, so the $\Rmsg^+$ value is exposed as a tag only on
the $n = 0$ path; since $n = \leaves(C)$ is fixed by $\|C\|$, the two
uses never coincide on a single ciphertext. We call \Rinit, $\Rad^+$,
$\Rmsg^-$, \Linit($j$), and $\Lmsg^-(j)$ the \emph{stream output
points}: by the schedule just described, these are the only classes at
which a construction computation requests keystream bytes, while tag
values are exposed only at $\Ltag(j)$ and $\Rtag$ points and, on the
$n = 0$ path, at $\Rmsg^+$.

The hidden-key encryption, verification, and mu-ind accounting all
reindex transcript addresses by the key field. We record that reindexing
once.

\begin{lemma}[Key-Indexed Transcript Renaming]
\label{lem:key-renaming}
For a key $K \in \bits^k$, let $\mathcal{Y}_K$ be the set of bridged
inputs $\flatmap(X)$ of well-formed \TW{} transcript prefixes $X$ whose
first-call key field is $K$---the counted hidden-key transcript addresses
of \Cref{tab:keyed-twout} for key $K$. Then:
\begin{enumerate}
\item \emph{Disjointness.} For distinct keys $K_0 \neq K_1$, the sets
  $\mathcal{Y}_{K_0}$ and $\mathcal{Y}_{K_1}$ are disjoint.
\item \emph{Renaming.} The map $\Phi_{K_0 \to K_1}$ that replaces the
  first-call key field $K_0$ by $K_1$ and leaves every other recovered
  field unchanged---nonce, leaf index, phase sequence, AD blocks,
  ciphertext blocks, aggregated leaf tag bytes, and root/leaf
  instance---induces through $\flatmap$ a bijection
  $\flatmap \circ \Phi_{K_0 \to K_1} \circ \flatmap^{-1} :
  \mathcal{Y}_{K_0} \to \mathcal{Y}_{K_1}$ preserving output-point
  classes and equality relations among addresses.
\item \emph{Direct-query avoidance.} A direct query input
  $Y \in \mathcal{Y}_K$ has $\KeyedTWOut(Y) = (K,\mathcal{C})$ for a
  counted class $\mathcal{C}$; hence if no prior direct query decodes to
  $K$ then $\mathcal{Y}_K$ is disjoint from all prior direct-query
  inputs.
\end{enumerate}
\end{lemma}

\begin{proof}
By bridge injectivity (\Cref{lem:bridge-decode}) and Transcript
Injectivity (\Cref{lem:transcript-inj}), a bridged input $\flatmap(X)$
determines the native flattened prefix and hence the first duplex call
$\prfxroot \concat K \concat U$ or
$\prfxleaf \concat K \concat U \concat \langle j \rangle_8$, in
particular its fixed-width key field $K$. \emph{(1)} If
$\flatmap(X) = \flatmap(X')$ with key fields $K_0$ and $K_1$ then
$K_0 = K_1$, so distinct keys give disjoint address sets. \emph{(2)}
$\Phi_{K_0 \to K_1}$ rewrites only the key field, leaving the recovered
duplex-call sequence otherwise intact, so it is a bijection on
well-formed prefixes preserving the output-point class
(\Cref{cor:output-sep}) and the equality pattern among prefixes;
$\flatmap$ transports it to the stated bijection on addresses. \emph{(3)}
By \Cref{lem:bridge-decode} a $Y = \flatmap(X) \in \mathcal{Y}_K$ decodes
under $\KeyedTWOut$ to $(K,\mathcal{C})$; the contrapositive gives the
avoidance claim.
\end{proof}

\subsection{Adaptive Candidate Bounds}
\label{app:lem-candidate-collisions}

The Adaptive Candidate Collision and Preimage lemma is stated as
\Cref{lem:adaptive-candidate} in \Cref{sec:root-tag-hash-security},
where a proof sketch also appears; this subsection supplies the full
proof.

\begin{proof}[Proof of \Cref{lem:adaptive-candidate}]
Both parts expose the interaction in chronological order and sample each
candidate's $\ROflat$ value lazily at its append step, which the premises
license: when $Y_j$ is appended its string is fixed by the prior state
(predictability) and no bit of $\ROflat(Y_j)$ has yet been
revealed (unrevealed value), so $\ROflat(Y_j)$ may be drawn then, uniform
and independent of everything so far.

\emph{Collision.} For a fixed position set $I = \{i_1 < \cdots < i_s\}
\subseteq \{1,\dots,L\}$, let $E_I$ be the event that these candidate
positions exist and their $\tau$-bit projections are all equal. Condition
on any positive-probability prefix immediately before position $i_r$ is
appended, for $r \geq 2$, and on the values sampled at the earlier
positions in $I$. By the above, $\ROflat(Y_{i_r})$ is then uniform and
independent of the conditioning, so its $\tau$-bit projection equals the
already fixed projection of $Y_{i_1}$ with probability $2^{-\tau}$.
Multiplying over $r = 2,\dots,s$ gives $\Pr[E_I] \leq 2^{-\tau(s-1)}$; if
some position in $I$ is never appended then $E_I$ is false and the bound
still holds. A union bound over the $\binom{L}{s}$ position sets gives
the claim.

\emph{Preimage.} For a fixed position $j \in \{1,\dots,L\}$, let $E_j$
be the event that position $j$ exists and
$\bigl\lfloor \ROflat(Y_j) \bigr\rfloor_\tau = t$. Condition on any
positive-probability prefix immediately before $Y_j$ is appended; the
target $t$ is a constant. By the above, $\ROflat(Y_j)$ is uniform and
independent of the conditioning, so its $\tau$-bit projection equals $t$
with probability $2^{-\tau}$. If position $j$ is never appended, $E_j$ is
false. A union bound over the $L$ positions gives the claim.

\emph{Second-preimage corollary.} If the designated position $j_0$ is
not appended, the event is false. Otherwise split the other candidates
according to whether they occur before or after $j_0$. For earlier
candidates, condition on any positive-probability prefix immediately
before $Y_{j_0}$ is appended, including the already sampled projections
of candidates $1,\dots,j_0-1$. The value $\ROflat(Y_{j_0})$ is then
fresh and uniform, so its $\tau$-bit projection hits one of those at most
$j_0-1$ earlier projections with probability at most
$(j_0-1)2^{-\tau}$. For a later candidate $Y_j$, $j>j_0$, condition on
the full prefix immediately before it is appended, including the fixed
projection $t_0 = \lfloor\ROflat(Y_{j_0})\rfloor_\tau$. The value
$\ROflat(Y_j)$ is fresh and uniform at its append step, so it hits
$t_0$ with probability $2^{-\tau}$. A union bound over the at most
$L-j_0$ later candidates gives an additional $(L-j_0)2^{-\tau}$.
Adding the two sides gives
$\Pr[\mathsf{sec}_{\tau,j_0}] \leq (L-1) \cdot 2^{-\tau}$; the
designated candidate is excluded, so its trivially equal projection does
not enter the count.
\end{proof}

\subsection{Adaptive Key Tests}
\label{app:lem-key-tests}

\begin{lemma}[Adaptive Key-Test Bound]
\label{lem:key-tests}
Consider a \TW{} random oracle model game with a single hidden context
component $Z$ sampled uniformly from $\bits^w$ for some $w \geq k$---the
standard case being the hidden key, $Z = K$ and $w = k$, used in the
AE-style proofs---direct adversary access to
$\ROflat$, and an immediate abort event $\badkey$ that fires on a
direct query $(Y,\ell)$ exactly when $Y$ decodes to some counted
output-point class $\mathcal{C}$ of \Cref{tab:keyed-twout} with
$\KeyedTWOut_Z(Y) = Z$. Suppose that, before the
abort, the following invariant holds for every direct query position
$i$: after conditioning on the full visible execution prefix and on
the set $S_i$ of distinct decoded candidate $Z$-values from earlier direct
queries, the hidden $Z$ is uniform on $\bits^w \setminus S_i$. Then
\[
  \Pr[\badkey]
  \leq \qRO / 2^w
  \leq \KeyGuess_k(\qRO)
  = \qRO / 2^k,
\]
where $\qRO$ is the execution-uniform count of direct $\ROflat$ queries.
\end{lemma}

\begin{proof}
If $\qRO \geq 2^w$ the right-hand side $\qRO/2^w$ is at least one, so the
claim is trivial. Assume $\qRO < 2^w$.

Order the direct $\ROflat$ queries. By the definition of
$\KeyedTWOut_Z$ and \Cref{lem:bridge-decode}, each query input either
decodes to $\bot$ or decodes to a unique candidate $Z$-value and counted
output-point class. Queries that decode to $\bot$, and repeated
queries whose decoded $Z$-value already lies in the current set $S_i$, do
not change the set of possible first successful tests. Removing
them can only shorten the sequence, so it suffices to analyze the
subsequence of distinct fresh decoded candidate $Z$-values
$Z'_1, \dots, Z'_m$, where $m \leq \qRO$.

Let $F_i$ be the event that the first $i$ fresh decoded candidate values
all miss the hidden $Z$. The invariant says that, conditioned on the
visible prefix before the $(i+1)$-st fresh test and on $F_i$, the
hidden $Z$ is uniform over the $2^w - i$ values not yet tested. The next
fresh candidate value $Z'_{i+1}$ is one of those remaining values and is
determined by the adversary's next direct query, so
\[
  \Pr[Z = Z'_{i+1} \mid F_i] = \frac{1}{2^w - i},
  \qquad
  \Pr[F_{i+1} \mid F_i] = 1 - \frac{1}{2^w - i}.
\]
Multiplying the miss probabilities gives
\[
  \Pr[F_m]
  = \prod_{i=0}^{m-1} \left(1 - \frac{1}{2^w - i}\right)
  = \prod_{i=0}^{m-1} \frac{2^w - i - 1}{2^w - i}
  = \frac{2^w - m}{2^w}.
\]
Thus the probability that some fresh decoded candidate value equals the
hidden $Z$ is $m / 2^w \leq \qRO / 2^w \leq \qRO/2^k$, the last step by
$w \geq k$. This event is exactly the first $\badkey$ abort, because
$\KeyedTWOut_Z$ returns either $\bot$ or a unique candidate $Z$-value.
\end{proof}

\subsection{Leaf Tag Collision Bound}
\label{app:lem-leaf-collision}

\begin{lemma}[Leaf Tag Collision Bound]
\label{lem:leaf-collision}
Consider any \TW{} random oracle model game with an adversary
$\mathcal{B}$ that has direct access to $\ROflat$ and construction
computations that may generate final leaf transcripts---by encryption,
plaintext-computing, or verification computations, whether they accept or
reject. Let
$\badleaf$ be the event that two distinct construction-generated final
leaf transcript prefixes $X \neq X'$ in leaf tag output-point classes
receive the same $\tauleaf$-bit projection from $\RF(\cdot)$, and let
$\nleaf(\mathcal{B})$ be the execution-uniform count of distinct
construction-generated final leaf transcripts (\Cref{sec:resources}).
The first bound applies the collision case ($s = 2$, $\tau = \tauleaf$)
of \Cref{lem:adaptive-candidate} to a chronological leaf tag candidate
sequence; the second is a same-slot freshness argument in an all-reject
game.
\begin{enumerate}
\item \emph{General count.} Unconditionally---in particular when some
  or all keys are adversary-supplied, so direct queries may themselves
  land on leaf tag output points---
  \[
    \Pr[\badleaf] \;\leq\;
    \LeafColl_{\tauleaf}\bigl(\qRO(\mathcal{B})+\nleaf(\mathcal{B})\bigr)
    = \binom{\qRO(\mathcal{B})+\nleaf(\mathcal{B})}{2} \big/ 2^{\tauleaf}.
  \]
\item \emph{Hidden key, same slot.} Suppose the game has a hidden key,
  nonce-respecting encryption queries, and \emph{all-reject}
  verification: every non-replay verification submission performs its
  construction computation and returns reject to the adversary. (The
  game \Grej{} of \Cref{def:adjacent-hybrid-games} and the
  \textsf{INT-RUP} continuation in the proof of
  \Cref{lem:int-rup-root-guess} have this shape.) Say two final
  leaf transcripts \emph{share a slot} if they recover the same key,
  nonce, and leaf index $(K, U, j)$ (\Cref{cor:output-sep}). Let
  $\badleaf$ be the event that some construction-generated final leaf
  transcript receives the same $\tauleaf$-bit projection as a distinct
  final leaf transcript that shares its slot and is generated by an
  encryption computation, and let $\mathsf{hit}_{\mathsf{leaf}}$ be the
  event that a direct adversary query hits the bridged input of a
  construction-generated hidden-key final leaf transcript before that
  transcript is first created. Then, in such a game,
  \[
    \Pr[\badleaf \text{ before } \mathsf{hit}_{\mathsf{leaf}}]
    \;\leq\; \frac{\nleaf(\mathcal{B})}{2^{\tauleaf}}.
  \]
\end{enumerate}
\end{lemma}

\begin{proof}
Build a chronological leaf tag candidate sequence: a candidate is the
bridged input of a leaf tag output point ($\Ltag(j)$), appended the first
time it is reached---by a direct query in $\mathcal{B}$'s oracle phase or
by a construction computation generating a final leaf transcript---and
not appended again. By \Cref{lem:bridge-decode}, \Cref{cor:output-sep},
and \Cref{lem:transcript-inj}, two reaches of the same bridged input come
from the same final leaf transcript, so repeated evaluations add no
candidate---the lazy $\ROflat$ value is reused---and distinct final leaf
transcripts give distinct candidates. No candidate's $\ROflat$ value is
read before its append step, so the sequence meets the predictability and
unrevealed-value premises of \Cref{lem:adaptive-candidate}, whose
collision case with $s = 2$ and $\tau = \tauleaf$ bounds the probability
that two candidates share a $\tauleaf$-bit projection by
$\binom{L}{2} \cdot 2^{-\tauleaf}$ for a length-$L$ sequence. Every
$\badleaf$ pair is a pair of construction-generated candidates.

\emph{(1) General count.} Direct queries may decode to a leaf
tag class $\Ltag(j)$ and are appended as candidates; together with the
$\nleaf(\mathcal{B})$ construction-generated transcripts the sequence has
length at most $\qRO(\mathcal{B}) + \nleaf(\mathcal{B})$, so the case
with $L = \qRO(\mathcal{B}) + \nleaf(\mathcal{B})$ gives
$\Pr[\badleaf] \leq
\binom{\qRO(\mathcal{B})+\nleaf(\mathcal{B})}{2} \cdot 2^{-\tauleaf}
= \LeafColl_{\tauleaf}(\qRO(\mathcal{B})+\nleaf(\mathcal{B}))$. The
count charges every direct query and every construction-generated final
leaf transcript regardless of key provenance, so no assumption on the
keys is used.

\emph{(2) Hidden key, same slot.} Every $\badleaf$ pair is a pair of
construction-generated hidden-key transcripts. Nonce-respecting
encryption uses each nonce in at most one encryption query, and within
one encryption computation the leaf index separates final leaf
transcripts, so each slot contains at most one encryption-generated
final leaf transcript---which may be a transcript first created by an
earlier verification or plaintext-computing computation, when the
encryption's absorbed leaf body reproduces an earlier absorbed body.
Each construction-generated final leaf transcript therefore belongs to
at most one $\badleaf$ pair, and at most $\nleaf(\mathcal{B})$ pairs can
occur.

Each pair's collision is decided at the first creation of its
second-created member. At that moment the earlier member's
$\tauleaf$-bit projection is a function of the prior execution, while
the later member's projection is a fresh uniform lazy sample: by
Transcript Injectivity (\Cref{lem:transcript-inj}) two construction
reaches of the same bridged input come from the same final leaf
transcript, so no construction computation has read that bridged input
before the creation, and a direct query reading it before the creation
would decode under $\KeyedTWOut$ to the hidden key and its $\Ltag(j)$
class (\Cref{lem:bridge-decode})---precisely
$\mathsf{hit}_{\mathsf{leaf}}$. Hence before
$\mathsf{hit}_{\mathsf{leaf}}$, conditioned on the entire prior
execution, each pair collides on its $\tauleaf$-bit projections with
probability $2^{-\tauleaf}$, and a union bound over the at most
$\nleaf(\mathcal{B})$ pairs gives the claim. No assumption on what the
visible interface exposes is needed: a projection revealed after its
sampling cannot alter a collision already decided. By Leaf-Transcript
Recovery (\Cref{lem:leaf-recovery}), on $\neg\badleaf$ equal flattened
final-root inputs of an encryption computation and any construction
computation identify the same leaf transcript sequence.
\end{proof}

\subsection{Encryption Transcript Marginal Uniformity}
\label{app:enc-transcript-marginal-uniformity}

\begin{lemma}[Encryption Transcript Marginal Uniformity]
\label{lem:enc-transcript-marginal-uniformity}
For any fixed nonce-respecting encryption query in the \TW{}
random oracle experiment, let $\mathsf{hit}_{\mathsf{fresh}}$ be the
event that, before this query reaches some counted transcript point,
a prior direct adversary query to $\ROflat$ has already hit the
bridged input of that hidden-key point. On
$\neg\mathsf{hit}_{\mathsf{fresh}}$, every transcript output of the
query that contributes to a returned ciphertext block, to an internal
leaf tag, or to the final root tag is either fresh when this query first
generates its bridged input, or reuses the lazy-table value sampled
uniformly when that input was first generated by an earlier construction
computation. The counted output points first generated by the fixed
query are pairwise distinct.
\end{lemma}

\begin{proof}
Work on an execution outside $\mathsf{hit}_{\mathsf{fresh}}$ for the
fixed query. This rules out prior adversary reads of the hidden-key
construction inputs considered below. A previous construction computation
may nevertheless have generated one of these inputs, for example through
a non-replay verification query using the same nonce. Such a lookup
samples the lazy-table value uniformly when it first occurs; if the fixed
encryption query later reaches the same bridged input, it reuses that
same value. Thus it remains only to show that the counted output points
newly generated by the fixed encryption query are distinct from one
another and from previous nonce-respecting encryption queries. Conditional
information leaked by earlier verification comparisons of a final root
tag is accounted for separately by \Cref{lem:ver-first-root-guess}, not
by this lemma.

\paragraph{First root keystream block.}
The first root keystream block, when present, is produced by the root
duplex after the initialization call $\prfxroot \concat K
\concat U$ and any AD absorption (elided when $A = \varepsilon$, in
which case the \siginit call itself emits the first root keystream
block). Since the nonce-respecting adversary never repeats
a nonce, no previous encryption query can have used the same root
initialization transcript. Within the current query,
\Cref{lem:transcript-inj} separates initialization, AD, message, and
aggregation calls by their suffixes and by the duplex-call encoding,
so the first root keystream transcript is distinct from the other
counted output points generated by this query. If the requested
keystream length is zero, no ciphertext block is returned at that point
and there is nothing to prove for that output.

\paragraph{First leaf keystream block.}
For each leaf chunk $j \geq 1$, the first leaf keystream block is
produced after the initialization call $\prfxleaf \concat
K \concat U \concat \langle j \rangle_8$. The prefix separates leaves
from the root, the nonce separates this encryption query from all
previous encryption queries, and $\langle j \rangle_8$ separates
leaves within this query. Hence every first leaf keystream transcript
newly generated by the query is distinct.

\paragraph{Later message outputs.}
Argue inductively within each root or leaf message stream. Suppose
the current message-output transcript is either first generated by this
query or reuses a value generated by an earlier construction computation.
If its output is used as a keystream block, write its $\ROflat$ value as
$Z_i$. Since the adaptive plaintext block $M_i$ is a deterministic
function of the prior view,
\[
  C_i = M_i \oplus Z_i
\]
is now fixed. \TW{} then absorbs $C_i$ with the appropriate message
phase suffix ($\sigmsgmore$ or
$\sigmsglast$) before requesting the next keystream block. Once $C_i$ is fixed, the next transcript is
fixed; \Cref{lem:transcript-inj}'s injectivity and domain separation
imply it has not appeared earlier unless the same duplex-call prefix
has appeared earlier. Nonce uniqueness excludes this across encryption
queries. By \Cref{cor:output-sep}, the next message-keystream output
point can coincide with a within-query earlier output point only if
both belong to the same output-point class on the same duplex
instance (root or leaf-$j$ with matching $j$) and have the same
duplex-call sequence; the current message-stream prefix is strictly
longer than every earlier prefix in that stream, so this is excluded.
Across distinct instances within the same query (root vs leaf, or
two leaves with distinct $j$), the prefixes already differ in the
first call's $\sigma$ by $\prfxroot /
\prfxleaf$ or $\langle j \rangle_8$ respectively. Each
later message-output transcript that emits keystream is therefore
distinct from the other points generated by this query, unless it is a
reuse of a point generated by an earlier construction computation. The
terminal message-output transcript of a leaf emits
the leaf tag rather than a following keystream block; the same
separation argument applies to that transcript, so each internal leaf
tag newly generated by the query is sampled at a distinct point. Empty
inputs still induce a duplex call, but with zero keystream output they
contribute no ciphertext block.

\paragraph{Final root tag transcript.}
The final root tag is sampled at the root tag output point selected by
$\|C\|$: the $\Rtag$ point after the root duplex absorbs the
aggregation transcript built from the leaf tags and $\langle n
\rangle_8$ when $n \geq 1$, and the $\Rmsg^+$ point closing the root
message phase when $n = 0$ (\Cref{sec:enc}). When $n \geq 1$, the leaf
tags are themselves $\ROflat$ outputs on final leaf transcripts;
condition on their sampled values, on the sampled
ciphertext, and on the adversary-visible prior view and $\ROflat$
table state, so that the root aggregation transcript is fixed. When
$n = 0$ there is no aggregation transcript; condition on the sampled
ciphertext and the adversary-visible prior view alone. By
\Cref{cor:output-sep}, the selected root tag output-point class is
separated from every other output-point class in the same query ---
including root initialization (suffix \siginit), AD-phase outputs,
the remaining message phase outputs, leaf-side output points (separated
by $\prfxleaf$), and, in the $n \geq 1$ case, earlier \Ragg-prefix
output points (separated by the \sigaggmore vs
\sigagglast suffix). Even if two encryption queries realize
the same aggregation transcript (or, for $n = 0$, the same root
ciphertext), the full root transcript also
contains the root initialization $\prfxroot \concat K
\concat U$, and nonce uniqueness separates that initialization from
all previous encryption queries. The final root tag transcript is
therefore either newly generated at a distinct point of the fixed query,
or is a reuse of a value first sampled by an earlier construction
computation. In the newly generated case, its $\tauroot$-bit $\ROflat$
output is a fresh lazy-table sample; in the reuse case, it is the same
uniform sample drawn when the bridged input was first generated.
\end{proof}

\subsection{Mu-Ind Hybrid Games and Verification Accounting}
\label{app:mu-verification-coupling}

The direct mu-ind adjacent-hybrid proof bounds one user-replacement step
of the hybrid through an explicit game sequence, defined next, using the
generic key-test and leaf-collision accounting lemmas together with the
verification-specific facts below. We use the root tag output point and
final root transcript classes of \Cref{def:root-tag-output-point}; in
particular, computations on a common ciphertext share the same root tag
output point because $\leaves(C)$ is fixed by $\|C\|$.

\begin{definition}[Adjacent-Hybrid Games]
\label{def:adjacent-hybrid-games}
Fix a flat-RO mu-ind adversary $\mathcal{B}$, the pre-sampled keys
$K_1, \dots, K_D$ of \Cref{sec:mu-setting}, and a user index
$d \in \{1, \dots, D\}$. On the branch where two sampled keys collide,
every game below outputs the same fixed bit; the remaining clauses
describe the branch with pairwise distinct keys. Every game maintains
one lazy table $\Theta$ implementing $\ROflat$: $\Theta[Y]$ samples a
fresh uniform infinite string at the first access to $Y$ and returns
the stored string afterwards---equivalently, $\Theta$ is sampled in
full at the outset and read on demand. Direct adversary queries
$(Y, \ell)$ are answered by $\lfloor \Theta[Y] \rfloor_\ell$; the
\emph{environment} answers users $e < d$ by $(\Rand, \Replay)$ and
users $e > d$ by $(\Enc_{K_e}^\RO, \Ver_{K_e}^\RO)$, reading
construction outputs from $\Theta$; and the $\mathsf{New}$ interface
and the validity and nonce-respecting checks of \Cref{sec:mu-setting}
are common to every game. Every game maintains the flags $\badkey$,
$\badleaf$, and $\forge$, initialized false; no oracle answer or branch
depends on a flag, and $\badkey$ is set whenever a direct query input
decodes under $\KeyedTWOut$ to $K_d$ and a counted output-point class
of \Cref{tab:keyed-twout}.
\begin{itemize}
\item Games \Greal{} and \Grej{} answer user $d$ by the shared oracle
  code of \Cref{alg:adjacent-hybrid-games} and differ exactly in the
  marked return statement: a non-replay verification whose comparison
  succeeds sets $\forge$ and returns accept in \Greal{} and reject in
  \Grej{}.
\item Game \Gideal{} answers user-$d$ encryption queries with fresh
  uniform byte strings of length $\|M\| + \tauroot/8$, recorded in
  $W$, and user-$d$ verification by the replay check alone; it performs
  no user-$d$ construction computation and sets neither $\badleaf$ nor
  $\forge$. \Gideal{} is the hybrid $H_d^\star$ of
  \Cref{lem:mu-adjacent-hybrid}, instrumented with $\badkey$.
\end{itemize}
The flag $\badleaf$ is set, in \Greal{} and \Grej{}, when a user-$d$
construction computation first creates a final leaf transcript whose
$\tauleaf$-bit projection equals that of a distinct, already created
final leaf transcript sharing its $(K_d, U, j)$ slot, at least one of
the two generated by an encryption computation---the same-slot event of
\Cref{lem:leaf-collision}, part~2, for user $d$. It guards no code
difference and is used only in the accounting.
\end{definition}

\begin{algorithm}[t]
  \caption{Shared user-$d$ oracle code of the games \Greal{} and
    \Grej{} (\Cref{def:adjacent-hybrid-games}). $\Theta$ is the shared
    lazy $\ROflat$ table, read at bridged transcript prefixes, and $W$
    the user-$d$ replay record. Inputs failing the validity or
    nonce-respecting checks of \Cref{sec:mu-setting} are rejected
    beforehand, as in every game. The two games differ exactly in the
    marked lines~\ref{line:ret-real} and~\ref{line:ret-rej}.}
  \label{alg:adjacent-hybrid-games}
  \begin{algorithmic}[1]
    \Procedure{RO}{$Y$, $\ell$} \Comment{direct adversary query}
      \State set $\badkey$ if $\KeyedTWOut(Y) = (K_d, \mathcal{C})$ for
        a counted class $\mathcal{C}$ (\Cref{tab:keyed-twout})
      \State \Return $\lfloor \Theta[Y] \rfloor_{\ell}$
    \EndProcedure
    \Statex
    \Procedure{Enc}{$U$, $A$, $M$}
      \State run $\Enc_{K_d}^{\RO}(U, A, M)$, reading each requested
        duplex output as the corresponding projection of $\Theta$ at
        its bridged transcript prefix (\Cref{sec:ro-model}); let the
        result be $C \concat T$
      \State set $\badleaf$ if a final leaf transcript first created
        here completes a same-slot collision
        (\Cref{def:adjacent-hybrid-games})
      \State record $(U, A, C \concat T)$ in $W$; \Return $C \concat T$
    \EndProcedure
    \Statex
    \Procedure{Vf}{$U$, $A$, $C \concat T$}
      \If{$(U, A, C \concat T) \in W$} \Return accept \EndIf
      \State determine the verification transcript of $(K_d, U, A, C)$;
        read from $\Theta$ the $\tauleaf$-bit leaf tags at its final
        leaf transcripts, setting $\badleaf$ as in \textsc{Enc}, and
        nothing at its stream output points
      \State $T' \gets \lfloor \Theta[\flatmap(R)] \rfloor_{\tauroot}$
        at the root tag output point $R$ of the computation
        (\Cref{def:root-tag-output-point})
      \If{$T' = T$}
        \State $\forge \gets \mathsf{true}$
        \State \Return accept \Comment{in \Greal{} only}
          \label{line:ret-real}
        \State \Return reject \Comment{in \Grej{} only}
          \label{line:ret-rej}
      \EndIf
      \State \Return reject
    \EndProcedure
  \end{algorithmic}
\end{algorithm}

\begin{lemma}[Environment Neutrality]
\label{lem:env-neutrality}
In each game of \Cref{def:adjacent-hybrid-games}, on the branch with
pairwise distinct keys, the environment's construction-internal
$\ROflat$ lookups---made while simulating $\Enc_{K_e}^\RO$ and
$\Ver_{K_e}^\RO$ for users $e > d$---neither set a user-$d$ flag nor
generate the bridged input of a user-$d$ counted transcript point.
\end{lemma}

\begin{proof}
Whenever such a lookup lies in a counted keyed transcript class,
\Cref{lem:bridge-decode} decodes its candidate key as $K_e$, which on
the distinct-keys branch differs from $K_d$; a lookup outside the
counted keyed language is irrelevant to the flags. By the key-field
disjointness of \Cref{lem:key-renaming}, the environment's addresses
are disjoint from user $d$'s. The flags are set only by direct
adversary queries decoding to $K_d$ ($\badkey$), by user-$d$ final leaf
transcript creations ($\badleaf$), and by user-$d$ verification
comparisons ($\forge$), none of which the environment performs. Ideal
users $e < d$ perform no construction lookups at all.
\end{proof}

\begin{lemma}[Verification Visible-Prefix Key Uniformity]
\label{lem:ver-key-uniformity}
In the stopped single-key \TW{} random oracle encryption/verification
interaction, or its \textsf{INT-RUP} extension with a
plaintext-computing oracle, with hidden key $K$, direct access to $\ROflat$, and
immediate abort on $\badkey$, order the direct $\ROflat$ queries. For a
direct query position $i$, let $S_i$
be the set of distinct candidate keys decoded by $\KeyedTWOut$ from
earlier direct-query inputs. Condition on any positive-probability
visible execution prefix immediately before query $i$, on no earlier
$\badkey$, and on $S_i$. Then the hidden key $K$
is uniform on $\bits^k \setminus S_i$. Consequently the uniformity
invariant required by \Cref{lem:key-tests} holds in the stopped
interaction.
\end{lemma}

\begin{proof}
The visible prefix contains the adversary's oracle queries and replies:
direct $\ROflat$ outputs, encryption outputs and the resulting replay
table, plaintext-computing replies when that oracle is present,
verification bits, and the adversary state determined by these
values. Fix such a prefix and two keys $K_0, K_1 \notin S_i$.

Let $\Phi_{0 \to 1} = \Phi_{K_0 \to K_1}$ be the key-renaming map of
\Cref{lem:key-renaming} and $\Phi_{1 \to 0}$ its inverse. By that lemma,
$\flatmap \circ \Phi_{0 \to 1} \circ \flatmap^{-1}$ is a bijection
$\mathcal{Y}_{K_0} \to \mathcal{Y}_{K_1}$ between the bridged hidden-key
construction addresses for $K_0$ and for $K_1$, preserving output-point
classes and equality relations among addresses.

This bijection is applied online, not to a pre-existing fixed address
set. Whenever a construction computation under $K_0$ would first query
a hidden-key address $Y$, the coupled computation under $K_1$ queries
$\Phi_{0 \to 1}(Y)$ and receives the same lazy-table value; repeated
queries are coupled consistently through the same map. If the queried
address is a root aggregation address containing leaf tag bytes, those
bytes are already coupled equal because the corresponding leaf tag
lookups were either repeated or first sampled earlier under the same
renaming rule. Thus adaptively generated message transcripts and root
aggregation transcripts are renamed step by step with the same visible
ciphertext, leaf tag, and root tag bytes.

Prior direct-query inputs are not renamed: they are adversary-chosen
visible strings. Since $K_0, K_1 \notin S_i$, neither
$\mathcal{Y}_{K_0}$ nor $\mathcal{Y}_{K_1}$ contains a prior direct-query
input (\Cref{lem:key-renaming}, direct-query avoidance). Direct-query
lazy-table values are therefore coupled identically and remain disjoint
from the renamed hidden
construction addresses.

Use the same adversary coins, the same direct-query lazy-table values,
and the online renamed construction-internal lazy-table values.
Encryption outputs have the same distribution under the coupling,
because every construction lookup that determines ciphertext, leaf
tags, or the final root tag is paired with a distinct renamed lookup
carrying the same sampled value. Plaintext-computing computations are
renamed in the same online way: a plaintext reply is the submitted
ciphertext block XORed with a coupled keystream projection at a stream
output point (\Cref{tab:keyed-twout}), so it carries the same visible
bytes under both keys, and the recomputed tag is not returned. Replay
verification decisions depend
only on the visible replay table. Non-replay verification computations
are renamed in the same online way as encryption computations, and the
final root tag comparison therefore gives the same accept/reject bit.
Verification does not reveal plaintext on rejection or acceptance; it
reveals only that bit.

Thus every visible prefix with hidden key $K_0 \notin S_i$ has the
same conditional probability as the corresponding execution with hidden
key $K_1 \notin S_i$. Since $K$ is initially uniform and the condition
``no earlier $\badkey$'' removes exactly the keys in
$S_i$, all keys in $\bits^k \setminus S_i$ remain equally likely.
\end{proof}

\paragraph{Foreign-key key-test claim.}
In each game of \Cref{def:adjacent-hybrid-games}, work in the original
experiment, where the key-collision branch outputs a fixed bit rather
than conditioning on pairwise distinct keys. Then
\[
  \Pr[\badkey] \leq \KeyGuess_k(\qRO(\mathcal{B})).
\]
Fix the foreign keys and leave their transcript addresses unchanged.
Before the first user-$d$ key hit, the visible prefix is invariant under
renaming $K_d$ among keys that have not already been tested and are not
equal to a fixed foreign key, by the same online coupling as
\Cref{lem:ver-key-uniformity} and the key-field disjointness of
\Cref{lem:key-renaming}. For each fresh candidate key tested by one of
$\mathcal{B}$'s direct $\ROflat$ queries, the contributing event is
contained in $\{K_d \text{ equals that candidate}\}$ and costs at most
$2^{-k}$ in the original pre-sampled-key experiment; a candidate equal
to a fixed foreign key contributes only on the fixed-output
key-collision branch and therefore contributes nothing to
$\badkey$. A union bound over at most $\qRO(\mathcal{B})$ direct
queries gives the claim.

\begin{lemma}[Verification-First Root Guessing]
\label{lem:ver-first-root-guess}
In the game \Grej{} of \Cref{def:adjacent-hybrid-games}, stopped at the
first occurrence of $\badkey \cup \badleaf$---with $\badleaf$ the
same-slot event of \Cref{lem:leaf-collision}, part~2, for user
$d$---let $\mathsf{win}_{\mathsf{vf}}$ be the event that some
non-replay user-$d$ verification comparison succeeds within a
verification-first final root transcript class. Then
\[
  \Pr[\mathsf{win}_{\mathsf{vf}}]
  \leq \qv^{(d)} / 2^{\tauroot},
\]
where $\qv^{(d)}$ is the user-$d$ non-replay verification cap of
\Cref{sec:mu-setting}.
\end{lemma}

\begin{proof}
Throughout, a \emph{non-replay verification submission} means a valid
non-replay user-$d$ verification query; by the validity convention, an
invalid verification query returns $0$ and performs no construction
computation, so it realizes no bridged final root input, belongs to no
final root transcript class, and contributes to no $G_{\mathcal{C}}$
below.

In \Grej{}, every non-replay submission performs its construction
computation and returns reject (\Cref{alg:adjacent-hybrid-games}), so
no comparison outcome is visible to the adversary, and each class's
root tag value is read from the lazy table only by the class's own
comparisons and by any later encryption computation reaching the class.

Fix a verification-first final root transcript class $\mathcal{C}$ and
follow the stopped execution until the first encryption query
generating the same class, if any (the cutoff). Let $G_{\mathcal{C}}$
be the set of distinct candidate tags submitted by non-replay
verification queries in class $\mathcal{C}$ before that cutoff;
repeated tags only shrink this set. The class's correct root tag
$U_{\mathcal{C}}$ is the $\tauroot$-bit projection of $\ROflat$ at the
class's bridged root tag input. Before the stop, no direct query has
read this hidden-key input: such a query would be decoded by
$\KeyedTWOut$ to the hidden key and would set $\badkey$; the
environment reads no user-$d$ address (\Cref{lem:env-neutrality}).
Before the cutoff, no encryption query has revealed $U_{\mathcal{C}}$
as an encryption tag. (When the class's root tag point is $\Rmsg^+$,
i.e.\ $n = 0$, an encryption query on a longer message that shares
chunk $0$ may touch this same address while closing its chunk-$0$
message phase, but it requests zero output there and continues into its
aggregation phase, so it does not reveal $U_{\mathcal{C}}$; its own tag
is emitted at a later $\Rtag$ point.) The comparisons against
$U_{\mathcal{C}}$ return the constant reject, so the visible execution
and the set $G_{\mathcal{C}}$ are functions only of the adversary's
coins, visible oracle replies, and lazy-table values at inputs other
than this class's root tag address. By the lazy-sampling order
exchange, conditioned on those values, $U_{\mathcal{C}}$ may
equivalently be sampled after $G_{\mathcal{C}}$ is fixed, and is
uniform independently of $G_{\mathcal{C}}$. The probability that some
pre-cutoff comparison in class $\mathcal{C}$ succeeds---that
$U_{\mathcal{C}} \in G_{\mathcal{C}}$---is therefore at most
$|G_{\mathcal{C}}|/2^{\tauroot}$.

After the cutoff, if it exists, class $\mathcal{C}$ admits no
successful non-replay comparison. Let the cutoff encryption query
return $(U_e,A_e,C_e,T_e)$. Any later verification submission in the
same final root transcript class has, by the final-root determinacy of
\Cref{lem:transcript-inj} and Leaf-Transcript Recovery
(\Cref{lem:leaf-recovery}) under the absence of $\badleaf$---which
supplies the recovery hypothesis because the cutoff computation is an
encryption computation---the same $(U,A,C)$ as that encryption
computation. If the submitted tag is $T_e$, the tuple is a recorded
encryption tuple and the query is answered by the replay check before
any comparison; if the submitted tag is not $T_e$, the comparison
fails.

Each non-replay verification submission deterministically reconstructs
one final root transcript at the root tag output point, so it
contributes its submitted tag to at most one set $G_{\mathcal{C}}$, and
only when its class is verification-first and the submission precedes
that class's cutoff. Thus
$\sum_{\mathcal{C}} |G_{\mathcal{C}}| \leq \qv^{(d)}$, and a union
bound over the verification-first classes gives the claim.
\end{proof}

\begin{lemma}[\textsf{INT-RUP} Root Guessing]
\label{lem:int-rup-root-guess}
In the stopped single-key \TW{} random oracle
encryption/plaintext-computing/verification interaction before
$\badkey \cup \badleaf$---with $\badleaf$ the same-slot event of
\Cref{lem:leaf-collision}, part~2---let $\forge$ be the event that some valid
non-replay verification query accepts. Then
\[
  \Pr[\forge \text{ before } \badkey \cup \badleaf]
  \leq \qv(\mathcal{B}) / 2^{\tauroot}.
\]
\end{lemma}

\begin{proof}
Throughout, a \emph{non-replay verification submission} means a valid
non-replay verification query; invalid verification queries return $0$,
perform no construction computation, realize no bridged final root
input, and contribute to no set of guessed tags.

Work in the game stopped at the first occurrence of
$\badkey$ or $\badleaf$. A final root transcript class is
\emph{tag-revealed} once an encryption query has returned the root tag
for that class. If a non-replay verification submission belongs to a
tag-revealed class, compare it with an encryption computation that
revealed the class. The two computations' final-root transcript
prefixes are well-formed typed transcript prefixes with the same
bridged final root input; by
\Cref{lem:bridge-decode}, they have the same native flattened input.
Then final-root determinacy (\Cref{lem:transcript-inj}) and
Leaf-Transcript Recovery (\Cref{lem:leaf-recovery}) under
$\neg\badleaf$---which supplies the recovery hypothesis because the
comparand is an encryption computation---recover the same $(U,A,C)$. If the submitted tag is the
revealed tag, the tuple is a replay-table hit; if it is any other tag,
the final comparison rejects. Thus any accepting non-replay
verification submission must lie in a class whose root tag has not yet
been revealed by encryption.

For tag-unrevealed classes, define an all-reject continuation:
encryption and plaintext-computing queries are answered normally,
replay verification queries accept normally, and each non-replay
verification submission performs the same construction computation but
returns reject to the adversary. Plaintext-computing queries may
generate the same stream and leaf tag transcripts as decryption, but
they do not reveal the final root tag value. Equivalently, the final
root tag sample for a tag-unrevealed class can be deferred until it is
first needed for an encryption answer or a verification comparison;
this does not change any plaintext reply. By Output-Point Separation
(\Cref{cor:output-sep}) and the keystream schedule recorded at
\Cref{tab:keyed-twout}, the keystream values used to compute returned
plaintext blocks are projections at stream output points, not at the
final root tag output point; in particular the chunk-$0$ $\Rmsg^+$ call
of an $n \geq 1$ computation requests zero output, so no plaintext
reply exposes an $\Rmsg^+$ tag value.

Fix a final root transcript class $\mathcal{C}$ whose first
construction generation is not by an encryption query; for a class
first generated by an encryption query the tag is revealed from that
generation onward, so every verification submission in it is covered by
the previous paragraph. Follow the all-reject path until the first
encryption query that reveals the class, if any (the cutoff; any such
query is later than the class's first generation). Let
$G_{\mathcal{C}}$ be the set of distinct
tags submitted by non-replay verification queries in class
$\mathcal{C}$ before that cutoff. Before the cutoff, no encryption query
has returned the class tag and no direct query has read the hidden-key
root tag point without triggering $\badkey$. Conditioned on the
all-reject path and on all lazy-table values except this final root tag
projection, the correct tag for $\mathcal{C}$ is uniform in
$\bits^{\tauroot}$ and independent of $G_{\mathcal{C}}$. The
probability that any pre-cutoff submitted tag in $G_{\mathcal{C}}$
equals the deferred class tag is therefore at most
$|G_{\mathcal{C}}|/2^{\tauroot}$.

Consider the event that the first accepting non-replay verification
submission in the stopped game belongs to class $\mathcal{C}$. Up to
that global first accept, the adversary's visible answers match the
all-reject continuation: every earlier non-replay submission rejects,
replay and plaintext-computing queries are answered identically in both
executions, and encryption answers are unchanged. The construction
computations also agree at every address except that the continuation
defers the final root tag value for class $\mathcal{C}$ until the
comparison is needed. Therefore the first accepting non-replay
submission can belong to class $\mathcal{C}$ only before the cutoff,
and only if its submitted tag equals the deferred class tag, which has
probability at most $|G_{\mathcal{C}}|/2^{\tauroot}$.

The execution-uniform count $\qv(\mathcal{B})$ applies to the
all-reject continuation. Each valid non-replay verification query
contributes its submitted tag
to at most one such set $G_{\mathcal{C}}$, and replay-table hits do not
contribute. Thus
$\sum_{\mathcal{C}} |G_{\mathcal{C}}| \leq \qv(\mathcal{B})$, and a
union bound over the classes not first generated by an encryption
query, each with its pre-cutoff guess set $G_{\mathcal{C}}$, gives the
claim.
\end{proof}

\subsection{Mu-Ind Adjacent-Hybrid Lemma}

The multi-user indistinguishability theorem (\Cref{thm:mu-ind}) telescopes a
hybrid over the $D$ users, replacing one real user by its ideal
$(\Rand,\Replay)$ interface at each step. The following lemma bounds one such
step through the game sequence of \Cref{def:adjacent-hybrid-games}: the
instrumented real game \Greal{}, the all-reject game \Grej{} differing from
it in a single return statement, and the ideal game \Gideal{}, with the
divergence events set as the flags $\badkey$, $\badleaf$, and $\forge$ by
game code.

\begin{lemma}[Mu-Ind Adjacent-Hybrid Bound]
\label{lem:mu-adjacent-hybrid}
Let $\mathcal{B}$ be a flat-RO mu-ind adversary in the pre-sampled-key
formulation of \Cref{sec:mu-setting}, with keys
$K_1,\dots,K_D$ sampled at the start of the experiment. Fix
$d \in \{1,\dots,D\}$. Let $H_{d-1}^\star$ and $H_d^\star$ be the
adjacent stopped hybrids that output a fixed bit if any two sampled keys
collide, and otherwise run $\mathcal{B}$ with one shared lazy
$\ROflat$ table as follows: users $e < d$ are answered by
$(\Rand,\Replay)$, users $e > d$ are answered by
$(\Enc_{K_e}^\RO,\Ver_{K_e}^\RO)$, and user $d$ is answered by
$(\Enc_{K_d}^\RO,\Ver_{K_d}^\RO)$ in $H_{d-1}^\star$ and by
$(\Rand,\Replay)$ in $H_d^\star$. If $\nleaf^{(d)}$ and $\qv^{(d)}$
are the user-$d$ resource caps, then
\[
  \bigl|\Pr[\mathcal{B}^{H_{d-1}^\star} \Rightarrow 1]
       - \Pr[\mathcal{B}^{H_d^\star} \Rightarrow 1]\bigr|
  \leq
  \KeyGuess_k(\qRO(\mathcal{B}))
  + \frac{\nleaf^{(d)}}{2^{\tauleaf}}
  + \frac{\qv^{(d)}}{2^{\tauroot}}.
\]
\end{lemma}

\begin{proof}
The fixed-output branch on a key collision is identical in
$H_{d-1}^\star$, $H_d^\star$, and the games of
\Cref{def:adjacent-hybrid-games}, so only the branch with pairwise
distinct keys contributes; work on that branch, with the keys fixed,
throughout. The proof identifies $H_{d-1}^\star$ with \Greal{}
(Claim~1), relates \Greal{} to \Grej{} pointwise (Claim~2) and \Grej{}
to $H_d^\star = \Gideal{}$ in distribution (Claim~3), and assembles the
bound by a single first-occurrence decomposition in \Grej{} (Claim~4).

\paragraph{Claim 1: \Greal{} realizes $H_{d-1}^\star$.}
The user-$d$ oracles of \Cref{alg:adjacent-hybrid-games} implement
$(\Enc_{K_d}^\RO, \Ver_{K_d}^\RO)$ with three distribution-preserving
refactors. First, reading every construction output as a projection of
the shared table $\Theta$ at its bridged transcript prefix is the lazy
implementation of $\ROflat$ (\Cref{sec:ro-model}); whether $\Theta$ is
sampled in full at the outset or entry by entry on first access is
immaterial. Second, \textsc{Vf} reads $\Theta$ only at leaf tag and
root tag output points. This is sound because decryption absorbs the
submitted ciphertext, so the verification transcript is determined by
$(K_d, U, A, C)$ independently of any stream output
(\Cref{lem:duplex-flat}, \Cref{sec:dec}), and verification discards the
recovered plaintext, so the keystream projections a real
$\Ver_{K_d}^\RO$ computation reads influence neither a visible answer
nor the value of any other table entry; leaving those entries unread
until some later computation first accesses them leaves the joint
distribution of all oracle answers unchanged. Third, the flags are
instrumentation: no oracle answer or branch depends on them. Hence
$\Pr[\mathcal{B}^{\Greal} \Rightarrow 1]
= \Pr[\mathcal{B}^{H_{d-1}^\star} \Rightarrow 1]$.

\paragraph{Claim 2: \Greal{} and \Grej{} agree until $\forge$.}
Run the two games on the same adversary coins, environment coins, and
table realization $\Theta$. The games share all code, including the
comparison that sets $\forge$, and differ in exactly one statement, the
marked return line of \Cref{alg:adjacent-hybrid-games}, which executes
only when $\forge$ has just been set. The two executions are therefore
identical---every oracle answer, table access, record in $W$, and flag
update---up to and including the query at which $\forge$ is first set,
and the first differing value is that query's returned bit. In
particular, on executions in which \Grej{} never sets $\forge$, the two
runs coincide, outputs included, and the flags $\badkey$ and $\badleaf$
are set at the same queries in both games up to the first $\forge$.

\paragraph{Claim 3: \Grej{} agrees with $H_d^\star$ until $\badkey$.}
Couple $H_d^\star = \Gideal{}$ to \Grej{} as follows: run \Grej{} on
the same adversary and environment coins, give the \Gideal{} coordinate
the same visible oracle answers for every query answered while
$\badkey$ is unset, and let the \Gideal{} coordinate continue
independently according to its own law from the query that sets
$\badkey$ onward. We claim the resulting \Gideal{} coordinate is distributed
exactly as \Gideal{}. The answer rules coincide query by query except
for user-$d$ encryption: non-replay user-$d$ verification returns
reject in \Grej{} and under $\Replay$ alike, and replay queries are
answered from the record $W$, which the coupling keeps equal in both
coordinates; the environment runs identical code in both games; and
direct queries read table entries that agree in both games before
$\badkey$---an entry written by a user-$d$ construction computation is
read by a direct query only if that query decodes under $\KeyedTWOut$
to $K_d$ and a counted class (\Cref{lem:key-renaming}), which sets
$\badkey$, while entries written by the environment are produced
identically in both games (\Cref{lem:env-neutrality}).

It remains to show that, before $\badkey$, each user-$d$ encryption
answer of \Grej{} is, conditioned on the entire prior visible
transcript, a fresh uniform byte string of length $\|M\| +
\tauroot/8$---the law of $\Rand$. Fix such a query $(U, A, M)$. Before
$\badkey$, no direct query has read the bridged input of a user-$d$
hidden-key counted point, so the hypothesis
$\neg\mathsf{hit}_{\mathsf{fresh}}$ of Encryption Transcript Marginal
Uniformity (\Cref{lem:enc-transcript-marginal-uniformity}) holds, and
every counted output point first generated by this query is pairwise
distinct and receives a fresh uniform sample. The stream output points
are always first generated by this query: user-$d$ verification
computations read no stream points (\Cref{alg:adjacent-hybrid-games}),
nonce-respecting encryption and Transcript Injectivity
(\Cref{lem:transcript-inj}) separate them from every earlier user-$d$
encryption computation, and the environment writes only foreign-key
addresses (\Cref{lem:env-neutrality}). The ciphertext body $C$ is
therefore a fresh uniform string. The internal leaf tags are either
fresh samples or reuse entries written by user-$d$ verification
computations; they are never returned and influence the answer only
through the root tag address. The root tag $T$ is read at the
computation's root tag output point (\Cref{def:root-tag-output-point}),
an address distinct from every stream point of the same query by
Output-Point Separation (\Cref{cor:output-sep}). If that entry is first
generated here, $T$ is a fresh uniform sample. Otherwise the entry was
written by an earlier user-$d$ verification computation---before
$\badkey$, the only other writer of user-$d$ addresses---whose reads of
it produced only the constant reject answers of \Grej{}; since no
oracle answer or branch depends on the flags, the visible transcript is
a function of the coins and of table entries other than this one, and
conditioned on the visible transcript the entry remains uniform. In
either case $(C, T)$ has the $\Rand$ law, and at most one user-$d$
encryption reaches any root tag address, since two would share the
nonce $U$, contradicting nonce-respecting encryption. By induction over
the query sequence, the coupled \Gideal{} coordinate has exactly the
law of $H_d^\star$, and its visible transcript agrees with \Grej{}'s
while $\badkey$ is unset.

\paragraph{Claim 4: assembly.}
Place the three games on one probability space: shared adversary and
environment coins and one table realization $\Theta$ drive \Greal{} and
\Grej{} pointwise, and the $H_d^\star$ coordinate is coupled to \Grej{}
as in Claim~3. Until the first occurrence of $\forge \vee \badkey$ in
\Grej{}, the three visible transcripts coincide: \Greal{}'s by Claim~2,
since $\forge$ has not yet been set, and $H_d^\star$'s by Claim~3,
since $\badkey$ has not. When neither flag is ever set, the final
outputs of $\mathcal{B}$ coincide as well, so the outputs can differ
only on the event $\forge \vee \badkey$. With Claim~1 identifying
$\Pr[\mathcal{B}^{H_{d-1}^\star} \Rightarrow 1] =
\Pr[\mathcal{B}^{\Greal} \Rightarrow 1]$ and Claim~3 identifying
$\Pr[\mathcal{B}^{H_d^\star} \Rightarrow 1]$ with the coupled
coordinate's output probability,
\[
  \bigl|\Pr[\mathcal{B}^{H_{d-1}^\star} \Rightarrow 1]
       - \Pr[\mathcal{B}^{H_d^\star} \Rightarrow 1]\bigr|
  \;\leq\;
  \Pr_{\Grej}[\forge \vee \badkey].
\]
Decompose by first occurrence, inserting $\badleaf$---which guards no
code difference and enters only this accounting:
\[
\begin{aligned}
  \Pr_{\Grej}[\forge \vee \badkey]
  \;\leq\;
  \Pr[\badkey]
  &+ \Pr[\badleaf \text{ before } \badkey \vee \forge] \\
  &+ \Pr[\forge \text{ before } \badkey \vee \badleaf].
\end{aligned}
\]
The three terms are bounded in \Grej{}. For the first, the foreign-key
key-test claim above gives
$\Pr[\badkey] \leq \KeyGuess_k(\qRO(\mathcal{B}))$. For the second,
\Grej{} restricted to user $d$ is a hidden-key game with
nonce-respecting encryption and all-reject verification; before
$\badkey$ no direct query has read a user-$d$ hidden final leaf
point---such a query would decode to $K_d$ and its $\Ltag(j)$ class and
set $\badkey$---and the environment creates no user-$d$ leaf transcript
(\Cref{lem:env-neutrality}), so the same-slot case of
\Cref{lem:leaf-collision} gives
$\Pr[\badleaf \text{ before } \badkey \vee \forge]
\leq \nleaf^{(d)}/2^{\tauleaf}$. For the third, before
$\badkey \vee \badleaf$ every successful comparison lies in a
verification-first class: a class first generated by a user-$d$
encryption computation admits none, since by the final-root determinacy
of \Cref{lem:transcript-inj} and Leaf-Transcript Recovery
(\Cref{lem:leaf-recovery})---whose hypothesis $\neg\badleaf$ supplies
because the comparand is the class's encryption computation---any
non-replay submission in the class carries the same $(U,A,C)$ as that
encryption and either matches the recorded tuple, in which case the
replay check answers it before any comparison, or fails the comparison.
\Cref{lem:ver-first-root-guess} then gives
$\Pr[\forge \text{ before } \badkey \vee \badleaf]
\leq \qv^{(d)}/2^{\tauroot}$.

Summing the three bounds gives
\[
  \bigl|\Pr[\mathcal{B}^{H_{d-1}^\star} \Rightarrow 1]
       - \Pr[\mathcal{B}^{H_d^\star} \Rightarrow 1]\bigr|
  \leq
  \KeyGuess_k(\qRO(\mathcal{B}))
  + \frac{\nleaf^{(d)}}{2^{\tauleaf}}
  + \frac{\qv^{(d)}}{2^{\tauroot}},
\]
with the user-$d$ key-test event $\badkey$ charged a single time for
the adjacent hybrid.
\end{proof}

\section{The Flat Sponge Lift}
\label{app:flat-sponge-lift}

This appendix supplies the flat sponge lift machinery
(\Cref{lem:lift-bound,lem:derived-adversary,cor:lift-application}) invoked
by the AEAD security bounds of \Cref{sec:aead-security}, the root tag bounds of
\Cref{sec:root-tag-hash-security}, and the committing-notion bound of
\Cref{cor:committing-combined}. The machinery applies uniformly to the games
named once here.

The lift has two steps in the win-event games and three in the
real-vs-ideal games. First, the \cite{BDPV08} indifferentiability theorem
replaces the real permutation used by the \TW{} construction and by direct
primitive queries with the simulator over a flat random oracle, at cost
$\Lift_c(\Ntot)$, where $\Ntot = N + \Nprim$ bounds the operational primitive
trace: construction-internal duplex calls induced by construction queries and
deterministic challenger computations, together with direct primitive queries.
The simulator-side construction trace is then identified with the flat-RO
construction interface by the operational trace lemma below.
Second, the simulator-world adversary is converted into a flat-RO
adversary whose direct random oracle query count is at most $\Nprim$. Thus the
random oracle bounds are instantiated at $\qRO \leq \Nprim$, while the public
permutation lift pays $\Lift_c(N+\Nprim)$. Third, in the real-vs-ideal games,
whose ideal world exposes the real permutation rather than the simulator
(\Cref{sec:games-pp}), the ideal-side simulator introduced by the first step
is exchanged back for $(\pi,\pi^{-1})$ at cost $\Lift_c(\Nprim)$
(\Cref{lem:sim-perm-proximity}); the exchange is possible because the ideal
construction oracles are independent of the primitive.

\begin{definition}[Single-Stage \TW{} Games]
\label{def:single-stage-games}
The \emph{single-stage \TW{} games} are the games of
\Cref{sec:games-pp,sec:mu-setting} in which one adversary interacts with the
challenger in a single stage, so that a \cite{BDPV08} replacement applies to
an entire world's public-permutation transcript at once, including
deterministic challenger checks after the adversary's final output. They are
all-in-one AE, mu-ind, \textsf{INT-RUP}, CMTDs-4, the $\HTW$ collision and
preimage games (standard and everywhere), and CMTs-$\ell$ for
$\ell \in \{1,3,4\}$. The win-event games among these have one world; the
real-vs-ideal games---all-in-one AE and mu-ind---have two, and the lift below
treats their real world by the full replacement and their ideal world by the
simulator--permutation exchange of \Cref{lem:sim-perm-proximity}.
\end{definition}

The derived adversary forwards the following game data to the corresponding
flat-RO experiment.
\begin{center}
\small
\begin{tabular}{@{}l p{0.62\linewidth}@{}}
  \toprule
  Game family & Forwarded data \\
  \midrule
  AE / mu-ind / \textsf{INT-RUP}
    & construction queries, including $\mathsf{New}$ where present and replay state \\
  $\HTW$ hash games & final structured hash game output \\
  CMTDs / CMTs & final openings and deterministic challenger checks \\
  \bottomrule
\end{tabular}
\end{center}

\subsection{Simulator Convention}
\label{app:simulator}

Throughout this appendix, $(\Sim^{\ROflat}, \Sim^{-1,\ROflat})$
denotes the public primitive simulator supplied by \cite[Theorem~2]{BDPV08}
--- the indifferentiability theorem for the simple
padded sponge over a random permutation --- after applying the
\cite[Theorem~3]{BDPV11} bridge of \Cref{lem:bridge-decode} to \TW{}'s native
$\padtenone$ transcript inputs. The superscript records that both
simulator directions share the same internal lazy $\ROflat$ table.

The simulator is a proof device: each game of
\Cref{sec:games-pp,sec:mu-setting} states its advantage over the real
permutation, and the reduction to the corresponding random oracle model
adversary passes through intermediate simulator worlds. In those worlds,
the construction interface is the one specified by the game
($(\Rand,\Replay)$, instantiated per user in mu-ind, on the ideal side)
and the public primitive interface is
$(\Sim^{\ROflat}, \Sim^{-1,\ROflat})$. The $\ROflat$ table in this
notation is internal to that system and is not an extra interface
exposed to the application adversary.

\begin{lemma}[At Most One Simulator RO Call]
\label{lem:sim-ro-call-filter}
For the simulator $(\Sim^{\ROflat},\Sim^{-1,\ROflat})$, every
primitive-interface query induces at most one simulator-side $\ROflat$ call.
More precisely, every inverse primitive-interface query and every forward
primitive-interface query violating one of the following conditions induces no
simulator $\ROflat$ call. The symbols $s$, $p$, $\mathcal{R}$, $\mathcal{O}$,
and $\mathcal{C}$ are the local internal notation of the BDPV simple-padded
sponge simulator, not \TW{} resource notation:
\begin{description}
  \item[(C1) New edge.] The simulator's internal graph has no
    outgoing edge from the input node $s$.
  \item[(C2) Unsaturated.] $\mathcal{R} \cup \mathcal{O} \neq
    \mathcal{C}$.
  \item[(C3) Rooted input.] The input node is rooted.
  \item[(C4) Paddable path.] The constructed path $p$ decomposes
    uniquely as $p = p' \concat 0^{rj}$ with $p'$ not ending in
    $0^r$, and the prefix $p'$ is a valid padded sponge input:
    equivalently, $p' = x \concat \padten[r](|x|)$ for an unpadded
    BDPV08 random oracle input $x$.
\end{description}
When all four conditions hold, the forward primitive-interface query
induces one simulator $\ROflat$ call. Consequently, each
primitive-interface query induces at most one simulator-side
$\ROflat$ call.
\end{lemma}

\begin{proof}
This is the branch structure of the \cite{BDPV08} simple-padded
permutation simulator. The simulator invokes the random oracle only on
the forward-query branch where the input
node is rooted, unsaturated, paddable to an unpadded random oracle
input $x$, and not already extended by an outgoing edge in the
simulator graph. Forward queries violating any
of these conditions and all inverse queries return outputs sampled or
read from the simulator's internal graph without invoking the
random oracle. Under the \cite[Theorem~3]{BDPV11} bridge fixed in
\Cref{app:simulator}, that simulator random oracle invocation is the
simulator-side $\ROflat$ call used throughout this appendix.
\end{proof}

\begin{lemma}[Operational Duplex Trace]
\label{lem:operational-duplex-trace}
In any execution of a single-stage \TW{} game
(\Cref{def:single-stage-games}), the construction computations counted
by $N$ admit an operational primitive trace of cost at most $N$:
each root or leaf duplex call contributes one forward primitive query,
and the requested output bits are the corresponding projection of the
state returned by that query. When this trace is run against the
\cite{BDPV08} simulator $(\Sim^{\ROflat},\Sim^{-1,\ROflat})$ in an
execution whose combined primitive trace has cost less than $2^c$, the
construction outputs are the flat-RO values
$\RF(X)=\ROflat(\flatmap(X))$ at the same well-formed \TW{} transcript
prefixes $X$.
\end{lemma}

\begin{proof}
A \TW{} duplex call absorbs a single suffixed input block, padded by
$\padtenone$, applies the permutation once, and returns at most $r$
rate bits from the resulting state. Thus the public-permutation
implementation of any root or leaf construction computation is already
a primitive trace with one forward query per duplex call. The definition
of $N$ in \Cref{sec:resources} counts exactly these root initialization,
associated data, message, leaf initialization, leaf message, and
aggregation duplex calls, including the deterministic challenger checks
of the hash and committing games.

\Cref{lem:duplex-flat} is the \cite[Lemma~3]{BDPV11} duplex-to-sponge
identity specialized to \TW{}: the output of each duplex call is the
output of the sponge after absorbing the accumulated flattened
duplex-call transcript, and no additional permutation call is made after
the current absorption block when the requested output length is at most
$r$. The bridge of \Cref{lem:bridge-decode}, fixed in
\Cref{app:simulator}, maps these $\padtenone$ transcript prefixes to the
\cite{BDPV08} padded-sponge input domain without changing the output bits
at $r=r_{\max}$.

When the same operational trace is run against
$(\Sim^{\ROflat},\Sim^{-1,\ROflat})$, every construction-internal forward
query lies on a rooted, paddable path determined by a well-formed \TW{}
prefix. In an execution whose combined primitive trace has cost less
than $2^c$, the simulator cannot have saturated: each primitive query
adds at most one capacity value to the simulator's internal
$\mathcal{R} \cup \mathcal{O}$ set, so fewer than $2^c$ queries cannot
fill the $2^c$-element capacity space. The simulator's
sponge-consistency invariant \cite[Lemma~2]{BDPV08} therefore makes the
rate part returned at that node agree with the corresponding lazy value
$\ROflat(\flatmap(X))$. Thus the simulator-side operational construction
has the same construction-output law as the flat-RO construction
interface used in \Cref{sec:ro-model}.
\end{proof}

\begin{lemma}[Flat Sponge Lift Bound]
\label{lem:lift-bound}
Let $\mathsf{G}$ be a single-stage \TW{} game
(\Cref{def:single-stage-games}). Let
$\mathcal{A}$ be a
public permutation $\mathsf{G}$ adversary with construction interface
cost at most $N$ and at most $\Nprim$ direct primitive-interface
queries. In the \cite{BDPV08} single-stage replacement, implement the
\TW{} construction computations by the operational duplex trace of
\Cref{lem:operational-duplex-trace} and keep the adversary's direct
primitive-interface queries as primitive queries. This primitive
transcript has \cite{BDPV08} query cost at most
\[
  \Ntot = N + \Nprim .
\]
Consequently replacing the real public permutation and \TW{} construction
by the simulator-world primitive interface
$(\Sim^{\ROflat},\Sim^{-1,\ROflat})$ and the corresponding flat-RO
construction interface costs at most $\Lift_c(\Ntot)$.
\end{lemma}

\begin{proof}
The \cite{BDPV08} replacement is applied before the derived-adversary
construction. View the entire single-stage game, including its
construction interface and deterministic challenger checks, as a
primitive-only distinguisher that runs the \TW{} algorithms internally.
By \Cref{lem:operational-duplex-trace}, the construction side contributes
at most $N$ forward primitive queries: one per duplex call. This includes
plaintext-computing computations, verification computations whether they
accept or reject, the per-user sum of construction calls in mu-ind, the
deterministic challenger decryption or encryption checks charged in the
committing games, and the hash games' challenger encryptions---the
collision game's $s$ submitted-input encryptions, the standard preimage
game's source encryption and post-output check, and the
everywhere-preimage game's post-output check.

Each direct adversary query to $(\pi,\pi^{-1})$ contributes one
primitive-interface query and is counted in $\Nprim$. These two counts
are disjoint in the public permutation experiment: $N$ counts
construction-internal use of the primitive, while $\Nprim$ counts only
the adversary's public primitive interface queries.

Construction work is modeled by the operational duplex trace, rather
than by a collection of standalone \cite{BDPV08} construction queries at
every flattened duplex prefix. The standalone-query model would assign a
nested-prefix transcript the sum of the individual sponge-evaluation
costs, while the operational trace is the same stateful computation the
construction actually performs, with the output of each call read from
the state just returned by that call. Hence the execution-uniform caps
bound the primitive-only query cost by $N+\Nprim=\Ntot$. If
$\Ntot \geq 2^c$ then $\Lift_c(\Ntot) \geq 1$ and the claimed bound holds
trivially, so assume $\Ntot < 2^c$, satisfying the hypothesis of
\cite[Theorem~2]{BDPV08}. Applying that theorem to this primitive-only
distinguisher gives the real-to-simulator replacement, and
\Cref{lem:operational-duplex-trace} identifies the simulator-side
construction outputs with the flat-RO construction interface in this
nontrivial case. The theorem gives distinguishing or
success-probability cost at most the random-permutation
differentiating bound $f_P(\Ntot)$ of \cite[Lemma~5]{BDPV08}. The
\cite[Lemmas~4 and~5]{BDPV08} product bounds give the
random-transformation bound $f_T(\Ntot) = 1 - \prod_{i=1}^{\Ntot}(1 -
i/2^c)$, and the factor of $f_P(\Ntot)$ corresponding to the $f_T$ factor
at index $i$ is that factor divided by a quantity
$1 - (i-1)/2^{r+c} \in (0, 1]$, hence at least as large, so
$f_P(\Ntot) \leq f_T(\Ntot)$. With $\Ntot < 2^c$ every factor lies in
$[0, 1]$ and
\[
  f_P(\Ntot) \leq f_T(\Ntot)
  = 1 - \prod_{i=1}^{\Ntot}\Bigl(1 - \frac{i}{2^c}\Bigr)
  \leq \sum_{i=1}^{\Ntot} \frac{i}{2^c}
  = \frac{\Ntot(\Ntot + 1)}{2^{c+1}}
  = \Lift_c(\Ntot),
\]
the second inequality by the Weierstrass product inequality
$\prod_i (1 - x_i) \geq 1 - \sum_i x_i$ for $x_i \in [0, 1]$.
\end{proof}

\begin{lemma}[Simulator--Permutation Proximity]
\label{lem:sim-perm-proximity}
Let $\mathcal{D}$ be a distinguisher with access only to a two-directional
primitive interface, making at most $q$ queries. Then
\[
  \bigl|
    \Pr[\mathcal{D}^{\Sim^{\ROflat},\, \Sim^{-1,\ROflat}} \Rightarrow 1]
    - \Pr[\mathcal{D}^{\pi,\, \pi^{-1}} \Rightarrow 1]
  \bigr|
  \;\leq\; \Lift_c(q),
\]
where $\pi \sample \Perm(\bits^{1600})$ and $\ROflat$ is the lazily
sampled random oracle internal to the simulator.
\end{lemma}

\begin{proof}
Consider the indifferentiability distinguisher $\mathcal{D}'$ that runs
$\mathcal{D}$, forwards its primitive-interface queries, and never queries
the construction interface. The \cite{BDPV08} query cost of $\mathcal{D}'$
is at most $q$: each forwarded query costs one, and no construction queries
arise. Because $\mathcal{D}'$ ignores the construction interface, its two
views are exactly $\mathcal{D}^{\pi,\pi^{-1}}$ in the real pair and
$\mathcal{D}^{\Sim^{\ROflat},\Sim^{-1,\ROflat}}$ in the ideal pair of
\cite[Theorem~2]{BDPV08}, applied under the \cite[Theorem~3]{BDPV11}
bridge fixed in \Cref{app:simulator}. If $q \geq 2^c$ then
$\Lift_c(q) \geq 1$ and the claim is trivial; otherwise that theorem
bounds the displayed difference by the random-permutation differentiating
bound $f_P(q)$ of \cite[Lemma~5]{BDPV08}, and the derivation in the proof
of \Cref{lem:lift-bound} gives $f_P(q) \leq \Lift_c(q)$.
\end{proof}

\begin{lemma}[Derived-Adversary Execution Contract]
\label{lem:derived-adversary}
Let $\mathsf{G}$ be a single-stage \TW{} game
(\Cref{def:single-stage-games}). Let
$\mathcal{A}$ be a
simulator-world $\mathsf{G}$ adversary whose public primitive
interface is answered by
$(\Sim^{\ROflat}, \Sim^{-1,\ROflat})$, with at most $\Nprim$ direct
primitive-interface queries. Define
$\mathcal{B}_{\mathrm{derived}}(\mathcal{A})$ to run $\mathcal{A}$
internally, answer each primitive-interface query by running the
same simulator against direct finite-output queries to its own
$\ROflat$ oracle, forward construction interface queries unchanged in games
with construction oracles, and forward the adversary's final structured output
unchanged in the output-checked $\HTW$ hash and committing games. Then the
simulator-world execution of $\mathcal{A}$ and the random oracle execution of
$\mathcal{B}_{\mathrm{derived}}(\mathcal{A})$ have identical joint
distributions, and
\[
  \qRO(\mathcal{B}_{\mathrm{derived}}) \leq \Nprim.
\]
In this coupling, simulator-induced $\ROflat$ calls and
construction-internal $\RF$ lookups use the same lazy table. Therefore,
if a simulator-induced query is the bridged input $\flatmap(X)$ of a
well-formed \TW{} transcript prefix $X$ and a forwarded construction
computation later reaches the same prefix $X$, both interfaces read the
same $\ROflat(\flatmap(X))$ value, projected to the requested output
length.
Moreover, $\mathcal{B}_{\mathrm{derived}}$ inherits the construction-side
caps of $\mathcal{A}$:
\begin{align*}
  \qv(\mathcal{B}_{\mathrm{derived}}) &= \qv(\mathcal{A})
    \quad\text{(AE, mu-ind, \textsf{INT-RUP})}, \\
  \nleaf(\mathcal{B}_{\mathrm{derived}}) &\leq \nleaf(\mathcal{A})
    \quad\text{(all games above)}.
\end{align*}
In the mu-ind game these resource inheritances hold per user, and the
$\mathsf{New}$-query transcript is unchanged.
\end{lemma}

\begin{proof}
Couple the two executions by using the same lazy $\ROflat$ table for
the simulator and for all construction-internal $\RF(\cdot)$ lookups.
In the simulator-world execution, $\mathcal{A}$ sees construction
answers from the selected game interface and public primitive answers
from $(\Sim^{\ROflat}, \Sim^{-1,\ROflat})$. In the derived
random oracle execution, $\mathcal{B}_{\mathrm{derived}}$ runs the
same $\mathcal{A}$, forwards construction interface queries or the final
structured output required by the selected game unchanged, and answers
primitive-interface queries by invoking the same simulator code with the same
lazy $\ROflat$ table. The construction interface, hash challenger, or
committing challenger therefore receives the same forwarded queries or final
output, the simulator receives the same primitive queries in the same order, and
the lazy table supplies the same values.
The full views of $\mathcal{A}$ in the simulator-world execution and of
$\mathcal{B}_{\mathrm{derived}}$ in the random oracle execution are identically
distributed.

The simulator's random oracle calls may be on arbitrary
\cite{BDPV08} $H$-input strings, not only bridged \TW{} transcript
inputs. These calls are the direct $\ROflat$ queries of
$\mathcal{B}_{\mathrm{derived}}$ and are counted in
$\qRO(\mathcal{B}_{\mathrm{derived}})$. Construction-internal
$\RF(\cdot)$ lookups made by $\Enc_K^\RO$, $\Dec_K^\RO$,
$\PDec_K^\RO$, $\Ver_K^\RO$, or a committing challenger are forwarded to the
same $\ROflat$ table but are not counted in $\qRO$, exactly as in
\Cref{sec:resources}.

By \Cref{lem:sim-ro-call-filter}, each primitive-interface query
induces at most one direct $\ROflat$ query by
$\mathcal{B}_{\mathrm{derived}}$, and inverse queries or
non-RO-touching forward queries induce none. Since $\mathcal{A}$ makes
at most $\Nprim$ direct primitive-interface queries, this gives
$\qRO(\mathcal{B}_{\mathrm{derived}}) \leq \Nprim$.

\Cref{lem:bridge-decode}'s bridge places \TW{} construction
restrictions and simulator-induced $\ROflat$ calls in a single lazy
table. If the simulator invokes $\ROflat$ on $\flatmap(X)$ for a
well-formed \TW{} transcript prefix $X$, and a forwarded construction
computation later reaches $X$, the construction lookup
$\RF(X)=\ROflat(\flatmap(X))$ reads the same table entry and returns the
requested projection. Since \TW{} uses $r = r_{\max}$, the
\cite[Theorem~3]{BDPV11} output post-processing does not alter \TW{}
output bits. Thus simulator-induced direct queries and
construction-internal lookups at the same bridged input are
transcript-consistent. In the AE-style games, a simulator call whose input is
accepted by $\KeyedTWOut$ for the hidden key is still a direct
$\ROflat$ query and is covered by the same
$\badkey$ accounting used by the direct random oracle mu-ind proof. The
derivation step itself does not
prove a key-guessing statement; it only supplies the derived
adversary and the query bound at which the random oracle bounds of
\Cref{sec:aead-security,sec:root-tag-hash-security} are instantiated.

Finally, the execution-uniform caps of \Cref{sec:resources} hold
along every oracle-answer trace of $\mathcal{A}$, including traces
generated by the simulator on the ideal public primitive interface.
The derived-adversary construction does not add construction interface
encryption, plaintext-computing, or verification queries and forwards
$\mathcal{A}$'s construction interface queries unchanged along each trace.
$\mathcal{B}_{\mathrm{derived}}$ therefore inherits the relevant
construction resources: $\qv(\mathcal{B}_{\mathrm{derived}}) =
\qv(\mathcal{A})$ in the AE, mu-ind, and \textsf{INT-RUP} games, and
$\nleaf(\mathcal{B}_{\mathrm{derived}}) \leq \nleaf(\mathcal{A})$ for
the distinct construction-generated final-leaf-transcript cap. In
mu-ind, construction interface forwarding preserves the
$\mathsf{New}$ calls and these caps user by user.
\end{proof}

\subsection{Lift Application}
\label{app:lift-corollaries}

The lift bound of \Cref{lem:lift-bound}, the simulator--permutation
proximity bound of \Cref{lem:sim-perm-proximity}, and the
derived-adversary construction of \Cref{lem:derived-adversary} combine
into a single template, stated once and then instantiated by the bounds of
\Cref{sec:aead-security,sec:root-tag-hash-security} and by the
committing-notion bounds of \Cref{cor:committing-combined}. Every
public-permutation bound is obtained by adding $\Lift_c(\Ntot)$---plus
$\Lift_c(\Nprim)$ in the real-vs-ideal games---and
substituting $\qRO \leq \Nprim$ in the corresponding random oracle theorem.

\begin{corollary}[Lift Application]
\label{cor:lift-application}
Let $\mathsf{G}$ be a single-stage \TW{} game
(\Cref{def:single-stage-games}). For every
public permutation $\mathsf{G}$-adversary $\mathcal{A}$ against \TW{}
with the resources of \Cref{sec:resources}, the derived adversary
$\mathcal{B} = \mathcal{B}_{\mathrm{derived}}(\mathcal{A})$ of
\Cref{lem:derived-adversary} is a random oracle $\mathsf{G}$-adversary
with $\qRO(\mathcal{B}) \leq \Nprim$ and the construction-side caps of
\Cref{lem:derived-adversary}---in particular $\qv(\mathcal{B}) \leq \Qv$
and $\nleaf(\mathcal{B}) \leq \Nleaf$, where these apply---such that,
when $\mathsf{G}$ is a win-event game,
\[
  \Adv^{\mathsf{G}}_\TW(\mathcal{A})
  \;\leq\;
  \Lift_c(\Ntot) + \Adv^{\mathsf{G}}_{\RO}(\mathcal{B}),
\]
and, when $\mathsf{G}$ is a real-vs-ideal game (all-in-one AE or mu-ind),
\[
  \Adv^{\mathsf{G}}_\TW(\mathcal{A})
  \;\leq\;
  \Lift_c(\Ntot) + \Adv^{\mathsf{G}}_{\RO}(\mathcal{B}) + \Lift_c(\Nprim).
\]
If moreover $\Adv^{\mathsf{G}}_{\RO}(\mathcal{B}) \leq
\varepsilon_{\mathsf{G}}(\qRO(\mathcal{B}), \nleaf(\mathcal{B}))$ for some
$\varepsilon_{\mathsf{G}}$ nondecreasing in each argument, then
$\varepsilon_{\mathsf{G}}(\Nprim, \Nleaf)$ may replace
$\Adv^{\mathsf{G}}_{\RO}(\mathcal{B})$ in the corresponding bound.
\end{corollary}

\begin{proof}
By the flat sponge lift bound (\Cref{lem:lift-bound}), replacing the real
public permutation and \TW{} construction by
$(\Sim^{\ROflat},\Sim^{-1,\ROflat})$ and the flat-RO construction
interface costs at most $\Lift_c(\Ntot)$, where $\Ntot = N + \Nprim$
counts both the construction calls of $N$---including any construction
oracle exposed during $\mathcal{A}$'s run and, in the win-event games,
the challenger's construction computations, pre- and post-output (in the
preimage game, the source encryption computing $T^*$ as well as the
post-output check)---and the
$\Nprim$ direct primitive queries. The derived-adversary contract
(\Cref{lem:derived-adversary}) expresses the resulting simulator-world
execution as the identically distributed random oracle execution of
$\mathcal{B} = \mathcal{B}_{\mathrm{derived}}(\mathcal{A})$, with
$\qRO(\mathcal{B}) \leq \Nprim$ and the stated construction-side
inheritances. For a win-event game this gives the first bound.

For a real-vs-ideal game, chain four worlds:
$\mathsf{W}_1$, the real construction interface with $(\pi,\pi^{-1})$;
$\mathsf{W}_2$, the flat-RO real construction interface with the
simulator; $\mathsf{W}_3$, the ideal construction interface with the
simulator; and $\mathsf{W}_4$, the ideal construction interface with
$(\pi,\pi^{-1})$. The game of \Cref{sec:games-pp,sec:mu-setting}
compares $\mathsf{W}_1$ with $\mathsf{W}_4$. The step from
$\mathsf{W}_1$ to $\mathsf{W}_2$ costs at most $\Lift_c(\Ntot)$ as
above, and under the coupling of \Cref{lem:derived-adversary} the pair
$(\mathsf{W}_2, \mathsf{W}_3)$ is precisely the real/ideal pair of the
random oracle game for $\mathcal{B}$, so
$|\Pr[\mathcal{A}^{\mathsf{W}_2} \Rightarrow 1] -
\Pr[\mathcal{A}^{\mathsf{W}_3} \Rightarrow 1]| \leq
\Adv^{\mathsf{G}}_{\RO}(\mathcal{B})$. For the step from $\mathsf{W}_3$
to $\mathsf{W}_4$, the ideal construction oracles---$(\Rand,\Replay)$,
instantiated per user in mu-ind, together with the $\mathsf{New}$
interface and the validity checks---are independent of the primitive, so
a distinguisher can answer them from its own coins, run $\mathcal{A}$,
and forward the at most $\Nprim$ primitive-interface queries;
\Cref{lem:sim-perm-proximity} bounds this step by $\Lift_c(\Nprim)$.
The triangle inequality combines the three estimates into the second
bound.

The final claim follows by the
monotonicity of $\varepsilon_{\mathsf{G}}$. When $\mathsf{G}$ is
parameterized by a challenge fixed before $\mathcal{A}$ runs---the
everywhere-preimage target tag, or the standard preimage source $X^*$,
whose tag $T^* = \HTW(X^*)$ each coupled world computes---the coupling
holds for each fixed challenge and the challenge-independent bound holds
for the maximum or average over challenges defining
$\Adv^{\mathsf{G}}_\TW(\mathcal{A})$.
\end{proof}

\section{Vectorized Kernel Details}
\label{app:kernels}

This appendix details the architecture-specific cores summarized in
\Cref{sec:impl}: the instruction selection, register layout, and
remainder handling of the \texttt{amd64} and \texttt{arm64} kernels.

\paragraph{\texttt{amd64}.}
The \texttt{amd64} implementation provides AVX-512 and AVX2 cores and selects
between them once at startup from the \texttt{AVX512VL} CPUID feature bit. The
single-duplex permutation has two variants: a general-purpose-register core
adapted from the standard library's \Keccak{} permutation (using the
lane-complement trick so that $\chi$ requires no separate negation), and an
AVX-512 variant; the root duplex uses whichever matches the selected core. The
eight-way AVX-512 core packs the lane-major state into 25 ZMM registers, one per
lane with all eight instances in the register's eight 64-bit slots, and performs
the permutation with no cross-lane shuffles: rotations use \texttt{VPROLQ}, and
\texttt{VPTERNLOGQ} fuses the $\theta$ $D$-term formation and the $\chi$
\mbox{and-not} into single three-input logic instructions. The eight-way AVX2
core, lacking 512-bit registers, runs the eight instances as two groups of four
(four 64-bit lanes per YMM register), rotating by shift-or and computing $\chi$
with \texttt{VPANDN}. The fused chunk kernels combine absorption with the
permutation. The AVX-512 steady-state kernel reads the eight chunks with
contiguous loads, transposes each rate block into the lane-major layout in
registers (\texttt{VPUNPCKLQDQ}/\texttt{VPUNPCKHQDQ} and \texttt{VSHUFI64X2}),
absorbs on the lane-major state, and transposes each output block back for a
contiguous store, leaving only the $7$-byte partial rate lane to a masked gather;
these sequential accesses are prefetcher-friendly, whereas per-lane gathers and
scatters at the chunk stride $\beta$
(\texttt{VPGATHERQQ}/\texttt{VPSCATTERQQ}) would stall on cold memory and cap
the long-message rate. The register-resident kernel that drains a
two-to-seven-chunk remainder instead uses that per-lane gather/scatter form
under a mask, reading the active chunks in place: the transpose's shuffle work
is constant regardless of the chunk count and would be spent partly on inactive
lanes, whereas the gathers scale with the active element count. The AVX2 path
drains the same remainders with a variant of its grouped-by-four kernels that
aims each inactive instance's loads and stores at a dummy block inside the
kernel frame with a zero advance stride, so inactive instances touch no memory
beyond the remainder and the hot loop stays branch-free.

\paragraph{\texttt{arm64}.}
The \texttt{arm64} path uses NEON plus the Armv8.2 SHA-3 extension
(\texttt{FEAT\_SHA3}), available on Apple Silicon and on Neoverse V1/V2 and N2.
The permutation is expressed with the extension's fused instructions:
\texttt{EOR3} (three-input XOR) for the $\theta$ column parities, \texttt{RAX1}
(rotate-by-one and XOR) for the $\theta$ $D$-terms, \texttt{XAR} (XOR then
rotate) to fuse the $\theta$ accumulation with the $\rho$ rotation, and
\texttt{BCAX} (bit-clear and XOR) for $\chi$. The single-duplex core places one
instance per 64-bit \texttt{D} lane across the 25 vector registers; the
eight-way core packs two instances per 128-bit register and processes the batch
of eight as four such pairs, using \texttt{ZIP1} and duplicating moves to
interleave and extract the two instances of each lane in the fused chunk kernel.
The same pair machinery backs a $2$-wide kernel that drains a leaf-chunk
remainder two chunks at a time---one instance per $64$-bit half of a $128$-bit
register---at roughly twice the single-duplex per-chunk rate. Both batched
kernels store the permutation state back into the batch after the closing
block---the $\times 8$ kernel for its first pair, the $2$-wide kernel for both
instances---which is the writeback that permits the root duplex to occupy
lane~$0$.

\section{Test Vectors}
\label{app:test-vectors}

The following test vectors cover the corner cases of TW128 encryption, as
specified in \cref{sec:tw128}.

All byte strings are written in hexadecimal. For each vector we list the key
$K$, nonce $N$, associated data $A$, message $M$, ciphertext $C$, and tag $T$;
the sealed output of encryption is $C \concat T$. The empty string is written
$\varepsilon$. The leaf index is encoded as a little-endian $64$-bit integer and
the phase suffixes are appended least significant bit first.

To keep the longer vectors compact, any value exceeding $32$ bytes is given by
its length in bytes together with its $\mathrm{SHA}\text{-}256$ digest rather
than in full. Every message and associated data byte is defined by
$m_i = (17 i + 31) \bmod 256$ for $i = 0, 1, \dots$, so each such input is
determined by its length alone and can be reconstructed and checked against the
stated digest; a ciphertext given by its length and digest can likewise be
checked by encrypting the corresponding message and hashing the result.

% >>> generated test vectors (regenerate with UPDATE_TW128=1) >>>
% Generated by test_tw128.py; do not edit by hand. The full byte
% strings for every vector are in tw128_tv.json at the repository root.

\noindent
{\footnotesize
\begin{tabular}{@{}l p{0.88\textwidth}@{}}
\toprule
\multicolumn{2}{@{}p{0.94\textwidth}@{}}{\textbf{Vector 1.}~ Empty associated data and empty message: the associated data and aggregation phases are both elided, and the root tag is emitted by the single closing message call with no leaves.} \\
\midrule
  $K$ & \texttt{000102030405060708090a0b0c0d0e0f101112131415161718191a1b1c1d1e1f} \\
  $N$ & \texttt{202122232425262728292a2b2c2d2e2f303132333435363738393a3b3c3d3e3f} \\
  $A$ & $\varepsilon$ \\
  $M$ & $\varepsilon$ \\
  $C$ & $\varepsilon$ \\
  $T$ & \texttt{da65bbda0b20b1e936f5326458166633d1e5a3f7aebc0b318d24b051a4efa697} \\
\end{tabular}}

\medskip

\noindent
{\footnotesize
\begin{tabular}{@{}l p{0.88\textwidth}@{}}
\toprule
\multicolumn{2}{@{}p{0.94\textwidth}@{}}{\textbf{Vector 2.}~ Multi-block associated data with an empty message: the associated data spans two rate blocks and no message bytes are processed.} \\
\midrule
  $K$ & \texttt{000102030405060708090a0b0c0d0e0f101112131415161718191a1b1c1d1e1f} \\
  $N$ & \texttt{202122232425262728292a2b2c2d2e2f303132333435363738393a3b3c3d3e3f} \\
  $A$ & 168 bytes, SHA-256\newline \texttt{4c7d125a3488a7603e2f403f0010a6060b8321b60b0fd2471120ae627bbe759a} \\
  $M$ & $\varepsilon$ \\
  $C$ & $\varepsilon$ \\
  $T$ & \texttt{2093878418f3fca96ad3a8d1748a9be4d2a8b518e917d988fd703dcb7bf72173} \\
\end{tabular}}

\medskip

\noindent
{\footnotesize
\begin{tabular}{@{}l p{0.88\textwidth}@{}}
\toprule
\multicolumn{2}{@{}p{0.94\textwidth}@{}}{\textbf{Vector 3.}~ Empty associated data and a single-block message: the root message phase only, no leaves.} \\
\midrule
  $K$ & \texttt{000102030405060708090a0b0c0d0e0f101112131415161718191a1b1c1d1e1f} \\
  $N$ & \texttt{202122232425262728292a2b2c2d2e2f303132333435363738393a3b3c3d3e3f} \\
  $A$ & $\varepsilon$ \\
  $M$ & \texttt{1f30415263748596a7b8c9daebfc0d1e2f405162738495a6b7c8d9eafb0c1d2e} \\
  $C$ & \texttt{676a1441e20daea0f1ee4571ee4543842c946d59e0d3df6f1dd2f92dc4f876aa} \\
  $T$ & \texttt{1873332d359ff2fe5197488224af7bf6ad821875fdbd964100de424aa77342be} \\
\end{tabular}}

\medskip

\noindent
{\footnotesize
\begin{tabular}{@{}l p{0.88\textwidth}@{}}
\toprule
\multicolumn{2}{@{}p{0.94\textwidth}@{}}{\textbf{Vector 4.}~ Empty associated data and a two-block message: the root keystream is chained across two rate blocks.} \\
\midrule
  $K$ & \texttt{000102030405060708090a0b0c0d0e0f101112131415161718191a1b1c1d1e1f} \\
  $N$ & \texttt{202122232425262728292a2b2c2d2e2f303132333435363738393a3b3c3d3e3f} \\
  $A$ & $\varepsilon$ \\
  $M$ & 168 bytes, SHA-256\newline \texttt{4c7d125a3488a7603e2f403f0010a6060b8321b60b0fd2471120ae627bbe759a} \\
  $C$ & 168 bytes, SHA-256\newline \texttt{df3f0a98a91a69ed474edbc52e0dc1d8e54bfd6fc514966ce7b878d48be0bf9a} \\
  $T$ & \texttt{aab2413074a202cfadd5eb784663b1831eb330025972d163d03e9ed5a964b765} \\
\end{tabular}}

\medskip

\noindent
{\footnotesize
\begin{tabular}{@{}l p{0.88\textwidth}@{}}
\toprule
\multicolumn{2}{@{}p{0.94\textwidth}@{}}{\textbf{Vector 5.}~ A message of exactly one chunk: the root is full, there is still no leaf, and the aggregation phase is elided so the closing root message call emits the tag.} \\
\midrule
  $K$ & \texttt{000102030405060708090a0b0c0d0e0f101112131415161718191a1b1c1d1e1f} \\
  $N$ & \texttt{202122232425262728292a2b2c2d2e2f303132333435363738393a3b3c3d3e3f} \\
  $A$ & $\varepsilon$ \\
  $M$ & 8183 bytes, SHA-256\newline \texttt{6532a370a4ab6b5f4094dd8a6016512ab8e8416abe910bc1a0b532979224151d} \\
  $C$ & 8183 bytes, SHA-256\newline \texttt{8d58f8aa3413ea0ba1b8cc3ed4e9a334521d8dafe9694a0661fb8251049907e9} \\
  $T$ & \texttt{c2201637f1116c1303b8982622756770d23acb85384f4658fcb63c37e4b629dd} \\
\end{tabular}}

\medskip

\noindent
{\footnotesize
\begin{tabular}{@{}l p{0.88\textwidth}@{}}
\toprule
\multicolumn{2}{@{}p{0.94\textwidth}@{}}{\textbf{Vector 6.}~ A message of one chunk and one further byte: one root chunk and one leaf, exercising leaf tag aggregation.} \\
\midrule
  $K$ & \texttt{000102030405060708090a0b0c0d0e0f101112131415161718191a1b1c1d1e1f} \\
  $N$ & \texttt{202122232425262728292a2b2c2d2e2f303132333435363738393a3b3c3d3e3f} \\
  $A$ & $\varepsilon$ \\
  $M$ & 8184 bytes, SHA-256\newline \texttt{c4c383bca0808e2ce44d481ddfe56efdcd7a147dcd5a14c898ef5423f2604cd7} \\
  $C$ & 8184 bytes, SHA-256\newline \texttt{ef1b2f7607ade503ab497fc61e62cef2b44c6328fa06e499718eb7da5fdd04d3} \\
  $T$ & \texttt{c96bea48e07a8431fd19721ffbe8afe6a73047f96d953e3281f10126eeca06a1} \\
\end{tabular}}

\medskip

\noindent
{\footnotesize
\begin{tabular}{@{}l p{0.88\textwidth}@{}}
\toprule
\multicolumn{2}{@{}p{0.94\textwidth}@{}}{\textbf{Vector 7.}~ A multi-leaf message of two chunks and one further byte: the full tree path with two leaves.} \\
\midrule
  $K$ & \texttt{000102030405060708090a0b0c0d0e0f101112131415161718191a1b1c1d1e1f} \\
  $N$ & \texttt{202122232425262728292a2b2c2d2e2f303132333435363738393a3b3c3d3e3f} \\
  $A$ & $\varepsilon$ \\
  $M$ & 16367 bytes, SHA-256\newline \texttt{17e31f3700d0602aab8a686024ecc05b8192e80a531483c67dd90d5a26557c5e} \\
  $C$ & 16367 bytes, SHA-256\newline \texttt{6d83d8fbdc2a7c2ef590d29a36527a45019e50d1f8c8db037959081685a46908} \\
  $T$ & \texttt{cdc2e50ac5380419f06772a2f3af2146050bd0340b35f61b67f839351767804e} \\
\end{tabular}}

\medskip

\noindent
{\footnotesize
\begin{tabular}{@{}l p{0.88\textwidth}@{}}
\toprule
\multicolumn{2}{@{}p{0.94\textwidth}@{}}{\textbf{Vector 8.}~ Maximum associated data with a short message: the associated data blocks dominate the cost.} \\
\midrule
  $K$ & \texttt{000102030405060708090a0b0c0d0e0f101112131415161718191a1b1c1d1e1f} \\
  $N$ & \texttt{202122232425262728292a2b2c2d2e2f303132333435363738393a3b3c3d3e3f} \\
  $A$ & 337 bytes, SHA-256\newline \texttt{04ca7b24930bc6b5fe8b95ff1eb0933907657c1283cc3a7e69c5f7f02220f281} \\
  $M$ & \texttt{1f30415263748596a7b8c9daebfc0d1e} \\
  $C$ & \texttt{0dd615a0a05ed82a78c3999b9e80cbed} \\
  $T$ & \texttt{22e4e6725f548ad0660041d66cf2f18dc45c69b447cdaae1f147d33a3b6aa1a9} \\
\end{tabular}}

\medskip

\noindent
{\footnotesize
\begin{tabular}{@{}l p{0.88\textwidth}@{}}
\toprule
\multicolumn{2}{@{}p{0.94\textwidth}@{}}{\textbf{Vector 9.}~ An all-ones key and nonce with a multi-chunk message.} \\
\midrule
  $K$ & \texttt{ffffffffffffffffffffffffffffffffffffffffffffffffffffffffffffffff} \\
  $N$ & \texttt{ffffffffffffffffffffffffffffffffffffffffffffffffffffffffffffffff} \\
  $A$ & 168 bytes, SHA-256\newline \texttt{4c7d125a3488a7603e2f403f0010a6060b8321b60b0fd2471120ae627bbe759a} \\
  $M$ & 8184 bytes, SHA-256\newline \texttt{c4c383bca0808e2ce44d481ddfe56efdcd7a147dcd5a14c898ef5423f2604cd7} \\
  $C$ & 8184 bytes, SHA-256\newline \texttt{49818c320ab282296639798d840cbd60751e1623f449dcbf9423371d44145dfb} \\
  $T$ & \texttt{310c9fff20e9651f8ed47eed9bb182a086ed1af1b3f1799e35307a4139865fc7} \\
\bottomrule
\end{tabular}}

\medskip
% <<< generated test vectors <<<


\end{document}
